\documentclass[11pt, a4paper, oneside]{report}
% English edition of the A4 course "Architecture des ordinateurs".
% Same palette and boxes as the Arabic edition; compile from en/tex/.

\usepackage[a4paper, top=2.2cm, bottom=2.2cm, left=2cm, right=2cm,
    headheight=14pt]{geometry}
\usepackage{amsmath, amssymb}
\usepackage{booktabs, colortbl, tabularx, multirow, array}
\usepackage{tikz}
\usepackage[skins, breakable]{tcolorbox}
\usepackage{fontawesome5}
\usepackage{listings}
\usepackage{enumitem}
\usepackage{titlesec}
\usepackage{fancyhdr}

\usetikzlibrary{positioning, calc, arrows.meta, shapes.geometric, fit,
    backgrounds, decorations.pathreplacing, matrix}

%------------------------------------------------
% Fonts (English, left to right)
%------------------------------------------------
\usepackage{fontspec}
\setmainfont{texgyrepagella}[Extension=.otf, UprightFont=*-regular, BoldFont=*-bold, ItalicFont=*-italic, BoldItalicFont=*-bolditalic]
\setsansfont{texgyreheros}[Extension=.otf, UprightFont=*-regular, BoldFont=*-bold, ItalicFont=*-italic, BoldItalicFont=*-bolditalic]
\newfontfamily\khatt{texgyreheros}[Extension=.otf, UprightFont=*-regular, BoldFont=*-bold]
\newcommand{\textenglish}[1]{#1}
% Latin monospace for code
\newfontfamily\codefont[Scale=0.92, BoldFont=lmmonolt10-bold.otf,
    ItalicFont=lmmono10-italic.otf]{lmmono10-regular.otf}

\linespread{1.12}
\setlength{\parindent}{0pt}
\setlength{\parskip}{5pt}

%------------------------------------------------
% Colors
%------------------------------------------------
\definecolor{brandTeal}{HTML}{5CC8E6}
\definecolor{brandOrange}{HTML}{FFA27F}
\definecolor{brandLight}{HTML}{F8F9FA}
\definecolor{brandDark}{HTML}{263238}
\definecolor{brandMuted}{HTML}{6C757D}
\definecolor{brandGreen}{HTML}{7BC47F}
\definecolor{brandPurple}{HTML}{B39DDB}
\colorlet{tealDeep}{brandTeal!60!black}
\colorlet{orangeDeep}{brandOrange!70!black}
\definecolor{codeGrey}{HTML}{F1F3F5}
\color{brandDark}

%------------------------------------------------
% Chapter and section titles
%------------------------------------------------
\titleformat{\chapter}[display]
    {\khatt\bfseries}
    {\begin{tikzpicture}
        \node[fill=brandTeal, text=white, rounded corners=6pt, inner xsep=12pt,
            inner ysep=6pt, font=\khatt\bfseries\LARGE]
            {\chaptertitlename\ \thechapter};
     \end{tikzpicture}}
    {10pt}
    {\Huge\color{brandDark}}
    [{\color{brandTeal}\titlerule[2pt]}]
\titlespacing*{\chapter}{0pt}{-10pt}{24pt}
\titleformat{\section}
    {\khatt\bfseries\Large\color{tealDeep}}{\thesection}{0.6em}{}
\titleformat{\subsection}
    {\khatt\bfseries\large\color{brandDark}}{\thesubsection}{0.6em}{}
\titleformat{\subsubsection}
    {\bfseries\color{orangeDeep}}{}{0em}{\faAngleRight\ }

%------------------------------------------------
% Header and footer
%------------------------------------------------
\pagestyle{fancy}
\fancyhf{}
\fancyhead[L]{\small\color{brandMuted}\leftmark}
\fancyhead[R]{\small\color{brandMuted}Computer Architecture}
\fancyfoot[C]{\small\color{brandMuted}\thepage}
\renewcommand{\headrulewidth}{0.4pt}
\renewcommand{\headrule}{{\color{brandTeal!60}\hrule width\headwidth height\headrulewidth}}
\renewcommand{\chaptermark}[1]{\markboth{\chaptertitlename\ \thechapter: #1}{}}
\fancypagestyle{plain}{\fancyhf{}\fancyfoot[C]{\small\color{brandMuted}\thepage}
    \renewcommand{\headrulewidth}{0pt}}

%------------------------------------------------
% Lists
%------------------------------------------------
\setlist[itemize]{label={\color{brandTeal}\small\faChevronRight}, itemsep=2pt, topsep=3pt}
\setlist[enumerate]{label={\color{tealDeep}\bfseries\arabic*.}, itemsep=2pt, topsep=3pt}

%------------------------------------------------
% Boxes
%------------------------------------------------
\tcbset{
    base/.style={enhanced, breakable, arc=4pt, boxrule=1.2pt,
        coltitle=white, coltext=brandDark, fonttitle=\khatt\bfseries,
        toptitle=1mm, bottomtitle=1mm, top=2mm, bottom=2mm,
        left=3mm, right=3mm, before skip=8pt, after skip=8pt}}
% Definition
\newtcolorbox{tarif}[1][]{base, colback=brandTeal!8, colframe=brandTeal,
    title={\faLightbulb\quad Definition#1}}
% Example
\newtcolorbox{mithal}[1][]{base, colback=white, colframe=brandTeal!70!black,
    title={\faPencil*\quad Example#1}}
% Warning
\newtcolorbox{intabih}[1][]{base, colback=brandOrange!10, colframe=brandOrange,
    title={\faExclamationTriangle\quad Watch out#1}}
% Analogy from everyday life
\newtcolorbox{tashbih}[1][]{base, colback=brandPurple!10, colframe=brandPurple,
    title={\faSmile\quad Everyday analogy#1}}
% Step by step
% (unbreakable: a split step box could leave an empty page)
\newtcolorbox{khatwa}[1][]{base, unbreakable, colback=brandGreen!8, colframe=brandGreen,
    title={\faShoePrints\quad Step by step#1}}
% Summary
\newtcolorbox{khulasa}[1][]{base, colback=brandTeal!8, colframe=brandTeal,
    title={\faCheckCircle\quad Summary#1}}
% Exercise and solution
\newcounter{tamrin}[chapter]
\newtcolorbox{tamrin}[1][]{base, colback=white, colframe=brandMuted,
    title={\stepcounter{tamrin}\faEdit\quad Exercise \thechapter.\arabic{tamrin}#1}}
\newtcolorbox{hall}[1][]{base, colback=brandGreen!6, colframe=brandGreen!80!black,
    title={\faCheck\quad Solution#1}}
% Did you know?
\newtcolorbox{hal}[1][]{enhanced, breakable, arc=4pt, boxrule=0pt,
    colback=brandOrange!8, left=4mm, right=4mm, top=2mm, bottom=2mm,
    before skip=8pt, after skip=8pt,
    overlay={\node[anchor=north east, fill=brandOrange, text=white,
        rounded corners=3pt, font=\small\khatt\bfseries, inner sep=3pt]
        at ([xshift=-3mm, yshift=2mm]frame.north east) {\faQuestionCircle\ Did you know?};},
    top=4mm}

%------------------------------------------------
% Terms, inline code, numbers
%------------------------------------------------
% \term{English}{French}: the French name helps with the official syllabus
\newcommand{\term}[2]{\textbf{#1} {\small\color{brandMuted}(Fr.\ \emph{#2})}}
\newcommand{\code}[1]{\mbox{\codefont\color{tealDeep}#1}}
\newcommand{\hex}[1]{\mbox{#1\textsubscript{H}}}
\newcommand{\bin}[1]{\mbox{#1\textsubscript{2}}}
% register transfer: dest <- src
\newcommand{\nql}[2]{\mbox{#1\,$\leftarrow$\,#2}}

%------------------------------------------------
% Tables
%------------------------------------------------
\newcolumntype{R}{>{\raggedright\arraybackslash}X}
\newcolumntype{C}{>{\centering\arraybackslash}X}
\newcommand{\ras}[1]{\cellcolor{brandTeal!25}\textbf{#1}}
\newcommand{\cell}[2]{\leavevmode\textcolor{#1}{#2}}
\renewcommand{\arraystretch}{1.3}
\setlength{\arrayrulewidth}{0.5pt}
\arrayrulecolor{brandMuted!60}

%------------------------------------------------
% Code listings (always inside the english environment)
%------------------------------------------------
\lstdefinelanguage{asm86}{
    morekeywords={mov, add, adc, sub, sbb, inc, dec, neg, mul, imul, div, idiv,
        cmp, and, or, xor, not, test, shl, shr, sal, sar, rol, ror, rcl, rcr,
        jmp, je, jne, jz, jnz, jg, jge, jl, jle, ja, jae, jb, jbe, jc, jnc, jo,
        js, jns, jcxz, loop, loope, loopne, call, ret, iret, int, push, pop,
        pushf, popf, lea, xchg, in, out, nop, hlt, cbw, cwd, movsb, movsw,
        lodsb, stosb, rep, cld, std, cli, sti, clc, stc, xlat, syscall,
        movzx, movsx, cdq, cqo, addps, mulps, movaps, movups, vaddps, vmovups,
        addss, movss},
    morekeywords=[2]{org, db, dw, dd, dq, equ, proc, endp, offset, ptr, byte,
        word, dword, qword, section, global, extern, dup, resb, resq, rel},
    sensitive=false, morecomment=[l]{;}, morestring=[b]', morestring=[b]"}
\lstdefinestyle{base}{
    basicstyle=\codefont\small, columns=fullflexible, keepspaces=true,
    showstringspaces=false, upquote=true,
    keywordstyle=\color{tealDeep}\bfseries,
    keywordstyle=[2]\color{brandPurple!60!black}\bfseries,
    commentstyle=\color{red!75!black},
    stringstyle=\color{orangeDeep},
    numbers=left, numberstyle=\scriptsize\color{brandMuted}, numbersep=8pt,
    xleftmargin=18pt, aboveskip=0pt, belowskip=0pt}
\lstdefinestyle{asm}{style=base, language=asm86}
\lstdefinestyle{c}{style=base, language=C, morekeywords={uint8_t, int64_t}}
\lstdefinestyle{out}{basicstyle=\codefont\small\color{brandDark}, columns=fullflexible,
    keepspaces=true, aboveskip=0pt, belowskip=0pt}

\newtcolorbox{codebox}[1][]{enhanced, colback=codeGrey, colframe=codeGrey,
    boxrule=0pt, arc=2pt, left=4pt, right=4pt, top=3pt, bottom=3pt,
    before skip=6pt, after skip=6pt,
    overlay={\fill[brandTeal!70] (frame.south west) rectangle ([xshift=3pt]frame.north west);}, #1}
\newtcolorbox{outbox}[1][]{enhanced, colback=brandDark, colframe=brandDark, boxrule=0pt,
    arc=2pt, left=6pt, right=6pt, top=3pt, bottom=3pt, before skip=4pt, after skip=8pt,
    colupper=white, #1}
% \asm{ch3/hello.asm}  \cfile{ch4/sum.c}  \screen{ch3/hello.out}
\newcommand{\asm}[2][]{\begin{codebox}\lstinputlisting[style=asm,#1]{../../code/#2}\end{codebox}}
\newcommand{\cfile}[2][]{\begin{codebox}\lstinputlisting[style=c,#1]{../../code/#2}\end{codebox}}
\newcommand{\screen}[2][]{\begin{outbox}{\scriptsize\faDesktop\ \textenglish{\codefont screen}}\par\lstinputlisting[style=out, basicstyle=\codefont\small\color{white},#1]{../../code/#2}\end{outbox}}

% Line-by-line explanation:  \begin{sharh} \s{1}{...} \s{2-3}{...} \end{sharh}
% (tabularx must be called as \tabularx ... \endtabularx inside an environment)
\newenvironment{sharh}
    {\par\small\tabularx{\linewidth}{|>{\centering\arraybackslash}p{1.3cm}|R|}
     \hline \ras{Line} & \ras{What does it do?} \\ \hline}
    {\endtabularx\par}
\newcommand{\s}[2]{#1 & #2 \\ \hline}

%------------------------------------------------
% TikZ styles
%------------------------------------------------
\tikzset{
    lawn/.style={text=#1},
    sahm/.style={-{Stealth[length=2.6mm]}, thick, color=brandMuted},
    sahmT/.style={-{Stealth[length=2.6mm]}, very thick, color=tealDeep},
    sahmO/.style={-{Stealth[length=2.6mm]}, very thick, color=orangeDeep},
    unit/.style={draw=#1, fill=#1!15, very thick, rounded corners=4pt,
        align=center, font=\bfseries\small, inner sep=5pt, text=brandDark},
    unit/.default=brandTeal,
    reg/.style={draw=brandDark, fill=white, thick, minimum height=0.65cm,
        font=\codefont\small, inner sep=2pt},
}


\begin{document}

%================================================
% Cover
%================================================
\begin{titlepage}
\thispagestyle{empty}
\begin{tikzpicture}[remember picture, overlay]
    \fill[brandTeal!12] (current page.north west) rectangle ([yshift=-8cm]current page.north east);
    \draw[brandTeal, line width=1.5pt, opacity=0.35]
        ([xshift=3cm, yshift=4cm]current page.south west) circle (6cm);
    \draw[brandOrange, line width=1pt, opacity=0.35]
        ([xshift=3cm, yshift=4cm]current page.south west) circle (7cm);
    \begin{scope}[shift={([xshift=-2.6cm, yshift=-2.5cm]current page.north east)}, scale=0.62]
        \fill[brandDark] (-1.5,-1.5) rectangle (1.5,1.5);
        \fill[brandTeal] (-1,-1) rectangle (1,1);
        \node[font=\codefont\bfseries\Large, text=white] at (0,0) {CPU};
        \foreach \i in {-1.2,-0.6,0,0.6,1.2} {
            \draw[brandMuted, line width=2pt] (\i,1.5) -- (\i,2);
            \draw[brandMuted, line width=2pt] (\i,-1.5) -- (\i,-2);
            \draw[brandMuted, line width=2pt] (1.5,\i) -- (2,\i);
            \draw[brandMuted, line width=2pt] (-1.5,\i) -- (-2,\i);
        }
    \end{scope}
\end{tikzpicture}

\vspace*{1.2cm}
\begin{flushleft}
    {\sffamily\bfseries\color{tealDeep}\Large Course UEF31}\\[0.6cm]
    {\sffamily\bfseries\color{brandDark}\fontsize{38}{46}\selectfont Computer Architecture}\\[0.4cm]
    {\large\color{brandMuted}Architecture des ordinateurs}
\end{flushleft}

\vspace{4.5cm}
\begin{tcolorbox}[enhanced, colback=white, colframe=brandTeal, arc=6pt, boxrule=1.2pt,
    width=0.85\linewidth]
    \large
    \begin{itemize}
        \item Chapter 1: General organization of the central unit
        \item Chapter 2: Inside the processor
        \item Chapter 3: Case study, the Intel 8086 processor
        \item Chapter 4: Case study, 32-bit and 64-bit processors
        \item Chapter 5: Modern processor architectures
    \end{itemize}
\end{tcolorbox}

\vfill
\begin{flushleft}
    \begin{tabular}{@{}ll}
        \textbf{Teacher:} & Dr. Benabderrezak Youcef \\
        \textbf{Coefficient:} & 3 \quad\textbf{Credits:} 5 \\
        \textbf{Assessment:} & exam 60\%, continuous assessment 40\% \\
    \end{tabular}
\end{flushleft}
\end{titlepage}

%================================================
% How to use this course
%================================================
\chapter*{How to use this course}
\addcontentsline{toc}{chapter}{How to use this course}
\markboth{How to use this course}{}

This course is written for you, even if this is the first time you look inside a computer.
We start with simple ideas. Step by step, we go up to real programs in assembly language.
At the end, we understand how the processors in our phones and laptops work.

\section*{The coloured boxes}
On every page you will find coloured boxes. Each colour has a job:

\begin{tarif}
    A new important word. We also give its French name, because you will meet it in the
    official syllabus and in many books.
\end{tarif}
\begin{tashbih}
    A difficult idea, explained with something you know from daily life: a kitchen, a school, a bus\dots
\end{tashbih}
\begin{khatwa}
    A calculation or an execution, done one small step at a time. We never skip a step.
\end{khatwa}
\begin{intabih}
    A common mistake, or a small detail that changes the result.
\end{intabih}
\begin{hal}
    A fun extra fact. You do not need it for the exam, but it helps you understand.
\end{hal}

\section*{The programs in this course}
\begin{itemize}
    \item Every program is complete, and comes with a table \textbf{``What does each line do?''}.
    \item The 8086 programs are written for the \textbf{EMU8086} emulator. The 64-bit programs are
          written for \textbf{NASM} on \textbf{Linux}.
    \item All the programs \textbf{were really run}. What you see in the black screen \faDesktop{} is the real output,
          and the register tables come from the real execution.
\end{itemize}

\begin{intabih}
    Do not only read! Copy each program into EMU8086, run it one instruction at a time
    (the \emph{Single step} button), and watch the registers change. This is how you learn this course.
\end{intabih}

\tableofcontents

\chapter{General Organization of the Central Unit}

\begin{khulasa}[: what will we learn in this chapter?]
\begin{itemize}
    \item What a computer is, and how it changed through history.
    \item The main parts: the processor, the memory, the input/output devices.
    \item The Von Neumann and Harvard models, and the difference between them.
    \item The types of memory, and how to calculate a memory size.
    \item The three buses: the data bus, the address bus, the control bus.
\end{itemize}
\end{khulasa}

%================================================
\section{What is a computer?}
%================================================

\begin{tarif}
    A \term{computer}{ordinateur} is an electronic machine. It receives \textbf{data},
    processes it with a \textbf{program} (a list of instructions), and gives \textbf{results}.
    It can also store data and results.
\end{tarif}

Everything a computer does fits in four verbs: it \textbf{receives}, it \textbf{processes},
it \textbf{stores}, it \textbf{outputs}.

\begin{center}
\begin{tikzpicture}[node distance=1.4cm]
    \node[unit=brandGreen, minimum width=2.8cm, minimum height=1.2cm] (in) {Input\\{\small keyboard, mouse}};
    \node[unit, minimum width=2.8cm, minimum height=1.2cm, right=of in] (cpu) {Processing\\{\small the processor}};
    \node[unit=brandOrange, minimum width=2.8cm, minimum height=1.2cm, right=of cpu] (out) {Output\\{\small screen, printer}};
    \node[unit=brandPurple, minimum width=2.8cm, minimum height=1.1cm, below=1cm of cpu] (mem) {Storage\\{\small memory, disk}};
    \draw[sahmT] (in) -- (cpu);
    \draw[sahmT] (cpu) -- (out);
    \draw[sahmT, <->] (cpu) -- (mem);
\end{tikzpicture}
\end{center}

Let us take a simple example: you use the calculator app on your computer to compute $25 \times 4$.
\begin{enumerate}
    \item \textbf{Input}: you type 25, then $\times$, then 4 on the keyboard. Each key press becomes a
          number code that goes into the computer.
    \item \textbf{Temporary storage}: the computer keeps the two numbers in memory, because it needs them soon.
    \item \textbf{Processing}: the processor reads the two numbers, runs the multiply instruction, and gets 100.
    \item \textbf{Output}: the computer sends the result to the screen, and you see it.
\end{enumerate}
Every program you use, even a complex one (a game, a browser, a photo app), is in fact millions of
these small steps, done very fast.

\begin{tashbih}
    A computer is like a \textbf{restaurant kitchen}. The ingredients (vegetables, meat) are the \textbf{data}.
    The recipe is the \textbf{program}. The cook is the \textbf{processor}. The fridge is the \textbf{memory}.
    The finished dish is the \textbf{result}. The cook does not decide alone: he follows the recipe step by step.
\end{tashbih}

\subsection{Hardware and software}
\begin{itemize}
    \item \term{Hardware}{matériel}: everything you can touch: the processor, the memory, the screen, the wires.
    \item \term{Software}{logiciel}: the programs that tell the hardware what to do: the operating
          system, the browser, games.
\end{itemize}
In this course we study \textbf{the hardware from the inside}, and we see how it runs programs.

\subsection{The layers of a computer}
We can see a computer as layers, one on top of the other. Each layer uses the services of the layer
below it, without knowing all its details:

\begin{center}
\begin{tikzpicture}[l/.style={draw=#1, fill=#1!15, very thick, minimum width=9cm, minimum height=0.75cm, font=\small}]
    \node[l=brandGreen] at (0,1.8) {Applications: browser, games, word processor};
    \node[l=brandPurple] at (0,0.9) {Operating system: Windows, Linux, Android};
    \node[l=brandOrange] at (0,0) {Machine language and assembly: what the processor understands};
    \node[l=brandTeal] at (0,-0.9) {Hardware: processor, memory, buses, electronic circuits};
    \draw[sahmT] (5,-1.1) -- (5,2) node[midway, right, font=\small] {from machine to human};
\end{tikzpicture}
\end{center}
In this course we work on the \textbf{two lowest layers}: the hardware, and the assembly language
that talks to it directly.

\subsection{How does a program run?}
When you double-click on a program icon, these steps happen:
\begin{enumerate}
    \item The program is stored on the \textbf{disk} (SSD or HDD) as a file of zeros and ones.
    \item The operating system \textbf{copies it into the main memory} (RAM), because the disk is very slow
          and the processor does not read from it directly.
    \item The processor reads the instructions from memory \textbf{one after the other} and runs them.
    \item When you close the program, its space in memory is freed for other programs.
\end{enumerate}

%================================================
\section{A short history of computers}
%================================================

The computer did not appear in one day. These are the main steps:

\begin{center}
\begin{tikzpicture}[
    yr/.style={circle, fill=brandTeal, minimum size=0.35cm, inner sep=0pt},
    txt/.style={text width=2.7cm, align=center, font=\small}]
    \draw[brandTeal, line width=3pt] (0,0) -- (14,0);
    \foreach \x/\y/\t/\p in {
        0/1642/{Pascal's machine\\ adds, subtracts}/above,
        2/1837/{Babbage's machine\\ idea of a program}/below,
        4/1945/{ENIAC\\ vacuum tubes}/above,
        6/1947/{invention of the\\ \textbf{transistor}}/below,
        8/1958/{\textbf{integrated circuit}\\ (IC)}/above,
        10/1971/{first microprocessor\\ Intel 4004}/below,
        12/1978/{the\\ Intel 8086}/above,
        14/2005/{\textbf{multi-core}\\ processors}/below} {
        \node[yr] (n) at (\x, 0) {};
        \node[txt, \p=0.25cm of n] {\textbf{\y}\\ \t};
    }
\end{tikzpicture}
\end{center}

\begin{table}[h]
\centering
\begin{tabular}{|c|c|c|c|}
    \hline
    \ras{Generation} & \ras{Period} & \ras{Main technology} & \ras{Example} \\ \hline
    1st & 1945--1955 & vacuum tubes & ENIAC (30 tons!) \\ \hline
    2nd & 1955--1965 & transistor & IBM 7090 \\ \hline
    3rd & 1965--1980 & integrated circuits & IBM 360 \\ \hline
    4th & 1980--today & microprocessor (VLSI) & personal computer, phone \\ \hline
\end{tabular}
\end{table}

Each new generation brought a technology that made the parts \textbf{smaller}, \textbf{faster}
and \textbf{less hungry for energy}:
\begin{itemize}
    \item \textbf{Vacuum tubes} were big, got very hot, and burned out all the time. Some had to be changed almost every day.
    \item A \textbf{transistor} does the same job (an electric switch that opens and closes), but it is hundreds of times smaller and does not burn out.
    \item An \textbf{integrated circuit} puts thousands of transistors on one small piece of silicon, instead of connecting them with wires.
    \item A \textbf{microprocessor} puts the \textbf{whole} processor on one chip. This gave us the personal computer, then the smartphone.
\end{itemize}

\begin{tarif}
    \term{Moore's law}{loi de Moore}: in 1965, Gordon Moore (a founder of Intel) noticed that the number
    of transistors on a chip \textbf{doubles about every two years}. This stayed true for more than fifty
    years: from 2\,300 transistors in the 4004 to tens of billions today.
\end{tarif}

\begin{hal}
    ENIAC filled a whole room and used 150 kilowatts. Your phone is millions of times more powerful and runs
    on a small battery. The reason: transistors became tiny, so one processor now holds billions of them.
\end{hal}

%================================================
\section{Quick reminder: how does a computer count?}
%================================================

A computer only knows two electric states: \textbf{current flows} or \textbf{current does not flow}.
We write them 1 and 0.

\begin{tarif}
    \begin{itemize}
        \item A \term{bit}{bit}: the smallest unit of information. Its value is 0 or 1.
        \item A \term{byte}{octet}: a group of \textbf{8 bits}.
        \item A \term{word}{mot}: the number of bits the processor handles at once (16 bits in the 8086).
    \end{itemize}
\end{tarif}

We write numbers in three systems: decimal (base 10), binary (base 2), and hexadecimal (base 16, its
digits are 0 to 9, then A to F).

\begin{center}
\small
\begin{tabular}{|c|c|c||c|c|c|}
    \hline
    \ras{Decimal} & \ras{Binary} & \ras{Hex} & \ras{Decimal} & \ras{Binary} & \ras{Hex} \\ \hline
    0 & 0000 & 0 & 8  & 1000 & 8 \\ \hline
    1 & 0001 & 1 & 9  & 1001 & 9 \\ \hline
    2 & 0010 & 2 & 10 & 1010 & A \\ \hline
    3 & 0011 & 3 & 11 & 1011 & B \\ \hline
    4 & 0100 & 4 & 12 & 1100 & C \\ \hline
    5 & 0101 & 5 & 13 & 1101 & D \\ \hline
    6 & 0110 & 6 & 14 & 1110 & E \\ \hline
    7 & 0111 & 7 & 15 & 1111 & F \\ \hline
\end{tabular}
\end{center}

\begin{khatwa}[: convert \hex{2F} to binary and to decimal]
    \begin{enumerate}
        \item Each hex digit is \textbf{4 bits}: $2 = 0010$ and $F = 1111$.
        \item Put them side by side: $\hex{2F} = \bin{0010\,1111}$.
        \item To decimal: $2 \times 16 + 15 = 32 + 15 = 47$.
    \end{enumerate}
\end{khatwa}

\begin{khatwa}[: convert the binary number 1011 to decimal]
    Each bit has a ``weight'', a power of 2. We start on the right with $2^0 = 1$:
    \begin{center}
    \begin{tabular}{|c|c|c|c|c|}
        \hline \ras{Bit} & 1 & 0 & 1 & 1 \\ \hline
        \ras{Weight} & $2^3 = 8$ & $2^2 = 4$ & $2^1 = 2$ & $2^0 = 1$ \\ \hline
        \ras{Bit $\times$ weight} & 8 & 0 & 2 & 1 \\ \hline
    \end{tabular}
    \end{center}
    We add: $8 + 0 + 2 + 1 = 11$. So $\bin{1011} = 11$.
\end{khatwa}

\begin{khatwa}[: convert the decimal number 13 to binary (divide by 2 again and again)]
    \begin{enumerate}
        \item $13 \div 2 = 6$, remainder \textbf{1}.
        \item $6 \div 2 = 3$, remainder \textbf{0}.
        \item $3 \div 2 = 1$, remainder \textbf{1}.
        \item $1 \div 2 = 0$, remainder \textbf{1}. We stop, because the quotient is 0.
        \item Read the remainders \textbf{from the last to the first}: \bin{1101}. Check: $8 + 4 + 0 + 1 = 13$.
    \end{enumerate}
\end{khatwa}

\begin{tashbih}
    Why use hexadecimal? Because binary numbers are long and tiring to read: \bin{1010\,1111\,0011\,0110}.
    Hexadecimal packs every 4 bits into one digit, so the number becomes \hex{AF36}. It is the
    ``short form'' programmers use to write binary quickly, like writing ``Dr.'' instead of ``Doctor''.
\end{tashbih}

\begin{intabih}
    In assembly language we write \code{h} after a hex number (\code{2Fh}), and \code{b} after a binary
    number (\code{00101111b}). A number without a letter is decimal.
\end{intabih}

\subsection{Memory size units}
\begin{center}
\begin{tabular}{|c|c|c|}
    \hline
    \ras{Unit} & \ras{Symbol} & \ras{Value} \\ \hline
    kilobyte & KB & $2^{10} = 1024$ bytes \\ \hline
    megabyte & MB & $2^{20} = 1\,048\,576$ bytes \\ \hline
    gigabyte & GB & $2^{30}$ bytes (about one billion) \\ \hline
    terabyte & TB & $2^{40}$ bytes \\ \hline
\end{tabular}
\end{center}

%================================================
\section{The main parts of a computer}
%================================================

Inside every computer, from the biggest server to the smallest smartwatch, we find three main parts:

\begin{center}
\begin{tikzpicture}
    \node[unit, minimum width=4cm, minimum height=2.6cm] (cpu) at (0,0) {};
    \node[font=\bfseries] at (0, 1) {Processor (CPU)};
    \node[unit=brandOrange, minimum width=1.6cm, font=\small] at (-0.95, 0) {UC\\control};
    \node[unit=brandGreen, minimum width=1.6cm, font=\small] at (0.95, 0) {UAL\\compute};
    \node[draw=brandDark, fill=white, font=\small, minimum width=3.4cm] at (0,-0.9) {registers};
    \node[unit=brandPurple, minimum width=3.4cm, minimum height=2.6cm] (mem) at (-6,0) {Main memory\\{\small RAM and ROM}};
    \node[unit=brandGreen, minimum width=3.4cm, minimum height=2.6cm] (io) at (6,0) {Input/output\\devices};
    \draw[sahmT, <->] (cpu) -- (mem);
    \draw[sahmT, <->] (cpu) -- (io);
\end{tikzpicture}
\end{center}

\begin{enumerate}
    \item The \term{processor}{processeur} (CPU): the ``brain'' of the computer. It reads the instructions and runs them.
    \item The \term{main memory}{mémoire centrale}: it keeps the program that is running now, and its data.
    \item The \term{input/output devices}{périphériques d'E/S}: they let the computer talk with the outside world.
\end{enumerate}

\begin{khatwa}[: the trip of a letter from the keyboard to the screen]
    You press the key \textbf{A} in a text editor:
    \begin{enumerate}
        \item The keyboard (\textbf{input}) sends the code of the key to the computer.
        \item The processor is interrupted (it stops its work for a moment), reads the code, and stores it in
              \textbf{memory}, inside the text you are writing.
        \item The text editor asks the processor to ``draw'' the letter: the processor computes the light dots that make the letter.
        \item These dots go to the graphics card, then to the screen (\textbf{output}).
    \end{enumerate}
    All this happens in a tiny part of a second, so the letter seems to appear at once.
\end{khatwa}

\begin{tarif}
    The \term{central unit}{unité centrale}: the processor plus the main memory. This is the part that really runs the programs.
\end{tarif}

%================================================
\section{The two basic models: Von Neumann and Harvard}
%================================================

\subsection{The Von Neumann model}

In 1945, the scientist John von Neumann proposed a simple and brilliant idea: we keep the
\textbf{program} and the \textbf{data} together in \textbf{the same memory}. Before that, changing the
program meant re-connecting wires by hand!

In Von Neumann's model, a computer has \textbf{five units}: input, output, memory, the arithmetic and
logic unit, and the control unit. The key idea is the \term{stored program}{programme enregistré}: instructions
are numbers, just like data, so we can keep them in the same memory, and read and change them the same way.

\begin{center}
\begin{tikzpicture}
    \node[unit, minimum width=4.2cm, minimum height=2.4cm] (cpu) at (0,0) {};
    \node[font=\bfseries] at (0,0.85) {Processor};
    \node[unit=brandOrange, font=\small, minimum width=1.7cm] at (-1,-0.15) {control unit};
    \node[unit=brandGreen, font=\small, minimum width=1.7cm] at (1,-0.15) {UAL};
    \node[unit=brandPurple, minimum width=4.2cm, minimum height=2.4cm, align=center] (mem) at (-7,0) {};
    \node[font=\bfseries] at (-7,0.85) {one memory};
    \node[draw=brandDark, fill=brandTeal!20, font=\small, minimum width=3.4cm] at (-7,0.1) {program (instructions)};
    \node[draw=brandDark, fill=brandOrange!20, font=\small, minimum width=3.4cm] at (-7,-0.6) {data};
    \draw[brandDark, line width=5pt, {Stealth[length=4mm]}-{Stealth[length=4mm]}] (-4.8,0) -- (-2.2,0);
    \node[font=\small, above] at (-3.5,0.2) {one bus};
    \node[unit=brandGreen, minimum width=2.6cm, minimum height=1.2cm] (io) at (5,0) {input and\\output};
    \draw[sahmT, <->] (cpu) -- (io);
\end{tikzpicture}
\end{center}

\begin{itemize}
    \item \textbf{Advantage}: simple and flexible. To change the program, we just load another program into memory.
    \item \textbf{Drawback}: one bus for instructions and data, so we cannot read an instruction and a data
          item at the same moment. This traffic jam is called the \term{Von Neumann bottleneck}{goulot de Von Neumann}.
\end{itemize}

\begin{tashbih}
    Imagine a library with \textbf{one narrow door}. The students who bring back books (the data) and the
    students who come for instructions (the program) all use the same door, one by one. Even if the staff
    inside is very fast, the door decides the speed. This is the bottleneck.
\end{tashbih}

\subsection{The Harvard model}
In this model we use \textbf{two separate memories}: one for instructions and one for data. Each one has its own bus.

\begin{center}
\begin{tikzpicture}
    \node[unit, minimum width=3.4cm, minimum height=2.2cm] (cpu) at (0,0) {Processor};
    \node[unit=brandTeal, minimum width=3.6cm, minimum height=1.6cm] (im) at (-6.5,0) {instruction memory\\{\small the program}};
    \node[unit=brandOrange, minimum width=3.6cm, minimum height=1.6cm] (dm) at (6.5,0) {data memory\\{\small the variables}};
    \draw[tealDeep, line width=4pt, -{Stealth[length=4mm]}] (im) -- (cpu) node[midway, above, font=\small, text=brandDark] {instruction bus};
    \draw[orangeDeep, line width=4pt, {Stealth[length=4mm]}-{Stealth[length=4mm]}] (cpu) -- (dm) node[midway, above, font=\small, text=brandDark] {data bus};
\end{tikzpicture}
\end{center}

\begin{itemize}
    \item \textbf{Advantage}: faster, because the processor reads the next instruction and a data item at the same time.
    \item \textbf{Drawback}: more complex and more expensive, and the size of each memory is fixed.
\end{itemize}

\begin{tashbih}
    In the Harvard library there are \textbf{two doors}: one for instructions and one for data. The two
    queues move at the same time, so people get in faster. But two doors cost more, and if one side is full
    we cannot use the free space of the other side.
\end{tashbih}

\subsection{Comparison and uses}
\begin{center}
\begin{tabularx}{\linewidth}{|>{\centering\arraybackslash}p{3cm}|C|C|}
    \hline
    \ras{} & \ras{Von Neumann} & \ras{Harvard} \\ \hline
    Memory   & one memory for program and data & two separate memories \\ \hline
    Buses    & one shared bus & two independent buses \\ \hline
    Speed    & slower (bottleneck) & faster (two reads at once) \\ \hline
    Complexity & simple and cheap & more complex, more expensive \\ \hline
    Used in  & personal computers, servers & microcontrollers (Arduino), signal processors (DSP) \\ \hline
\end{tabularx}
\end{center}

\begin{hal}
    The processors in your computer mix the two models. The main memory (RAM) is shared (Von Neumann), but
    inside the processor there are two small, fast, separate memories: one for instructions and one for data
    (Harvard). We call this \textbf{modified Harvard}. We come back to it in Chapter 5.
\end{hal}

%================================================
\section{The processor: role and parts}
%================================================

\begin{tarif}
    The \term{processor}{processeur, UCT} (CPU): the circuit that runs the instructions of a program, one after the other.
\end{tarif}

The processor has three main parts (we study them in detail in Chapter 2):

\begin{center}
\begin{tabularx}{\linewidth}{|>{\centering\arraybackslash}p{4.2cm}|R|R|}
    \hline
    \ras{Part} & \ras{Its job} & \ras{Analogy} \\ \hline
    control unit (UC) & reads the instruction, understands it, gives orders to the other parts & the orchestra conductor \\ \hline
    arithmetic and logic unit (UAL) & does arithmetic ($+$, $-$) and logic (AND, OR) & the calculator \\ \hline
    registers & very small and very fast memory inside the processor & the scrap paper on the desk \\ \hline
\end{tabularx}
\end{center}

\subsection{The clock and the speed of the processor}
The processor works to the rhythm of a \term{clock}{horloge} that sends regular pulses. At each pulse the
processor moves one small step forward in its work.

\begin{tarif}
    The \textbf{frequency} $f$: the number of pulses per second, measured in hertz (Hz).
    The \textbf{period} $T$: the length of one pulse, and $T = \dfrac{1}{f}$.
\end{tarif}

\begin{khatwa}[: a processor at 3 GHz]
    \begin{enumerate}
        \item $f = 3\ \text{GHz} = 3 \times 10^{9}$ pulses per second.
        \item $T = \dfrac{1}{3 \times 10^{9}} \approx 0.33 \times 10^{-9}$ second, so about 0.33 nanosecond.
        \item In the time you blink (about 0.1 second), this processor does 300 million cycles!
    \end{enumerate}
\end{khatwa}

\begin{intabih}
    The frequency alone does not decide the speed of a processor. What matters is the \textbf{number of
    instructions done per second}. A 3 GHz processor that needs 4 cycles per instruction runs
    $\frac{3 \times 10^9}{4} = 750$ million instructions per second. A 2 GHz processor that runs 2 instructions
    per cycle runs $4 \times 10^9$ instructions per second: it is faster, even with a lower frequency.
\end{intabih}

%================================================
\section{Memory}
%================================================

\subsection{What does memory look like inside?}
Imagine memory as a \textbf{big building full of boxes}. Each box has a number called the
\textbf{address}, and each box holds \textbf{one byte}.

\begin{center}
\begin{tikzpicture}[cellx/.style={draw=brandDark, fill=white, minimum width=2.4cm, minimum height=0.6cm, font=\codefont},
    addr/.style={font=\codefont\small, text=tealDeep}]
    \foreach \a/\v [count=\i from 0] in {0000/3A, 0001/FF, 0002/07, 0003/41, 0004/00} {
        \node[cellx] (c\i) at (0, -0.6*\i) {\v};
        \node[addr, left=0.2cm of c\i] {\a};
    }
    \node[font=\small\bfseries] at (-2.15, 0.6) {address};
    \node[font=\small\bfseries] at (0, 0.6) {content};
    \node[font=\small, text width=6.5cm, align=left] at (5.5,-1.2)
        {Address 0003 holds \hex{41}.\\
         The address never changes; the content can change.};
\end{tikzpicture}
\end{center}

\begin{tashbih}
    The address is like the \textbf{house number} in a street, and the content is like \textbf{the person
    who lives there}. The house number is fixed; the people living there can change.
\end{tashbih}

\subsection{Only two operations: read and write}
\begin{itemize}
    \item \textbf{Read}: the processor sends an address, and the memory gives back the content (the content does not change).
    \item \textbf{Write}: the processor sends an address and a value, and the memory stores the value instead of the old one.
\end{itemize}

\begin{mithal}
    Box 0002 holds \hex{07}. The processor runs ``write \hex{55} at 0002''. After that, the box holds
    \hex{55}, and \hex{07} is lost for ever. So we must be careful where we write: writing in the wrong place
    erases important data.
\end{mithal}

\begin{tarif}
    \term{Random Access Memory}{mémoire à accès aléatoire} (RAM): ``random access'' means we can reach
    \textbf{any box directly}, in the same time, without going through the boxes before it. This is the
    opposite of an old cassette tape, which you must wind until you reach the song.
\end{tarif}

\subsection{Types of memory}
\begin{center}
\begin{tikzpicture}[level distance=1.4cm,
    level 1/.style={sibling distance=7cm},
    level 2/.style={sibling distance=2.3cm},
    every node/.append style={unit, font=\small, minimum width=1.9cm},
    edge from parent/.style={draw, brandMuted, thick}]
    \node[unit=brandDark, fill=brandDark!10] {main memory}
        child {node[unit=brandOrange] {RAM\\{\scriptsize read and write}}
            child {node[unit=brandOrange] {SRAM\\{\scriptsize fast}}}
            child {node[unit=brandOrange] {DRAM\\{\scriptsize cheap}}}}
        child {node[unit=brandTeal] {ROM\\{\scriptsize read only}}
            child {node {PROM}}
            child {node {EPROM}}
            child {node {EEPROM\\{\scriptsize and Flash}}}};
\end{tikzpicture}
\end{center}

\begin{center}
\small
\begin{tabularx}{\linewidth}{|>{\centering\arraybackslash}p{2.2cm}|R|>{\centering\arraybackslash}p{2.4cm}|R|}
    \hline
    \ras{Type} & \ras{Description} & \ras{When power is off} & \ras{Used for} \\ \hline
    SRAM & transistors only, very fast and expensive & erased & cache memory inside the processor \\ \hline
    DRAM & a tiny capacitor per bit, must be ``refreshed'' all the time & erased & main memory (RAM sticks) \\ \hline
    ROM  & written once in the factory & kept & old boot programs \\ \hline
    EPROM & erased with ultraviolet light, then written again & kept & prototypes \\ \hline
    EEPROM / Flash & erased electrically & kept & BIOS/UEFI, USB keys, SSD \\ \hline
\end{tabularx}
\end{center}

\begin{tashbih}
    \textbf{SRAM} is like a light switch: it stays in its position (on or off) as long as there is power.
    \textbf{DRAM} is like a glass with a small hole: the water (the electric charge) leaks slowly, so we must
    fill it again thousands of times per second. This is called \textbf{refresh}. But the glass is much
    smaller than the switch, so we can put billions of them on a cheap memory stick.
\end{tashbih}

\begin{hal}
    Why does a computer need ROM or Flash? When you switch it on, the RAM is completely empty. The processor
    needs a first program to start with: the boot program (BIOS or UEFI), kept in a memory that is not erased.
    This program checks the hardware, then finds the operating system on the disk and loads it into RAM.
\end{hal}

\begin{tarif}
    A \term{volatile}{volatile} memory loses its content when the power is off (RAM).
    A \term{non-volatile}{non volatile} memory keeps it (ROM, Flash).
\end{tarif}

\subsection{Calculating the size of a memory}
If the processor has $n$ address lines (wires), it can write $2^n$ different addresses.

\begin{tarif}
    $\text{number of boxes} = 2^{n}$ \qquad and \qquad
    $\text{size} = 2^{n} \times \text{size of one box}$
\end{tarif}

\begin{khatwa}[: the memory size of the 8086]
    The 8086 has a \textbf{20-bit} address bus, and each box holds one byte.
    \begin{enumerate}
        \item Number of addresses: $2^{20}$.
        \item $2^{20} = 2^{10} \times 2^{10} = 1024 \times 1024 = 1\,048\,576$ addresses.
        \item Size $= 2^{20} \times 1$ byte $= 1$ megabyte (1 MB).
    \end{enumerate}
\end{khatwa}

\begin{khatwa}[: how many address lines for a 64 KB memory?]
    \begin{enumerate}
        \item $64\ \text{KB} = 64 \times 1024 = 2^{6} \times 2^{10} = 2^{16}$ bytes.
        \item So $2^{n} = 2^{16}$, and $n = 16$ lines.
    \end{enumerate}
\end{khatwa}

%================================================
\section{Buses}
%================================================

\begin{tarif}
    A \term{bus}{bus}: a group of parallel wires that carries information between the processor, the memory
    and the input/output devices. Each wire carries one bit at a time.
\end{tarif}

\begin{center}
\begin{tikzpicture}
    \node[unit, minimum width=2.8cm, minimum height=4cm] (cpu) at (-5,0) {Processor};
    \node[unit=brandPurple, minimum width=2.6cm, minimum height=1.2cm] (mem) at (0,2.6) {Memory};
    \node[unit=brandGreen, minimum width=2.6cm, minimum height=1.2cm] (io) at (4,2.6) {Input and\\output};
    \draw[brandOrange, line width=5pt] (-3.6,0.9) -- (6,0.9);
    \draw[brandTeal!80!black, line width=5pt] (-3.6,0) -- (6,0);
    \draw[brandMuted, line width=5pt] (-3.6,-0.9) -- (6,-0.9);
    \node[font=\small\bfseries, text=orangeDeep, anchor=west] at (6.1,0.9) {address bus};
    \node[font=\small\bfseries, text=tealDeep, anchor=west] at (6.1,0) {data bus};
    \node[font=\small\bfseries, text=brandMuted, anchor=west] at (6.1,-0.9) {control bus};
    \foreach \x in {0, 4} {
        \draw[brandOrange, line width=2pt, -{Stealth}] (\x-0.6,0.9) -- (\x-0.6,1.95);
        \draw[brandTeal!80!black, line width=2pt, {Stealth}-{Stealth}] (\x,0) -- (\x,1.95);
        \draw[brandMuted, line width=2pt, {Stealth}-{Stealth}] (\x+0.6,-0.9) -- (\x+0.6,1.95);
    }
\end{tikzpicture}
\end{center}

\begin{center}
\begin{tabularx}{\linewidth}{|>{\centering\arraybackslash}p{3cm}|>{\centering\arraybackslash}p{2.6cm}|R|}
    \hline
    \ras{Bus} & \ras{Direction} & \ras{What does it carry?} \\ \hline
    \textbf{address} bus & from the processor only & the number of the box the processor wants to read or write \\ \hline
    \textbf{data} bus & both ways & the value that is read or written \\ \hline
    \textbf{control} bus & both ways & signals such as read (RD), write (WR), clock, interrupt \\ \hline
\end{tabularx}
\end{center}

\begin{tashbih}
    Imagine you ask for a parcel at the post office. The \textbf{address} tells the clerk which shelf you mean
    (address bus). The \textbf{parcel} itself is the data (data bus). The words ``give me'' or ``take this''
    are the control signal (read or write).
\end{tashbih}

\begin{khatwa}[: the processor reads box 0003]
    \begin{enumerate}
        \item The processor puts the address 0003 on the \textbf{address bus}.
        \item It sends a \textbf{read} signal (RD) on the \textbf{control bus}.
        \item The memory puts the content of the box (\hex{41}) on the \textbf{data bus}.
        \item The processor takes the value \hex{41} and keeps it in one of its registers.
    \end{enumerate}
\end{khatwa}

\subsection{Bus width}
\begin{itemize}
    \item The \textbf{width of the address bus} decides the \textbf{maximum memory size}: $n$ lines give $2^n$ addresses.
    \item The \textbf{width of the data bus} decides \textbf{how many bits move at once}: a 16-bit bus moves two bytes at a time.
\end{itemize}

\begin{khatwa}[: moving 1000 bytes with different buses]
    \begin{enumerate}
        \item An \textbf{8-bit} data bus moves one byte at a time: we need 1000 transfers.
        \item A \textbf{16-bit} bus moves two bytes: $1000 \div 2 = 500$ transfers.
        \item A \textbf{64-bit} bus moves 8 bytes: $1000 \div 8 = 125$ transfers only.
    \end{enumerate}
    The wider the bus, the faster the data moves, just as more cars pass on an 8-lane motorway than on a one-lane road.
\end{khatwa}

\begin{center}
\begin{tabular}{|c|c|c|c|}
    \hline
    \ras{Processor} & \ras{Data bus} & \ras{Address bus} & \ras{Max memory} \\ \hline
    Intel 8086 & 16 bits & 20 bits & $2^{20}$ = 1 MB \\ \hline
    Intel 80386 & 32 bits & 32 bits & $2^{32}$ = 4 GB \\ \hline
    modern x86-64 & 64 bits & 48 bits (in practice) & $2^{48}$ = 256 TB \\ \hline
\end{tabular}
\end{center}

%================================================
\section{Exercises}
%================================================

\begin{tamrin}
    Convert these numbers: (a) \hex{3C} to decimal and binary. (b) The decimal number 200 to hexadecimal.
\end{tamrin}
\begin{hall}
    (a) $\hex{3C} = 3 \times 16 + 12 = 60$, and in binary $\bin{0011\,1100}$.\\
    (b) $200 = 12 \times 16 + 8$, and 12 is written C, so $200 = \hex{C8}$.
\end{hall}

\begin{tamrin}
    A processor has a 24-bit address bus, and each memory box holds one byte. What is the maximum memory size?
\end{tamrin}
\begin{hall}
    $2^{24} = 2^{4} \times 2^{20} = 16 \times 1\ \text{MB} = 16\ \text{MB}$.
\end{hall}

\begin{tamrin}
    We want a 4 GB memory. How many lines does the address bus need?
\end{tamrin}
\begin{hall}
    $4\ \text{GB} = 2^{2} \times 2^{30} = 2^{32}$ bytes, so we need \textbf{32 lines}.
\end{hall}

\begin{tamrin}
    True or false? Explain.
    (a) In the Von Neumann model, the program is kept in a different memory from the data.
    (b) The address bus carries information in both directions.
    (c) DRAM keeps its content after the computer is switched off.
\end{tamrin}
\begin{hall}
    (a) \textbf{False}: Von Neumann has one memory for program and data; separation is the Harvard model.\\
    (b) \textbf{False}: only the processor sends addresses, so it goes one way.\\
    (c) \textbf{False}: DRAM is volatile; it is erased when the power goes off.
\end{hall}

\begin{tamrin}
    A processor runs at 2 GHz, and each instruction needs 4 cycles on average. How many instructions does it run per second?
\end{tamrin}
\begin{hall}
    Cycles per second $= 2 \times 10^{9}$. Instructions $= \dfrac{2 \times 10^{9}}{4} = 0.5 \times 10^{9}$,
    that is \textbf{500 million instructions per second}.
\end{hall}

\chapter{Inside the Processor}

\begin{khulasa}[: what will we learn in this chapter?]
\begin{itemize}
    \item What is inside the processor, and the job of each part.
    \item Registers and their types.
    \item The arithmetic and logic unit: how the processor adds two numbers with logic gates, and what flags are.
    \item The control unit: who decides what happens at each moment.
    \item The instruction set and the addressing modes.
    \item The steps of every instruction: fetch, decode, execute, write back.
\end{itemize}
\end{khulasa}

%================================================
\section{A look inside the processor}
%================================================

Let us open the processor and look inside. The figure shows a simple processor, but it has all the main
ideas of real processors.

\begin{center}
\begin{tikzpicture}[r/.style={reg, minimum width=1.5cm},
    bus/.style={line width=3pt, color=brandTeal!70!black},
    link/.style={thick, color=brandTeal!70!black},
    sig/.style={dashed, thick, color=orangeDeep, -{Stealth[length=2mm]}}]
    \draw[brandTeal, very thick, rounded corners=8pt, fill=brandTeal!4] (-8.4,-3.3) rectangle (2.2,3.4);
    \node[font=\bfseries, text=tealDeep, anchor=north east] at (2.1,3.35) {Processor};
    \draw[bus] (-7.6,0) -- (1.8,0);
    \node[font=\scriptsize, text=tealDeep] at (0.2,-0.3) {internal bus};
    % control unit
    \draw[orangeDeep, very thick, rounded corners=5pt, fill=brandOrange!12] (-8.1,0.5) rectangle (-4.9,3.0);
    \node[font=\small\bfseries, text=orangeDeep] at (-6.5,2.7) {control unit (UC)};
    \node[r] (ri) at (-6.5,2.05) {RI};
    \node[reg, font=\scriptsize, minimum width=1.25cm] (dec) at (-7.3,1.05) {decoder};
    \node[reg, font=\scriptsize, minimum width=1.25cm] (seq) at (-5.7,1.05) {sequencer};
    \draw[sahm] (ri) -- (dec);
    \draw[sahm] (dec) -- (seq);
    \draw[link] (ri.east) -- ++(0.35,0) |- (-4.55,0.35) -- (-4.55,0);
    \node[unit=brandPurple, font=\small] (clk) at (-5.8,3.95) {clock};
    \draw[sahm] (clk) -- (-5.8,3.0);
    % registers above the bus
    \node[r] (co) at (-2.6,2.3) {CO};
    \node[r] (ram) at (0.6,2.3) {RAM};
    \node[r] (rim) at (0.6,1.1) {RIM};
    \draw[sahmT] (co) -- node[above, font=\scriptsize, text=brandDark] {address} (ram);
    \node[font=\scriptsize, text=brandMuted, anchor=west] at (-2.45,1.65) {+1 per fetch};
    \draw[link] (co.south) ++(0,-0.15) -- (-2.6,0);
    \draw[link] (rim.south) -- (0.6,0);
    % ALU with two inputs, ACC under it
    \draw[brandGreen!70!black, very thick, fill=brandGreen!20]
        (-5.4,-0.8) -- (-3.6,-0.8) -- (-4.0,-1.9) -- (-5.0,-1.9) -- cycle;
    \node[font=\bfseries] at (-4.5,-1.35) {UAL};
    \node[r] (acc) at (-4.5,-2.75) {ACC};
    \node[r, font=\small] (fl) at (-2.3,-2.75) {flags};
    \draw[sahmT] (-4.5,-1.9) -- (acc.north) node[pos=0.5, right, font=\scriptsize, text=brandDark] {result};
    \draw[sahmT] (-3.8,-1.6) -| (fl.north);
    \draw[sahm] (-3.9,0) -- (-3.9,-0.8) node[pos=0.5, right, font=\scriptsize, text=brandDark] {operand 2};
    \draw[sahm] (acc.west) -- ++(-0.9,0) |- (-6.0,-0.45) -| (-5.1,-0.8);
    \node[font=\scriptsize] at (-6.75,-1.6) {operand 1};
    \draw[link] (acc.south) -- (-4.5,-3.2) -- (-1.0,-3.2) -- (-1.0,0);
    % control signals
    \draw[sig] (seq.south) -- (-5.7,0.15);
    \node[font=\scriptsize, text=orangeDeep] at (-7.3,0.25) {control signals};
    % memory
    \node[unit=brandPurple, minimum width=2cm, minimum height=4cm] (mem) at (4.4,1.2) {main\\memory};
    \draw[sahmO] (ram.east) -- node[above, font=\scriptsize, text=brandDark] {address bus} ([yshift=1.1cm]mem.west);
    \draw[sahmT, {Stealth[length=2.6mm]}-{Stealth[length=2.6mm]}] (rim.east) -- node[below, font=\scriptsize, text=brandDark] {data bus} ([yshift=-0.1cm]mem.west);
\end{tikzpicture}
\end{center}

\begin{center}
\small
\begin{tabularx}{\linewidth}{|>{\centering\arraybackslash}p{4cm}|R|}
    \hline
    \ras{Part} & \ras{Its job in one sentence} \\ \hline
    CO (program counter) & keeps the address of the \textbf{next instruction} \\ \hline
    RI (instruction register) & keeps the \textbf{current instruction} being run \\ \hline
    RAM (memory address register) & keeps the address we send to memory \\ \hline
    RIM (memory data register) & keeps the data coming from memory or going to it \\ \hline
    ACC (accumulator) & keeps the result of the calculation \\ \hline
    flags & keep information about the last result: is it zero? negative? \dots \\ \hline
    UAL & does the arithmetic and logic operations \\ \hline
    decoder and sequencer & understand the instruction and send the orders in the right order \\ \hline
\end{tabularx}
\end{center}

\begin{intabih}
    Do not mix up RAM the memory (\emph{Random Access Memory}) and RAM the register (\emph{Registre d'Adresse
    Mémoire}) used in French books. In English this register is called \textbf{MAR}, RIM is called \textbf{MDR},
    and CO is called \textbf{PC} (\emph{Program Counter}). In the 8086 it is called \textbf{IP}.
    In this course we keep the French short names (CO, RI, RAM, RIM, UAL, UC) because they are used in the syllabus.
\end{intabih}

%================================================
\section{Registers}
%================================================

\begin{tarif}
    A \term{register}{registre}: a very small memory box (8, 16, 32 or 64 bits) \textbf{inside the processor}.
    It is the fastest memory in the whole computer.
\end{tarif}

\begin{tashbih}
    When you solve a maths exercise, you do not walk to the library each time you need a number! You write
    the numbers you need on \textbf{scrap paper} in front of you. Registers are the processor's scrap paper,
    and memory is the library.
\end{tashbih}

\subsection{Why are registers important for speed?}
\begin{center}
\begin{tabular}{|c|c|c|}
    \hline
    \ras{Place} & \ras{Access time (about)} & \ras{If one cycle took one second\dots} \\ \hline
    register & less than 1 ns & one second \\ \hline
    main memory (RAM) & about 80 ns & about a minute and a half \\ \hline
    SSD disk & about 100 µs & more than a day \\ \hline
\end{tabular}
\end{center}
So we keep in registers the data we use often, and we go to memory as little as possible.

\subsection{Types of registers}
\begin{enumerate}
    \item \term{Data registers}{registres de données}: they keep the numbers we compute with (like AX, BX in the 8086).
    \item \term{Address registers}{registres d'adresse}: they keep memory addresses (like SI, DI, SP).
    \item \term{Control and status registers}{registres de contrôle}: they help the processor organise its work: CO, RI, and the flags register.
\end{enumerate}

\begin{mithal}[: registers at work]
    The processor is now running the instruction at address 11, ``add ACC and the content of box 21''.
    \begin{itemize}
        \item CO holds 12: the address of the \textbf{next} instruction, not the current one.
        \item RI holds the current instruction, so the control unit can read it during the whole execution.
        \item RAM holds 21: the address we are sending to memory now.
        \item RIM receives the content of box 21 coming from memory.
        \item ACC holds the first number, and will receive the result.
    \end{itemize}
    Each register has a clear job. This is why we call them special registers, unlike the general registers
    that the programmer uses freely.
\end{mithal}

\subsection{The flags register}
\begin{tarif}
    A \term{flag}{indicateur, drapeau}: one bit that the arithmetic unit sets to 1 or 0 after each
    operation, depending on the result. We use flags later to make decisions (for example: ``if the result is zero, jump'').
\end{tarif}

\begin{center}
\begin{tabularx}{\linewidth}{|>{\centering\arraybackslash}p{2.6cm}|R|R|}
    \hline
    \ras{Flag} & \ras{Becomes 1 when\dots} & \ras{Example on 8 bits} \\ \hline
    ZF (zero) & the result is \textbf{zero} & $5 - 5 = 0$ \\ \hline
    CF (carry) & a carry comes out of the last bit (the result is too big) & $200 + 100 = 300 > 255$ \\ \hline
    SF (sign) & the last bit of the result is 1 (negative number) & $3 - 7 = -4$ \\ \hline
    OF (overflow) & the result is wrong if we read the numbers as signed & $100 + 50 = 150 > 127$ \\ \hline
\end{tabularx}
\end{center}

\subsection{The difference between CF and OF}
The same byte can be read in two ways, and the processor does not know which one the programmer means:
\begin{itemize}
    \item \textbf{unsigned}: values from 0 to 255.
    \item \textbf{signed}: values from $-128$ to $+127$. The last bit (bit 7) is the sign: 1 means negative.
\end{itemize}
So the processor computes both flags. CF matters for unsigned numbers, and OF matters for signed numbers.

\begin{khatwa}[: $100 + 50$ on 8 bits]
    \begin{enumerate}
        \item $100 = \bin{0110\,0100}$ and $50 = \bin{0011\,0010}$.
        \item The sum: $\bin{1001\,0110} = 150$ (unsigned). No carry out of bit 7, so CF = 0: correct for unsigned numbers.
        \item But bit 7 is now 1! Reading it as signed, we see a negative number: $150 - 256 = -106$.
        \item We added two positive numbers and got a negative one: impossible, so the signed result is wrong, and OF = 1.
    \end{enumerate}
\end{khatwa}

\begin{intabih}
    An easy rule for OF: overflow happens when we add two numbers \textbf{with the same sign} and get a result
    \textbf{with the other sign}. Adding a positive and a negative number never overflows.
\end{intabih}

%================================================
\section{The arithmetic and logic unit (UAL)}
%================================================

\begin{tarif}
    The \term{arithmetic and logic unit}{unité arithmétique et logique, UAL} (ALU in English): the part
    that does arithmetic (add, subtract, compare\dots) and logic (AND, OR, XOR, NOT, shifts).
\end{tarif}

\begin{center}
\begin{tikzpicture}
    \draw[brandGreen!70!black, very thick, fill=brandGreen!20]
        (-2.2,1) -- (-0.4,1) -- (0,0.4) -- (0.4,1) -- (2.2,1) -- (1.2,-1) -- (-1.2,-1) -- cycle;
    \node[font=\bfseries\large] at (0,-0.2) {UAL};
    \draw[sahmT] (-1.3,2) node[above, font=\small] {operand A} -- (-1.3,1);
    \draw[sahmT] (1.3,2) node[above, font=\small] {operand B} -- (1.3,1);
    \draw[sahmO] (-3.6,0) node[left, font=\small] {operation (from the control unit)} -- (-1.7,0);
    \draw[sahmT] (0,-1) -- (0,-2) node[below, font=\small] {result};
    \draw[sahm] (1.6,-0.2) -- (3.4,-0.2) node[right, font=\small] {flags ZF, CF, SF, OF};
\end{tikzpicture}
\end{center}

\subsection{How do logic gates add? The half adder}
Let us add two bits $A$ and $B$. There are only four cases:

\begin{center}
\begin{tabular}{|c|c||c|c|}
    \hline
    \ras{$A$} & \ras{$B$} & \ras{sum $S$} & \ras{carry $C$} \\ \hline
    0 & 0 & 0 & 0 \\ \hline
    0 & 1 & 1 & 0 \\ \hline
    1 & 0 & 1 & 0 \\ \hline
    1 & 1 & 0 & \cell{orangeDeep}{\textbf{1}} \\ \hline
\end{tabular}
\end{center}

We see that $S$ is 1 when exactly one bit is 1: this is the \textbf{XOR} gate.
And $C$ is 1 when both bits are 1: this is the \textbf{AND} gate.

\begin{tarif}
    \term{Half adder}{demi-additionneur}: \quad $S = A \oplus B$ \quad and \quad $C = A \cdot B$
\end{tarif}

\begin{center}
\begin{tikzpicture}[gate/.style={draw=brandDark, fill=white, thick, minimum width=1.5cm,
        minimum height=0.9cm, font=\codefont\bfseries, rounded corners=2pt}]
    \node[gate] (x) at (0,0.8) {XOR};
    \node[gate] (a) at (0,-0.8) {AND};
    \node[font=\bfseries] (A) at (-3.2,0.4) {A};
    \node[font=\bfseries] (B) at (-3.2,-0.4) {B};
    \fill (-2,0.4) circle (2pt); \fill (-1.5,-0.4) circle (2pt);
    \draw[thick] (A) -- (-2,0.4) |- ([yshift=0.15cm]x.west);
    \draw[thick] (-2,0.4) |- ([yshift=0.15cm]a.west);
    \draw[thick] (B) -- (-1.5,-0.4) |- ([yshift=-0.15cm]x.west);
    \draw[thick] (-1.5,-0.4) |- ([yshift=-0.15cm]a.west);
    \draw[sahmT] (x.east) -- (2,0.8) node[right] {S};
    \draw[sahmO] (a.east) -- (2,-0.8) node[right] {C};
\end{tikzpicture}
\end{center}

\subsection{The full adder}
To add numbers with many bits, we must also add the \textbf{incoming carry} from the previous bit.
A full adder has three inputs: $A$, $B$ and the carry in $C_{in}$.

\begin{tarif}
    \term{Full adder}{additionneur complet}: \quad
    $S = A \oplus B \oplus C_{in}$ \quad and \quad
    $C_{out} = A \cdot B + C_{in} \cdot (A \oplus B)$
\end{tarif}

\begin{center}
\small
\begin{tabular}{|c|c|c||c|c|}
    \hline
    \ras{$A$} & \ras{$B$} & \ras{$C_{in}$} & \ras{$S$} & \ras{$C_{out}$} \\ \hline
    0&0&0&0&0 \\ \hline 0&0&1&1&0 \\ \hline 0&1&0&1&0 \\ \hline 0&1&1&0&1 \\ \hline
    1&0&0&1&0 \\ \hline 1&0&1&0&1 \\ \hline 1&1&0&0&1 \\ \hline 1&1&1&1&1 \\ \hline
\end{tabular}
\end{center}

\subsection{A 4-bit adder}
We put 4 full adders side by side: the carry out of each one goes into the next one, exactly as we do
when we add with pen and paper.

\begin{center}
\begin{tikzpicture}[fa/.style={draw=brandTeal, fill=brandTeal!15, very thick, minimum width=1.6cm,
        minimum height=1.1cm, font=\codefont\bfseries, rounded corners=3pt}]
    \foreach \i/\x in {0/0, 1/-3, 2/-6, 3/-9} {
        \node[fa] (f\i) at (\x,0) {FA};
        \draw[thick] ([xshift=-0.35cm]f\i.north) -- ++(0,0.6) node[above, font=\small] {A\i};
        \draw[thick] ([xshift=0.35cm]f\i.north) -- ++(0,0.6) node[above, font=\small] {B\i};
        \draw[sahmT] (f\i.south) -- ++(0,-0.6) node[below, font=\small] {S\i};
    }
    \draw[sahmO] (1.7,0) node[right, font=\small] {0} -- (f0.east);
    \draw[sahmO] (f0.west) -- (f1.east);
    \draw[sahmO] (f1.west) -- (f2.east);
    \draw[sahmO] (f2.west) -- (f3.east);
    \draw[sahmO] (f3.west) -- ++(-0.9,0) node[left, font=\small] {CF};
\end{tikzpicture}
\end{center}
Bit 0 (the smallest) is on the right, like when we write numbers. The carry travels from right to left.

\begin{khatwa}[: add $0111 + 0110$ (that is $7 + 6$) with the adder]
    \begin{center}
    \begin{tabular}{|c|c|c|c|c|c|}
        \hline
        \ras{Bit} & \ras{$A$} & \ras{$B$} & \ras{$C_{in}$} & \ras{$S$} & \ras{$C_{out}$} \\ \hline
        0 & 1 & 0 & 0 & 1 & 0 \\ \hline
        1 & 1 & 1 & 0 & 0 & 1 \\ \hline
        2 & 1 & 1 & 1 & 1 & 1 \\ \hline
        3 & 0 & 0 & 1 & 1 & 0 \\ \hline
    \end{tabular}
    \end{center}
    Read $S$ from bit 3 down to bit 0: $1101 = 13$, which is indeed $7 + 6$. The last carry is 0, so CF = 0.
\end{khatwa}

\subsection{And subtraction? The same adder!}
The processor does not need a special circuit to subtract. It uses the \term{two's complement}{complément à deux}:
$A - B = A + (-B)$, and we get $-B$ by flipping all the bits of $B$, then adding 1.

\begin{khatwa}[: $7 - 5$ on 4 bits]
    \begin{enumerate}
        \item $5 = 0101$. Flip the bits: $1010$. Add 1: $1011$. So $-5 = 1011$.
        \item Add: $0111 + 1011 = 1\,0010$. We drop the carry out, so $0010 = 2$ is left.
        \item The result: $7 - 5 = 2$. Correct! The same circuit added and subtracted.
    \end{enumerate}
\end{khatwa}

\subsection{Comparing by subtracting}
The processor does not need a special circuit to compare either: to compare $A$ and $B$ it computes
$A - B$ \textbf{without keeping the result}, then looks at the flags:
\begin{itemize}
    \item if ZF = 1 then $A = B$.
    \item if CF = 1 (a ``borrow'' happened) then $A < B$ (for unsigned numbers).
\end{itemize}

\subsection{A real test: how do the flags change?}
We ran the program below on the 8086. Do not worry if you do not understand all the instructions yet; we
study them in Chapter 3. Just look at the flags.

\asm{ch2/flags.asm}

\begin{center}
\small
\begin{tabular}{|c|l|c|c|c|c|c|}
    \hline
    \ras{Line} & \ras{Instruction} & \ras{AL} & \ras{CF} & \ras{ZF} & \ras{SF} & \ras{OF} \\ \hline
    4 & \code{add al, 100} ($200 + 100$) & $\hex{2C} = 44$ & \cell{orangeDeep}{\textbf{1}} & 0 & 0 & 0 \\ \hline
    5 & \code{mov al, 100} & $\hex{64} = 100$ & 1 & 0 & 0 & 0 \\ \hline
    6 & \code{add al, 50} ($100 + 50$) & $\hex{96} = 150$ & 0 & 0 & 1 & \cell{orangeDeep}{\textbf{1}} \\ \hline
    8 & \code{sub al, 5} ($5 - 5$) & $\hex{00} = 0$ & 0 & \cell{orangeDeep}{\textbf{1}} & 0 & 0 \\ \hline
    10 & \code{sub al, 7} ($3 - 7$) & $\hex{FC} = 252$ & \cell{orangeDeep}{\textbf{1}} & 0 & \cell{orangeDeep}{\textbf{1}} & 0 \\ \hline
\end{tabular}
\end{center}

\begin{khatwa}[: why did $200 + 100$ give 44?]
    \begin{enumerate}
        \item $200 + 100 = 300$, but AL holds only 8 bits, that is from 0 to 255.
        \item $300 - 256 = 44$: AL keeps 44, and the ``256'' that did not fit goes into the flag CF, so CF = 1.
    \end{enumerate}
\end{khatwa}

\begin{intabih}
    Look at line 5: the instruction \code{mov} \textbf{does not change the flags}, so CF stayed 1. Only
    arithmetic and logic instructions change the flags.
\end{intabih}

%================================================
\section{The control unit (UC)}
%================================================

\begin{tarif}
    The \term{control unit}{unité de commande, UC}: the part that reads the instructions and understands
    them, then sends \textbf{signals} to the other parts at each clock pulse (open this register, add, read
    memory\dots) in the right order.
\end{tarif}

\begin{tashbih}
    The control unit is the \textbf{orchestra conductor}: it does not play any instrument itself, but it tells
    each musician when to start and when to stop. Without it, everybody plays at the same time and it becomes noise.
\end{tashbih}

Its main parts:
\begin{itemize}
    \item the \textbf{instruction register} (RI): keeps the current instruction.
    \item the \term{decoder}{décodeur}: reads the operation code (opcode) and finds what is asked: add? move? jump?
    \item the \term{sequencer}{séquenceur}: sends the right signals, one after the other, to the rhythm of the clock.
\end{itemize}

\subsection{Micro-operations, pulse by pulse}
Each step of an instruction is split into very small steps called \term{micro-operations}{micro-opérations}.
At each clock pulse, the sequencer sends signals that open ``doors'' between registers. Here is how an instruction is fetched:

\begin{center}
\small
\begin{tabularx}{\linewidth}{|c|R|R|}
    \hline
    \ras{Pulse} & \ras{Micro-operation} & \ras{Signal sent by the sequencer} \\ \hline
    T1 & \nql{RAM}{CO} & ``copy CO into RAM'' \\ \hline
    T2 & \nql{RIM}{[RAM]}, \nql{CO}{CO + 1} & ``read memory'' and ``increase CO'' \\ \hline
    T3 & \nql{RI}{RIM} & ``copy RIM into RI'' \\ \hline
    T4\dots & depends on the instruction & the decoder chooses the next signals \\ \hline
\end{tabularx}
\end{center}
The fetch is the same for all instructions; what comes after depends on the opcode.

\subsection{Two kinds of control units}
\begin{center}
\begin{tabularx}{\linewidth}{|>{\centering\arraybackslash}p{3.2cm}|R|R|}
    \hline
    \ras{} & \ras{Hardwired (câblée)} & \ras{Microprogrammed (microprogrammée)} \\ \hline
    Idea & a fixed logic circuit makes the signals & each instruction is a small program stored in an internal memory \\ \hline
    Speed & very fast & a little slower \\ \hline
    Changes & hard (we must change the circuit) & easy (we change the microprogram) \\ \hline
    Found in & RISC processors & CISC processors like the 8086 \\ \hline
\end{tabularx}
\end{center}

%================================================
\section{The instruction set}
%================================================

\begin{tarif}
    The \term{instruction set}{jeu d'instructions}: the list of all the instructions a processor can run.
    Each processor family has its own set.
\end{tarif}

\subsection{The shape of an instruction}
Every machine-language instruction has two parts:

\begin{center}
\begin{tikzpicture}
    \node[draw=brandOrange, fill=brandOrange!20, thick, minimum width=3.5cm, minimum height=1cm, font=\bfseries] (op) at (0,0) {opcode};
    \node[draw=brandTeal, fill=brandTeal!20, thick, minimum width=5cm, minimum height=1cm, font=\bfseries, anchor=west] (oper) at (op.east) {operands};
    \node[below=0.15cm of op, font=\small] {what do we do?};
    \node[below=0.15cm of oper, font=\small] {on what do we work?};
    \node[above=0.2cm of op, font=\codefont] {ADD};
    \node[above=0.2cm of oper, font=\codefont] {AX, 5};
\end{tikzpicture}
\end{center}

\begin{tarif}
    \term{Assembly language}{langage d'assemblage}: we write instructions with short words that humans
    understand (\code{ADD}, \code{MOV}) instead of binary numbers. A program called the
    \term{assembler}{assembleur} turns them into machine language (zeros and ones).
\end{tarif}

\begin{mithal}
    On the 8086, the instruction \code{add al, 5} is written in machine language with the two bytes 04 05, that is
    0000\,0100 0000\,0101: the first byte is the opcode ``add an immediate value to AL'', and the second one is the
    number 5. The programmer writes \code{add al, 5}, and the assembler writes the zeros and ones (we see this in
    detail in Chapter 3).
\end{mithal}

\subsection{How many instructions can a processor have?}
If the opcode uses $k$ bits, there are $2^k$ different codes, so at most $2^k$ instructions.

\begin{khatwa}
    A 6-bit opcode: $2^6 = 64$ possible instructions. An 8-bit opcode: $2^8 = 256$ instructions.
\end{khatwa}

\subsection{Number of addresses in an instruction}
Let us compute $A = B + C$ on four kinds of machines. The fewer operands in an instruction, the shorter the
instruction, but the more instructions we need:

\begin{center}
\small
\begin{tabularx}{\linewidth}{|>{\centering\arraybackslash}p{3.6cm}|X|>{\centering\arraybackslash}p{3cm}|}
    \hline
    \ras{Kind of machine} & \ras{Program} & \ras{Instructions} \\ \hline
    3 addresses & \code{ADD A, B, C} & 1 \\ \hline
    2 addresses & \code{MOV A, B} then \code{ADD A, C} & 2 \\ \hline
    1 address (accumulator) & \code{LOAD B} then \code{ADD C} then \code{STORE A} & 3 \\ \hline
    0 address (stack) & \code{PUSH B}, \code{PUSH C}, \code{ADD}, \code{POP A} & 4 \\ \hline
\end{tabularx}
\end{center}
The 8086 is a \textbf{2-address} machine: \code{add ax, bx} means $\text{AX} \leftarrow \text{AX} + \text{BX}$.

\subsection{Classes of instructions}
\begin{center}
\begin{tabularx}{\linewidth}{|>{\centering\arraybackslash}p{3.5cm}|R|>{\centering\arraybackslash}p{4.2cm}|}
    \hline
    \ras{Class} & \ras{Job} & \ras{Examples (8086)} \\ \hline
    data transfer & copy a value from one place to another & \code{MOV}, \code{PUSH}, \code{POP} \\ \hline
    arithmetic & add, subtract, multiply, divide & \code{ADD}, \code{SUB}, \code{MUL}, \code{INC} \\ \hline
    logic & operations on bits & \code{AND}, \code{OR}, \code{XOR}, \code{SHL} \\ \hline
    flow control & jump and change the order of execution & \code{JMP}, \code{JE}, \code{CALL}, \code{RET} \\ \hline
    input/output & talk with devices & \code{IN}, \code{OUT}, \code{INT} \\ \hline
\end{tabularx}
\end{center}

\subsection{CISC and RISC}
\begin{itemize}
    \item \textbf{CISC} (Complex Instruction Set Computer): many complex instructions; some do a lot of work. Example: Intel x86.
    \item \textbf{RISC} (Reduced Instruction Set Computer): few, simple, fast instructions, all of the same length.
          Example: ARM in phones, RISC-V.
\end{itemize}

%================================================
\section{Addressing modes}
%================================================

\begin{tarif}
    An \term{addressing mode}{mode d'adressage}: the way we say \textbf{where the data is} that the instruction works on.
\end{tarif}

\begin{tashbih}
    You want to give a book to a friend:
    \textbf{(1)} you put the book in his hand (immediate),
    \textbf{(2)} you say: ``the book is on shelf number 12'' (direct),
    \textbf{(3)} you say: ``look at the paper in my pocket, it has the shelf number'' (indirect),
    \textbf{(4)} you say: ``the third shelf after shelf number 10'' (indexed).
\end{tashbih}

\subsection{The basic modes}
\begin{center}
\small
\begin{tabularx}{\linewidth}{|>{\centering\arraybackslash}p{2.4cm}|>{\centering\arraybackslash}p{3.3cm}|R|}
    \hline
    \ras{Mode} & \ras{8086 example} & \ras{Where is the data?} \\ \hline
    immediate & \code{mov ax, 5} & inside the instruction itself \\ \hline
    register & \code{mov bx, ax} & in a register \\ \hline
    direct & \code{mov cx, [val]} & in memory; its address is written in the instruction \\ \hline
    indirect & \code{mov dx, [si]} & in memory; its address is in a register \\ \hline
    indexed & \code{mov ax, [tab+di]} & address = start address + value of an index register \\ \hline
    relative & \code{jmp next} & address = current address + offset (for jumps) \\ \hline
\end{tabularx}
\end{center}

\subsection{When do we use each mode?}
\begin{itemize}
    \item \textbf{Immediate}: for constants (number of loop turns, start value). The fastest mode, because the value
          came with the instruction, so we do not go to memory. But we cannot change it while the program runs.
    \item \textbf{Register}: for fast calculations between values already inside the processor. Very fast too.
    \item \textbf{Direct}: for variables that have a fixed place in memory (like a global variable). One memory read.
    \item \textbf{Indirect}: when we do not know the address when we write the program, but compute it while it runs.
          It is the basis of \textbf{pointers} in C.
    \item \textbf{Indexed}: for arrays. We fix the start address and change the index at each turn of the loop, to visit every element.
    \item \textbf{Relative}: for jumps. Because the distance is relative, the program works even if it is loaded at another place in memory.
\end{itemize}

\begin{khatwa}[: how many times do we go to memory?]
    \begin{center}
    \begin{tabular}{|l|c|}
        \hline \ras{Instruction} & \ras{Memory reads to get the data} \\ \hline
        \code{mov ax, 5} (immediate) & 0 \\ \hline
        \code{mov ax, bx} (register) & 0 \\ \hline
        \code{mov ax, [val]} (direct) & 1 \\ \hline
        \code{mov ax, [si]} (indirect with a register) & 1 \\ \hline
        memory indirect (the address itself is in memory) & 2 \\ \hline
    \end{tabular}
    \end{center}
    Every trip to memory costs time, so programmers try to keep important data in registers.
\end{khatwa}

\subsection{The modes in a real program}
\asm{ch2/modes.asm}

\begin{center}
\begin{tikzpicture}[cellx/.style={draw=brandDark, fill=white, minimum width=2cm, minimum height=0.55cm, font=\codefont\small},
    addr/.style={font=\codefont\small, text=tealDeep}]
    \foreach \a/\v/\n [count=\i from 0] in {0117/1234/val, 0119/000A/tab, 011B/0014/{}, 011D/001E/{}} {
        \node[cellx] (c\i) at (0, -0.55*\i) {\v};
        \node[addr, left=0.15cm of c\i] {\a};
        \node[font=\codefont\small, right=0.15cm of c\i] {\n};
    }
    \node[font=\small\bfseries] at (0,0.55) {memory};
    \draw[sahmO] (-2.4,-0.55) -- (-2.4,-1.65) node[midway, left, font=\scriptsize] {+4};
\end{tikzpicture}
\end{center}

\begin{sharh}
    \s{1}{A program of type .COM: it starts at address \hex{100}.}
    \s{3}{\textbf{Immediate}: the number 5 is written inside the instruction. Result: AX = 0005.}
    \s{4}{\textbf{Register}: copy AX into BX. Result: BX = 0005.}
    \s{5}{\textbf{Direct}: read memory at the address of \code{val} (\hex{0117}). Result: CX = 1234.}
    \s{6}{\code{lea} puts the \textbf{address} of \code{val} (not its value) into SI: SI = 0117.}
    \s{7}{\textbf{Indirect}: read memory at the address held in SI. Result: DX = 1234.}
    \s{8}{DI = 4: each element of \code{tab} is a word (two bytes), so the third element is 4 bytes away.}
    \s{9}{\textbf{Indexed}: address $= \hex{0119} + 4 = \hex{011D}$, which holds 30. Result: AX = 001E.}
    \s{10}{End of the program, back to the system.}
\end{sharh}

\begin{intabih}
    The brackets \code{[ ]} always mean: ``the \textbf{content of memory} at this address''. \code{mov dx, si}
    copies the number \hex{0117} itself, but \code{mov dx, [si]} reads what is stored at address \hex{0117}.
\end{intabih}

%================================================
\section{The steps of an instruction}
%================================================

Every instruction, even a simple one, goes through the same steps. This is called the
\term{instruction cycle}{cycle d'instruction}.

\begin{center}
\begin{tikzpicture}
    \foreach \a/\t/\c in {90/{1. Fetch}/brandTeal, 0/{2. Decode}/brandOrange,
                          -90/{3. Execute}/brandGreen, 180/{4. Write back}/brandPurple} {
        \node[unit=\c, minimum width=2.9cm, minimum height=1.1cm] (n\a) at (\a:2.6) {\t};
    }
    \draw[sahmT] (n90) to[bend left=25] (n0);
    \draw[sahmT] (n0) to[bend left=25] (n-90);
    \draw[sahmT] (n-90) to[bend left=25] (n180);
    \draw[sahmT] (n180) to[bend left=25] (n90);
    \node[font=\small, text width=3cm, align=center] at (0,0) {then again\\with the next instruction};
\end{tikzpicture}
\end{center}

\begin{enumerate}
    \item \textbf{Fetch}: CO is sent to RAM, memory is read, the instruction arrives in RIM, then in RI.
          Then CO increases to point to the next instruction.
    \item \textbf{Decode}: the decoder reads the opcode in RI and finds the operands.
    \item \textbf{Execute}: if we need data from memory we read it, then the UAL computes the result.
    \item \textbf{Write back}: we store the result in a register or in memory, and we update the flags.
\end{enumerate}

\subsection{A complete example, step by step}
We take a simple machine with an accumulator ACC, and we want to compute $z = x + y$. The variable $x$
is at address 20, $y$ at 21, and $z$ at 22.

\begin{center}
\begin{tabular}{|c|l|l|}
    \hline
    \ras{Address} & \ras{Content} & \ras{Meaning} \\ \hline
    10 & \code{LOAD 20}  & $\text{ACC} \leftarrow [20]$ \\ \hline
    11 & \code{ADD 21}   & $\text{ACC} \leftarrow \text{ACC} + [21]$ \\ \hline
    12 & \code{STORE 22} & $[22] \leftarrow \text{ACC}$ \\ \hline
    20 & 7 & $x$ \\ \hline
    21 & 5 & $y$ \\ \hline
    22 & 0 & $z$ (before running) \\ \hline
\end{tabular}
\end{center}

We follow the second instruction \code{ADD 21} step by step (at the start, CO = 11 and ACC = 7):

\begin{center}
\small
\begin{tabularx}{\linewidth}{|>{\centering\arraybackslash}p{1.6cm}|R|c|c|c|c|c|}
    \hline
    \ras{Step} & \ras{What happens?} & \ras{CO} & \ras{RAM} & \ras{RIM} & \ras{RI} & \ras{ACC} \\ \hline
    fetch & \nql{RAM}{CO} & 11 & \cell{orangeDeep}{11} & -- & -- & 7 \\ \cline{2-7}
     & read memory: \nql{RIM}{[11]} & 11 & 11 & \cell{orangeDeep}{\code{ADD 21}} & -- & 7 \\ \cline{2-7}
     & \nql{RI}{RIM}, and \nql{CO}{CO + 1} & \cell{orangeDeep}{12} & 11 & \code{ADD 21} & \cell{orangeDeep}{\code{ADD 21}} & 7 \\ \hline
    decode & the decoder understands: ``add the content of address 21'' & 12 & 11 & \code{ADD 21} & \code{ADD 21} & 7 \\ \hline
    execute & \nql{RAM}{21}, then read: \nql{RIM}{[21]} & 12 & \cell{orangeDeep}{21} & \cell{orangeDeep}{5} & \code{ADD 21} & 7 \\ \cline{2-7}
     & UAL: $7 + 5$ & 12 & 21 & 5 & \code{ADD 21} & 7 \\ \hline
    write back & \nql{ACC}{12}, and update the flags & 12 & 21 & 5 & \code{ADD 21} & \cell{orangeDeep}{\textbf{12}} \\ \hline
\end{tabularx}
\end{center}

\begin{khulasa}
    After this instruction: ACC = 12 and CO = 12, so the next instruction is \code{STORE 22}, which writes 12
    in memory, and $z$ becomes 12.
\end{khulasa}

\begin{hal}
    Modern processors do not wait for one instruction to finish before starting the next. While one instruction
    runs, the next one is decoded and the third one is fetched. This is the \textbf{pipeline}, which we study in Chapter 4.
\end{hal}

%================================================
\section{Exercises}
%================================================

\begin{tamrin}
    Compute with the full adder (bit by bit) $1011 + 0110$ on 4 bits. What is the result, and the value of CF?
\end{tamrin}
\begin{hall}
    \begin{tabular}{|c|c|c|c|c|c|}
        \hline \ras{Bit} & \ras{$A$} & \ras{$B$} & \ras{$C_{in}$} & \ras{$S$} & \ras{$C_{out}$} \\ \hline
        0 & 1 & 0 & 0 & 1 & 0 \\ \hline 1 & 1 & 1 & 0 & 0 & 1 \\ \hline
        2 & 0 & 1 & 1 & 0 & 1 \\ \hline 3 & 1 & 0 & 1 & 0 & 1 \\ \hline
    \end{tabular}\\
    The result is $0001$ and CF = 1: $11 + 6 = 17 = 16 + 1$, and the ``16'' went into CF.
\end{hall}

\begin{tamrin}
    On 8 bits, give AL, CF and ZF after: (a) $255 + 1$ (b) $10 - 10$ (c) $10 - 20$.
\end{tamrin}
\begin{hall}
    (a) $256 - 256 = 0$: AL = 0, CF = 1, ZF = 1.\\
    (b) AL = 0, CF = 0, ZF = 1.\\
    (c) $10 - 20 = -10$, and on 8 bits $256 - 10 = 246 = \hex{F6}$: CF = 1 (borrow), ZF = 0.
\end{hall}

\begin{tamrin}
    Give the addressing mode of each instruction: (a) \code{mov al, [bx]} (b) \code{mov al, 10h} (c) \code{mov al, [10h]} (d) \code{mov al, bl}.
\end{tamrin}
\begin{hall}
    (a) register indirect \quad (b) immediate \quad (c) direct \quad (d) register.
\end{hall}

\begin{tamrin}
    In the example above ($z = x + y$), write the table of steps for the third instruction \code{STORE 22}.
\end{tamrin}
\begin{hall}
    \textbf{Fetch}: \nql{RAM}{12}, \nql{RIM}{STORE 22}, \nql{RI}{RIM}, CO = 13.
    \textbf{Decode}: ``write ACC at address 22''.
    \textbf{Execute and write}: \nql{RAM}{22}, \nql{RIM}{ACC} (12), then a write signal: $[22] = 12$.
\end{hall}

\chapter{Case Study: the Intel 8086 Processor}

\begin{khulasa}[: what will we learn in this chapter?]
\begin{itemize}
    \item The 80x86 family, and the internal structure of the 8086.
    \item The 8086 registers, and segmented addressing (how to compute a physical address).
    \item Writing a first program, and how an instruction runs inside the processor.
    \item The instructions: transfer, arithmetic, logic, jumps, loops, then how an instruction becomes machine code.
    \item Input/output, the stack, procedures, and interrupts.
\end{itemize}
\end{khulasa}

%================================================
\section{The 80x86 family}
%================================================

In 1978 Intel released the 8086. It was a big success because IBM chose its brother, the 8088, for the first
personal computer (IBM PC) in 1981. Since then, every new Intel and AMD processor can still run 8086 programs.
This family is called \textbf{x86}.

\begin{tarif}
    \term{Backward compatibility}{compatibilité ascendante}: a new processor understands all the instructions
    of the older processors of its family. A program written for the 8086 in 1980 can run on a modern Core i7.
    This is why the x86 family has lived for more than 45 years.
\end{tarif}

\begin{center}
\small
\begin{tabular}{|c|c|c|c|l|}
    \hline
    \ras{Processor} & \ras{Year} & \ras{Bits} & \ras{Max memory} & \ras{What is new} \\ \hline
    8086 / 8088 & 1978 & 16 & 1 MB & first processor of the family \\ \hline
    80286 & 1982 & 16 & 16 MB & protected mode \\ \hline
    80386 & 1985 & 32 & 4 GB & first 32-bit processor (IA-32) \\ \hline
    80486 & 1989 & 32 & 4 GB & cache memory and floating-point unit inside \\ \hline
    Pentium & 1993 & 32 & 4 GB & runs two instructions at once \\ \hline
    Athlon 64 (AMD) & 2003 & 64 & $2^{48}$ in practice & the AMD64 (x86-64) architecture \\ \hline
    Core i7 and later & 2008 & 64 & huge & several cores on the same chip \\ \hline
\end{tabular}
\end{center}

\begin{tarif}
    \textbf{General features of the 8086}:
    \begin{itemize}
        \item a \textbf{16-bit} processor: its registers and its data bus are 16 bits wide.
        \item a \textbf{20-bit} address bus: it can address 1 MB of memory.
        \item about 29\,000 transistors, and a frequency between 5 and 10 MHz.
        \item about \textbf{100 different instructions} (a CISC processor).
    \end{itemize}
\end{tarif}

\begin{hal}
    The only difference between the 8086 and the 8088 is the \textbf{external data bus}: 16 bits in the 8086,
    only 8 bits in the 8088. Inside, they are identical. IBM chose the 8088 because 8-bit support chips were
    cheaper at that time.
\end{hal}

%================================================
\section{Inside the 8086}
%================================================

The 8086 is split into two units that work \textbf{at the same time}:

\begin{center}
\begin{tikzpicture}[r/.style={reg, minimum width=1.25cm, font=\codefont\small}]
    \draw[brandTeal, very thick, rounded corners=6pt, fill=brandTeal!8] (0.3,-3.2) rectangle (6.6,2.6);
    \node[font=\bfseries, text=tealDeep] at (3.45,2.25) {Execution Unit (EU)};
    \foreach \n/\y in {AX/1.4, BX/0.8, CX/0.2, DX/-0.4} \node[r] at (1.3,\y) {\n};
    \foreach \n/\y in {SI/1.4, DI/0.8, BP/0.2, SP/-0.4} \node[r] at (2.7,\y) {\n};
    \draw[brandGreen!70!black, very thick, fill=brandGreen!20] (4.1,1.2) -- (5.9,1.2) -- (5.5,-0.2) -- (4.5,-0.2) -- cycle;
    \node[font=\bfseries\small] at (5,0.5) {UAL};
    \node[r, minimum width=2.5cm, font=\small] at (5,-1.0) {FLAGS};
    \node[unit=brandOrange, font=\small, minimum width=2.6cm] at (2.2,-2.1) {control unit};
    \draw[brandOrange, very thick, rounded corners=6pt, fill=brandOrange!8] (-6.3,-3.2) rectangle (-0.3,2.6);
    \node[font=\bfseries, text=orangeDeep] at (-3.3,2.25) {Bus Interface Unit (BIU)};
    \node[r] at (-4.6,1.4) {CS}; \node[r] at (-4.6,0.8) {DS};
    \node[r] at (-4.6,0.2) {SS}; \node[r] at (-4.6,-0.4) {ES};
    \node[r] at (-4.6,-1.0) {IP};
    \node[draw=brandDark, fill=brandOrange!25, thick, minimum width=1.4cm, minimum height=0.6cm, font=\small] at (-2,1.1) {adder $\Sigma$};
    \node[font=\scriptsize, align=center] at (-2,0.35) {computes the\\physical address};
    \foreach \i in {0,...,5} \node[draw=brandDark, fill=white, minimum width=0.55cm, minimum height=0.5cm] at (-4.2+0.6*\i,-2.1) {};
    \node[font=\scriptsize] at (-2.7,-2.75) {instruction queue (6 bytes)};
    \draw[sahmT, <->] (-0.3,-2.1) -- (0.3,-2.1);
    \draw[brandDark, line width=4pt, {Stealth[length=4mm]}-{Stealth[length=4mm]}] (-6.3,0.2) -- (-7.8,0.2);
    \node[font=\small, align=center] at (-7.1,-0.5) {system\\bus};
\end{tikzpicture}
\end{center}

\begin{itemize}
    \item The \term{Bus Interface Unit}{unité d'interface de bus} (BIU): it talks with memory. It computes the
          addresses, fetches instructions in advance into a \textbf{queue} of 6 bytes, and reads and writes data.
    \item The \term{Execution Unit}{unité d'exécution} (EU): it takes the instructions from the queue, decodes
          them, and runs them with the arithmetic and logic unit.
\end{itemize}

\begin{tashbih}
    In a restaurant, the \textbf{waiter} (BIU) brings orders from the customers and puts them in a queue in front of
    the cook, while the \textbf{cook} (EU) prepares the current order. The cook never waits for the waiter, because
    the next order is always in the queue. This is the first idea of a pipeline!
\end{tashbih}

%================================================
\section{The 8086 registers}
%================================================

The 8086 has \textbf{14 registers} of 16 bits.

\subsection{General registers}
Four data registers. Each one can be split into two halves of 8 bits:

\begin{center}
\begin{tikzpicture}[h/.style={draw=brandDark, thick, minimum width=2.2cm, minimum height=0.7cm, font=\codefont\bfseries}]
    \foreach \r/\n/\d [count=\i from 0] in {A/AX/{accumulator: arithmetic, multiply, divide, input/output},
                                         B/BX/{base: keeps an address in memory},
                                         C/CX/{counter: number of loop turns},
                                         D/DX/{data: with AX for multiply and divide}} {
        \node[font=\codefont\bfseries, text=tealDeep] at (-1.6,-0.9*\i) {\n};
        \node[h, fill=brandTeal!25] (hi\i) at (0,-0.9*\i) {\r H};
        \node[h, fill=brandTeal!10] (lo\i) at (2.2,-0.9*\i) {\r L};
        \node[font=\small, anchor=west] at (3.6,-0.9*\i) {\d};
    }
    \node[font=\scriptsize] at (0,0.6) {bits 15 to 8};
    \node[font=\scriptsize] at (2.2,0.6) {bits 7 to 0};
\end{tikzpicture}
\end{center}

\begin{mithal}
    If AX = \hex{1234} then AH = \hex{12} and AL = \hex{34}. If we change AL to \hex{FF}, AX becomes \hex{12FF}.
\end{mithal}

Why split the registers? Because we often work with small values (a character, a number from 0 to 255).
Instead of wasting a whole 16-bit register for one character, we can put two characters in one register:
one in AH and one in AL.

\begin{intabih}
    These registers are ``general'' because we can use them in any calculation, but some instructions force a
    register: \code{loop} always uses CX, \code{mul} and \code{div} always use AX and DX, and the
    \code{int 21h} services expect the function number in AH. You must remember these roles.
\end{intabih}

\subsection{Pointer and index registers}
\begin{center}
\begin{tabular}{|c|l|l|}
    \hline
    \ras{Register} & \ras{Name} & \ras{Job} \\ \hline
    SI & Source Index & index of the source (arrays and strings) \\ \hline
    DI & Destination Index & index of the destination \\ \hline
    SP & Stack Pointer & points to the top of the stack \\ \hline
    BP & Base Pointer & to reach data inside the stack \\ \hline
\end{tabular}
\end{center}

\begin{mithal}
    To walk through an array of 10 elements, we put its start address in SI, then increase SI at each turn
    (\code{inc si}). At the first turn SI points to the first element, at the second turn to the second one,
    and so on. SI plays the role of the \textbf{finger} you put on the element you are reading now.
\end{mithal}

\subsection{Segment registers}
\begin{center}
\begin{tabular}{|c|l|l|}
    \hline
    \ras{Register} & \ras{Name} & \ras{Points to the start of\dots} \\ \hline
    CS & Code Segment & the program segment (the instructions) \\ \hline
    DS & Data Segment & the data segment (the variables) \\ \hline
    SS & Stack Segment & the stack segment \\ \hline
    ES & Extra Segment & an extra segment (for copying strings, for example) \\ \hline
\end{tabular}
\end{center}

\subsection{Instruction pointer and flags}
\begin{itemize}
    \item \textbf{IP} (Instruction Pointer): the address of the next instruction inside the code segment. It is the
          CO we saw in Chapter 2. We cannot change it directly with \code{mov}, only with jumps and calls.
    \item \textbf{FLAGS}: 9 bits are used out of 16:
\end{itemize}

\begin{center}
\begin{tikzpicture}[b/.style={draw=brandDark, minimum width=0.75cm, minimum height=0.7cm, font=\codefont\scriptsize\bfseries}]
    \foreach \n [count=\i from 0] in {{}, {}, {}, {}, OF, DF, IF, TF, SF, ZF, {}, AF, {}, PF, {}, CF} {
        \pgfmathtruncatemacro\bit{15-\i}
        \node[b, fill=brandTeal!15] at (0.75*\i, 0) {\n};
        \node[font=\codefont\tiny] at (0.75*\i, 0.55) {\bit};
    }
\end{tikzpicture}
\end{center}

\begin{center}
\small
\begin{tabularx}{\linewidth}{|c|R||c|R|}
    \hline
    \ras{Flag} & \ras{Meaning} & \ras{Flag} & \ras{Meaning} \\ \hline
    CF & carry or borrow out of the last bit & OF & overflow: wrong result for signed numbers \\ \hline
    PF & the number of 1s in the low byte is even & DF & direction for strings (0: forward) \\ \hline
    AF & carry from bit 3 to bit 4 & IF & hardware interrupts allowed \\ \hline
    ZF & the result is zero & TF & step-by-step execution (for debugging) \\ \hline
    SF & the result is negative (last bit is 1) & & \\ \hline
\end{tabularx}
\end{center}

%================================================
\section{Segmented addressing}
%================================================

\subsection{The problem}
The 8086 registers are 16 bits wide, so they can write addresses from 0 to $2^{16} - 1 = 65535$ only
(that is 64 KB). But the address bus is 20 bits wide (1 MB). How can we write a 20-bit address with 16-bit registers?

\subsection{The solution: segment and offset}
We use \textbf{two registers}: a \textbf{segment} register gives the start of a zone of memory, and the
\textbf{offset} gives the distance inside this zone. We write the logical address as \texttt{segment:offset}.

\begin{tarif}
    \centering
    $\text{physical address} = \text{segment} \times 16 + \text{offset}$
    \par\small
    Multiplying by 16 = adding one zero on the right in hexadecimal.
\end{tarif}

\begin{khatwa}[: the physical address for CS = 1000H and IP = 0234H]
    \begin{enumerate}
        \item Multiply the segment by 16 (add a zero): \hex{1000} becomes \hex{10000}.
        \item Add the offset:
        \begin{center}
        {\codefont
        \begin{tabular}{r}
              10000 \\
            + \ 0234 \\ \hline
              10234
        \end{tabular}}
        \end{center}
        \item The physical address is \hex{10234}, a 20-bit address (5 hex digits).
    \end{enumerate}
\end{khatwa}

\begin{center}
\begin{tikzpicture}
    \draw[thick, fill=brandLight] (0,0) rectangle (2.5,-7);
    \node[font=\small\bfseries] at (1.25,0.35) {memory (1 MB)};
    \node[font=\codefont\scriptsize, anchor=east] at (-0.1,-0.1) {00000};
    \node[font=\codefont\scriptsize, anchor=east] at (-0.1,-6.9) {FFFFF};
    \fill[brandTeal!30] (0,-2) rectangle (2.5,-4.5);
    \draw[thick] (0,-2) rectangle (2.5,-4.5);
    \node[font=\scriptsize, align=center] at (1.25,-3.4) {segment\\(64 KB)};
    \draw[sahmO] (4,-2) node[right, font=\small] {segment start: 10000} -- (2.5,-2);
    \draw[sahmT] (4,-2.8) node[right, font=\small] {address 10234} -- (2.5,-2.8);
    \draw[decorate, decoration={brace, amplitude=5pt, mirror}, thick] (-0.1,-2) -- (-0.1,-2.8)
        node[midway, left=6pt, font=\small] {offset 0234};
\end{tikzpicture}
\end{center}

\subsection{Why multiply by 16?}
If we used the segment as it is (without multiplying), we could not go beyond $\hex{FFFF} + \hex{FFFF} = \hex{1FFFE}$,
about 128 KB only. With the multiplication by 16, a segment can point anywhere in 1 MB: $\hex{FFFF} \times 16 = \hex{FFFF0}$.
So each segment always starts at an address that can be divided by 16. This block of 16 bytes is called a
\term{paragraph}{paragraphe}.

\begin{khatwa}[: which zone does the segment DS = 2000H cover?]
    \begin{enumerate}
        \item The smallest offset is \hex{0000}: $\hex{20000} + \hex{0000} = \hex{20000}$.
        \item The biggest offset is \hex{FFFF}: $\hex{20000} + \hex{FFFF} = \hex{2FFFF}$.
        \item The segment covers the zone from \hex{20000} to \hex{2FFFF}, that is $2^{16} = 65536$ bytes = 64 KB.
    \end{enumerate}
\end{khatwa}

\begin{tashbih}
    Imagine a town with many streets. The full address of a house is ``\textbf{street} number 12, \textbf{house}
    number 34''. The segment register is the street number (it gives the start of a zone), and the offset is the
    house number inside the street. The postman needs both pieces of information to find the house.
\end{tashbih}

\begin{intabih}
    One physical address can be written in several ways: 1000:0234 and 1023:0004 give the same address
    \hex{10234}. Check it yourself: $\hex{10230} + \hex{0004} = \hex{10234}$.
\end{intabih}

\subsection{Default pairs}
Each offset register is used by default with a given segment register:

\begin{center}
\begin{tabular}{|c|c|l|}
    \hline
    \ras{Offset} & \ras{Default segment} & \ras{Use} \\ \hline
    IP & CS & fetching instructions \\ \hline
    SP, BP & SS & the stack \\ \hline
    BX, SI, direct address & DS & variables \\ \hline
    DI (in string instructions) & ES & destination of a copy \\ \hline
\end{tabular}
\end{center}

\begin{hal}
    In the .COM programs we write in this chapter, the four segments CS = DS = SS = ES are equal: the whole
    program fits in one 64 KB segment, and starts at offset \hex{100}. This is the meaning of the line
    \code{org 100h} that we find in all our programs.
\end{hal}

%================================================
\section{8086 addressing modes}
%================================================

\begin{tarif}
    The \term{effective address}{adresse effective} (EA): the offset the processor computes for a memory operand.
    Then it computes the physical address $= \text{DS} \times 16 + \text{EA}$.
\end{tarif}

\begin{center}
\small
\begin{tabularx}{\linewidth}{|>{\centering\arraybackslash}p{2.8cm}|>{\centering\arraybackslash}p{3.4cm}|R|}
    \hline
    \ras{Mode} & \ras{Example} & \ras{Effective address EA} \\ \hline
    immediate & \code{mov ax, 5} & none: the value is in the instruction \\ \hline
    register & \code{mov ax, bx} & none: the value is in a register \\ \hline
    direct & \code{mov ax, [0200h]} & EA = \hex{0200} \\ \hline
    register indirect & \code{mov ax, [bx]} & EA = BX (or SI, DI, BP) \\ \hline
    based & \code{mov ax, [bx+4]} & EA = BX + 4 \\ \hline
    indexed & \code{mov ax, [si+4]} & EA = SI + 4 \\ \hline
    based indexed & \code{mov ax, [bx+si+4]} & EA = BX + SI + 4 \\ \hline
\end{tabularx}
\end{center}

\begin{khatwa}[: where does \code{mov ax, [bx+si+4]} read?]
    We suppose: DS = \hex{2000}, BX = \hex{0100}, SI = \hex{0020}.
    \begin{enumerate}
        \item Effective address: $\text{EA} = \hex{0100} + \hex{0020} + 4 = \hex{0124}$.
        \item Physical address: $\hex{2000} \times 16 + \hex{0124} = \hex{20000} + \hex{0124} = \hex{20124}$.
        \item The processor reads a word (two bytes): the byte at \hex{20124} goes into AL, and the byte at
              \hex{20125} goes into AH.
    \end{enumerate}
\end{khatwa}

\begin{khatwa}[: \code{mov ax, [bp+2]} uses the segment SS]
    We suppose SS = \hex{3000} and BP = \hex{FFF0}.
    \begin{enumerate}
        \item Effective address: $\text{EA} = \hex{FFF0} + 2 = \hex{FFF2}$.
        \item Because we used BP, the default segment is SS, not DS.
        \item Physical address: $\hex{30000} + \hex{FFF2} = \hex{3FFF2}$.
    \end{enumerate}
    This is important in procedures: BP is used to read the parameters that are in the stack.
\end{khatwa}

\begin{intabih}
    The 8086 stores a word in memory with the \textbf{small byte first} (little endian): the word \hex{1234}
    is stored as 34 then 12. Also, two memory operands in the same instruction are not allowed: \code{mov [a], [b]} is wrong.
\end{intabih}

%================================================
\section{A first program}
%================================================

\subsection{The tools}
We write the program in \textbf{EMU8086}: \emph{File $\rightarrow$ New $\rightarrow$ COM template}, then we press
\emph{Emulate} and run one instruction at a time with the \emph{Single step} button.

Every assembly program goes through three steps:
\begin{center}
\begin{tikzpicture}[node distance=1.9cm]
    \node[unit=brandTeal, minimum width=2.8cm] (src) {source file\\\code{hello.asm}};
    \node[unit=brandOrange, minimum width=2.8cm, right=of src] (bin) {executable file\\\code{hello.com}};
    \node[unit=brandGreen, minimum width=2.8cm, right=of bin] (run) {running\\in memory};
    \draw[sahmT] (src) -- node[above, font=\scriptsize, text=brandDark] {assembler} (bin);
    \draw[sahmT] (bin) -- node[above, font=\scriptsize, text=brandDark] {loading} (run);
\end{tikzpicture}
\end{center}
\begin{enumerate}
    \item \textbf{Writing}: we type the instructions in a text file (the source).
    \item \textbf{Assembling}: the assembler turns each instruction into machine-code bytes, and computes the
          addresses of the variables and labels.
    \item \textbf{Loading and running}: the system copies the bytes into memory starting at offset \hex{100},
          sets IP = \hex{100}, and the processor starts.
\end{enumerate}

\subsection{The program}
\asm{ch3/hello.asm}
\screen{ch3/hello.out}

\begin{sharh}
    \s{1}{A comment: everything after \code{;} is not run by the processor; it is an explanation for the reader.}
    \s{2}{\code{org 100h}: a .COM program starts at offset \hex{100}.}
    \s{4}{Put the address of the message \code{msg} into DX.}
    \s{5}{Choose DOS function \hex{09}: ``print the string whose address is in DX''.}
    \s{6}{\code{int 21h}: ask DOS to run the function chosen in AH.}
    \s{7}{\code{ret}: end of the program, back to the system.}
    \s{9}{Define the variable \code{msg} with \code{db} (define byte). The sign \code{\$} marks the end of the text.}
\end{sharh}

\begin{intabih}
    We put the data \textbf{after} \code{ret}. If we put it before the instructions, the processor would try
    to run the letters of the message as if they were instructions!
\end{intabih}

\begin{khatwa}[: the program in memory after loading]
    These are the real bytes made by the assembler:
    \begin{center}
    \small
    \begin{tabular}{|c|l|l|}
        \hline \ras{Offset} & \ras{Bytes} & \ras{Content} \\ \hline
        0100 & {\codefont 8D 16 09 01} & \code{lea dx, msg} (address \hex{0109}) \\ \hline
        0104 & {\codefont B4 09} & \code{mov ah, 09h} \\ \hline
        0106 & {\codefont CD 21} & \code{int 21h} \\ \hline
        0108 & {\codefont C3} & \code{ret} \\ \hline
        0109 & {\codefont 53 61 6C 61 6D \dots} & the letters \code{S a l a m \dots} \\ \hline
    \end{tabular}
    \end{center}
    The whole program is 9 bytes of instructions, and the message starts right after, at \hex{0109}. Notice
    that the assembler computed the address of \code{msg} by itself and put it inside the \code{lea} instruction.
\end{khatwa}

\subsection{Defining variables}
\begin{center}
\begin{tabular}{|l|l|}
    \hline
    \ras{Writing} & \ras{Meaning} \\ \hline
    \code{x db 5} & one byte with value 5 \\ \hline
    \code{y dw 1234h} & one word (two bytes) with value \hex{1234} \\ \hline
    \code{t db 14, 3, 27} & an array of 3 bytes \\ \hline
    \code{buf db 10 dup(0)} & 10 bytes, all 0 \\ \hline
    \code{s db 'Hi\$'} & a string of characters \\ \hline
\end{tabular}
\end{center}

%================================================
\section{How an instruction runs inside the 8086}
%================================================

Let us follow the instruction \code{mov cx, [si]} from the program \code{transfer.asm} (next section). In
memory, this instruction is written with the bytes 8B 0C at offset \hex{010C}, and SI = \hex{0112}.

\begin{center}
\small
\begin{tabularx}{\linewidth}{|>{\centering\arraybackslash}p{1.3cm}|>{\centering\arraybackslash}p{1.2cm}|R|}
    \hline
    \ras{Step} & \ras{Unit} & \ras{What happens} \\ \hline
    1 & BIU & computes the physical address $\text{CS} \times 16 + \text{IP}$, and fetches the bytes 8B 0C into the queue (often before we need them). \\ \hline
    2 & EU & takes 8B from the queue: ``move a word from memory to a register''. Then 0C: ``register CX, address in SI''. IP becomes \hex{010E}. \\ \hline
    3 & EU & computes the effective address: EA = SI = \hex{0112}. \\ \hline
    4 & BIU & computes $\text{DS} \times 16 + \hex{0112}$, sends it on the address bus, with a read signal. \\ \hline
    5 & BIU & the word \hex{ABCD} arrives on the data bus (CD then AB). \\ \hline
    6 & EU & writes the result into CX: CX = \hex{ABCD}. The flags do not change with \code{mov}. \\ \hline
\end{tabularx}
\end{center}

%================================================
\section{The instruction set in detail}
%================================================

We use these short names in the tables: reg = register, mem = memory box, imm = immediate value.

\subsection{Transfer instructions}
\begin{center}
\small
\begin{tabular}{|l|l|l|}
    \hline
    \ras{Instruction} & \ras{Meaning} & \ras{Example} \\ \hline
    \code{mov dest, src} & $\text{dest} \leftarrow \text{src}$ & \code{mov ax, bx} \\ \hline
    \code{xchg a, b} & swap the two values & \code{xchg ax, bx} \\ \hline
    \code{lea reg, mem} & reg gets the address of the variable & \code{lea si, tab} \\ \hline
    \code{push src} & put on the stack & \code{push ax} \\ \hline
    \code{pop dest} & take from the stack & \code{pop ax} \\ \hline
\end{tabular}
\end{center}

\asm{ch3/transfer.asm}

\begin{center}
\small
\begin{tabular}{|c|l|c|c|c|c|}
    \hline
    \ras{Line} & \ras{Instruction} & \ras{AX} & \ras{BX} & \ras{CX} & \ras{SI} \\ \hline
    4 & \code{mov ax, 1234h} & \cell{orangeDeep}{1234} & 0000 & 0000 & 0000 \\ \hline
    5 & \code{mov bl, ah} & 1234 & \cell{orangeDeep}{0012} & 0000 & 0000 \\ \hline
    6 & \code{mov bh, al} & 1234 & \cell{orangeDeep}{3412} & 0000 & 0000 \\ \hline
    7 & \code{xchg ax, bx} & \cell{orangeDeep}{3412} & \cell{orangeDeep}{1234} & 0000 & 0000 \\ \hline
    8 & \code{lea si, num} & 3412 & 1234 & 0000 & \cell{orangeDeep}{0112} \\ \hline
    9 & \code{mov cx, [si]} & 3412 & 1234 & \cell{orangeDeep}{ABCD} & 0112 \\ \hline
    10 & \code{mov [res], ax} & 3412 & 1234 & ABCD & 0112 \\ \hline
\end{tabular}
\end{center}
\begin{sharh}
    \s{5--6}{Copy the two halves of AX into BX, swapped: BH = 34 and BL = 12.}
    \s{8}{\code{lea} gives the \textbf{address} \hex{0112}, not the value.}
    \s{9}{Read the word stored at address \hex{0112}: \hex{ABCD}.}
    \s{10}{The only line that writes into memory: \code{res} becomes \hex{3412}.}
\end{sharh}

\subsection{Arithmetic instructions}
\begin{center}
\small
\begin{tabularx}{\linewidth}{|l|R|l|}
    \hline
    \ras{Instruction} & \ras{Meaning} & \ras{Flags} \\ \hline
    \code{add d, s} & $d \leftarrow d + s$ & all \\ \hline
    \code{adc d, s} & $d \leftarrow d + s + \text{CF}$ (to add big numbers) & all \\ \hline
    \code{sub d, s} & $d \leftarrow d - s$ & all \\ \hline
    \code{inc d} & $d \leftarrow d + 1$ & all except CF \\ \hline
    \code{dec d} & $d \leftarrow d - 1$ & all except CF \\ \hline
    \code{neg d} & $d \leftarrow -d$ (two's complement) & all \\ \hline
    \code{cmp a, b} & computes $a - b$ \textbf{without keeping the result} & all \\ \hline
    \code{mul s} & unsigned multiply (\code{imul} for signed) & CF, OF \\ \hline
    \code{div s} & unsigned divide (\code{idiv} for signed) & undefined \\ \hline
\end{tabularx}
\end{center}

\asm{ch3/arith.asm}

\begin{center}
\small
\begin{tabular}{|c|l|c|c|c|}
    \hline
    \ras{Line} & \ras{Instruction} & \ras{Result} & \ras{Decimal} & \ras{CF} \\ \hline
    4 & \code{mov al, 25} & AL = 19 & 25 & 0 \\ \hline
    5 & \code{add al, 17} & AL = 2A & 42 & 0 \\ \hline
    6 & \code{sub al, 2} & AL = 28 & 40 & 0 \\ \hline
    7 & \code{inc al} & AL = 29 & 41 & 0 \\ \hline
    8 & \code{dec al} & AL = 28 & 40 & 0 \\ \hline
    9 & \code{neg al} & AL = D8 & $-40$ & 1 \\ \hline
    10 & \code{mov bx, 1000} & BX = 03E8 & 1000 & 1 \\ \hline
    11 & \code{add bx, 2500} & BX = 0DAC & 3500 & 0 \\ \hline
\end{tabular}
\end{center}

\begin{khatwa}[: why is $-40$ written \hex{D8}? (two's complement)]
    \begin{enumerate}
        \item Write 40 in binary: \bin{0010\,1000}.
        \item Flip all the bits: \bin{1101\,0111}.
        \item Add 1: $\bin{1101\,1000} = \hex{D8}$.
    \end{enumerate}
    Check: $40 + (-40) = \hex{28} + \hex{D8} = \hex{100}$, and on 8 bits only 00 is left. Correct!
\end{khatwa}

\subsubsection{Multiplication and division}
The processor chooses the place depending on the size of the operand:
\begin{center}
\small
\begin{tabular}{|c|c|c|c|}
    \hline
    \ras{Instruction} & \ras{Operand} & \ras{Operation} & \ras{Result} \\ \hline
    \code{mul bl} & 8 bits & $\text{AL} \times \text{BL}$ & AX \\ \hline
    \code{mul cx} & 16 bits & $\text{AX} \times \text{CX}$ & DX:AX (32 bits) \\ \hline
    \code{div bl} & 8 bits & $\text{AX} \div \text{BL}$ & AL = quotient, AH = remainder \\ \hline
    \code{div cx} & 16 bits & $\text{DX:AX} \div \text{CX}$ & AX = quotient, DX = remainder \\ \hline
\end{tabular}
\end{center}

\asm{ch3/muldiv.asm}

\begin{sharh}
    \s{6}{$12 \times 11 = 132 = \hex{84}$: AX = \hex{0084}.}
    \s{10}{$300 \times 500 = 150000$, which is bigger than 65535, so it is split over two registers:
           DX = \hex{0002} and AX = \hex{49F0} (because $\hex{249F0} = 150000$).}
    \s{14}{$47 = 5 \times 9 + 2$: AL = 09 (quotient) and AH = 02 (remainder), so AX = \hex{0209}.}
\end{sharh}

\begin{intabih}
    Before \code{div cx} we must put DX = 0 if the number to divide is only in AX; otherwise the processor takes DX
    as part of the number. Dividing by 0 stops the program (interrupt number 0).
\end{intabih}

\subsubsection{Adding numbers bigger than 16 bits: \code{adc}}
The 8086 adds 16 bits at a time. To add two 32-bit numbers, we split each number into two words (high and
low). We add the low words with \code{add}, then the high words with \code{adc}, which adds the carry CF
from the first addition, exactly like we carry ``one'' to the next column when we add with a pen.

\asm{ch3/add32.asm}

\begin{center}
\small
\begin{tabular}{|c|l|c|c|c|}
    \hline
    \ras{Line} & \ras{Instruction} & \ras{DX} & \ras{AX} & \ras{CF} \\ \hline
    5--6 & load the first number & 0001 & FFFF & 0 \\ \hline
    7 & \code{add ax, 0001h} & 0001 & \cell{orangeDeep}{0000} & \cell{orangeDeep}{\textbf{1}} \\ \hline
    8 & \code{adc dx, 0000h} & \cell{orangeDeep}{0002} & 0000 & 0 \\ \hline
\end{tabular}
\end{center}
The result is DX:AX = 0002:0000 = \hex{00020000}. If we used \code{add} on line 8, the carry would be lost and
we would get \hex{00010000}, a wrong result.

\subsection{Logic instructions}
\begin{center}
\small
\begin{tabularx}{\linewidth}{|l|R|l|}
    \hline
    \ras{Instruction} & \ras{Usual use} & \ras{Example} \\ \hline
    \code{and} & switch bits off (a mask) & \code{and al, 0Fh} \\ \hline
    \code{or} & switch bits on & \code{or al, 80h} \\ \hline
    \code{xor} & flip bits, or clear a register & \code{xor cx, cx} \\ \hline
    \code{not} & flip all the bits & \code{not al} \\ \hline
    \code{test} & \code{and} without keeping the result (check a bit) & \code{test al, 1} \\ \hline
    \code{shl}, \code{shr} & shift: multiply or divide by 2 & \code{shl bl, 1} \\ \hline
    \code{rol}, \code{ror} & rotate: the bit that leaves comes back on the other side & \code{rol al, 1} \\ \hline
\end{tabularx}
\end{center}

\asm{ch3/logic.asm}

\begin{center}
\small
\begin{tabular}{|c|l|c|c|}
    \hline
    \ras{Line} & \ras{Instruction} & \ras{AL in binary} & \ras{Hex} \\ \hline
    4 & \code{mov al, 10110110b} & {\codefont 1011 0110} & B6 \\ \hline
    5 & \code{and al, 00001111b} & {\codefont 0000 0110} & 06 \\ \hline
    6 & \code{or  al, 10000000b} & {\codefont 1000 0110} & 86 \\ \hline
    7 & \code{xor al, 11111111b} & {\codefont 0111 1001} & 79 \\ \hline
    8 & \code{not al} & {\codefont 1000 0110} & 86 \\ \hline
\end{tabular}
\end{center}

\begin{sharh}
    \s{5}{\code{and} with 1 keeps the bit, and with 0 switches it off: we keep only the 4 low bits.}
    \s{6}{\code{or} with 1 switches the bit on: bit 7 is now 1.}
    \s{10--12}{BL: $5 \rightarrow 10 \rightarrow 20 \rightarrow 10$. A shift left multiplies by 2, and a shift right divides by 2.}
    \s{13}{\code{xor cx, cx}: every bit with itself gives 0, so CX = 0. Faster and shorter than \code{mov cx, 0}.}
\end{sharh}

\subsubsection{Practical uses of masks}
\begin{khatwa}[: is the number in AL even?]
    A number is even if bit 0 is 0. We write \code{test al, 1}: this instruction computes $\text{AL} \cdot 1$
    without changing AL.
    \begin{itemize}
        \item If AL = 6 = \bin{0000\,0110}: bit 0 is 0, the result is 0, and ZF = 1: the number is even.
        \item If AL = 7 = \bin{0000\,0111}: bit 0 is 1, the result is 1, and ZF = 0: the number is odd.
    \end{itemize}
    Then we write \code{jz even} to jump when the number is even.
\end{khatwa}

\begin{khatwa}[: turn a small letter into a capital letter with a mask]
    In the ASCII table: \code{'a'} $= \hex{61} = \bin{0110\,0001}$ and \code{'A'} $= \hex{41} = \bin{0100\,0001}$.
    The only difference is bit 5. So \code{and al, 11011111b} (that is \code{0DFh}) switches bit 5 off and turns
    any small letter into a capital letter, and leaves a capital letter as it is. This is another way than
    \code{sub al, 32}, which we see in the input section.
\end{khatwa}

\subsection{Flow control: jumps}
Usually the processor runs the instructions in order. A \textbf{jump} changes IP, so execution goes somewhere else.

\begin{center}
\small
\begin{tabularx}{\linewidth}{|l|R|l|}
    \hline
    \ras{Instruction} & \ras{Jumps if\dots} & \ras{Flags} \\ \hline
    \code{jmp L} & always (unconditional jump) & -- \\ \hline
    \code{je}, \code{jz} & equal / result is zero & ZF = 1 \\ \hline
    \code{jne}, \code{jnz} & not equal & ZF = 0 \\ \hline
    \code{jg}, \code{jge} & greater / greater or equal (\textbf{signed}) & ZF, SF, OF \\ \hline
    \code{jl}, \code{jle} & less / less or equal (\textbf{signed}) & SF, OF \\ \hline
    \code{ja}, \code{jae} & above / above or equal (\textbf{unsigned}) & CF, ZF \\ \hline
    \code{jb}, \code{jbe} & below / below or equal (\textbf{unsigned}) & CF, ZF \\ \hline
\end{tabularx}
\end{center}

\begin{tashbih}
    \code{cmp} is like \textbf{putting two things on a scale}, and a conditional jump is like saying: ``\textbf{if}
    the right side goes down, go right''. The scale (\code{cmp}) sets the flags, and the jump looks at them.
\end{tashbih}

\subsubsection{Example: if \dots\ else}
\asm{ch3/ifelse.asm}
\screen{ch3/ifelse.out}

\begin{center}
\begin{tikzpicture}[node distance=0.6cm]
    \node[unit] (c) {\code{cmp al, [b]}};
    \node[diamond, aspect=2.2, draw=brandOrange, fill=brandOrange!15, thick, below=of c, font=\small] (q) {$a > b$?};
    \node[unit=brandGreen, below right=0.4cm and 0.6cm of q] (y) {\code{lea dx, msg1}};
    \node[unit=brandPurple, below left=0.4cm and 0.6cm of q] (n) {\code{lea dx, msg2}};
    \node[unit, below=1.8cm of q] (p) {\code{int 21h} (print)};
    \draw[sahm] (c) -- (q);
    \draw[sahm] (q) -| node[pos=0.2, above, font=\small] {yes: \code{jg first}} (y);
    \draw[sahm] (q) -| node[pos=0.2, above, font=\small] {no} (n);
    \draw[sahm] (y) |- (p);
    \draw[sahm] (n) |- node[pos=0.7, below, font=\small] {\code{jmp show}} (p);
\end{tikzpicture}
\end{center}

\begin{sharh}
    \s{5}{\code{cmp} computes $7 - 12 = -5$ and does not keep it, but it sets the flags (SF = 1).}
    \s{6}{\code{jg}: is $a > b$? No, so we do not jump and continue at line 7.}
    \s{7--8}{Choose the second message, then \code{jmp show} is \textbf{needed} so that line 10 is not run.}
    \s{12--13}{The printing is shared by the two cases.}
\end{sharh}

\begin{intabih}
    We often forget \code{jmp show} on line 8! Without it, the program also runs \code{lea dx, msg1}, so it
    always prints the first message.
\end{intabih}

\subsubsection{Signed or unsigned? \code{jg} or \code{ja}}
The same bits can be read in two ways (we saw this in Chapter 2). AL = \hex{C8} is 200 unsigned, and $-56$
signed. So the processor has two kinds of jumps: \code{ja}/\code{jb} for unsigned numbers, and \code{jg}/\code{jl}
for signed numbers.

\asm{ch3/signed.asm}
\screen{ch3/signed.out}

\begin{sharh}
    \s{4--6}{\code{ja}: compare 200 with 100 as positive numbers. $200 > 100$, so we jump.}
    \s{14--16}{\code{jg}: the same bits \hex{C8}, but read as signed: $-56 > 100$? No, so we do not jump.}
\end{sharh}

\begin{intabih}
    Choose the jump from the \textbf{kind of data}: ages, addresses, character codes, sizes are unsigned
    (\code{ja}, \code{jb}). Temperatures, balances, differences are signed (\code{jg}, \code{jl}).
\end{intabih}

\subsubsection{From the algorithm to assembly}
Each control structure from the algorithms course has a fixed ``translation'' in assembly. The idea is
always the same: we \textbf{compare} with \code{cmp}, then we \textbf{jump} depending on the result.

\newcommand{\stk}[1]{\begin{tabular}[t]{@{}l@{}}#1\end{tabular}}
\begin{center}
\small
\begin{tabular}{|l|l|l|}
    \hline
    \ras{Structure} & \ras{In C} & \ras{In assembly} \\ \hline
    if & \code{if (a == b) x = 1;} &
        \stk{\code{    cmp a, b}\\ \code{    jne end}\\ \code{    mov x, 1}\\ \code{end:}} \\ \hline
    while & \code{while (cx > 0) \{ ... \}} &
        \stk{\code{L:  cmp cx, 0}\\ \code{    jle end}\\ \code{    ...}\\ \code{    jmp L}\\ \code{end:}} \\ \hline
    for & \code{for (cx = 10; cx > 0; cx--)} &
        \stk{\code{    mov cx, 10}\\ \code{L:  ...}\\ \code{    loop L}} \\ \hline
\end{tabular}
\end{center}

\begin{intabih}
    In the translation of \code{if}, we jump with the \textbf{opposite} condition (\code{jne} instead of \code{je}).
    We tell the processor: ``if they are not equal, skip the assignment''. This makes the program shorter.
\end{intabih}

\subsection{Loops and repetition}
\begin{tarif}
    \code{loop L} is like two instructions: \code{dec cx}, then ``if CX $\neq 0$, jump to L''.
    So CX is the loop counter.
\end{tarif}

\subsubsection{The sum of the numbers from 1 to 10}
\asm{ch3/loopsum.asm}

\begin{center}
\small
\begin{tabular}{|c|c|c|c|c|c|c|c|c|c|c|}
    \hline
    \ras{Turn} & 1 & 2 & 3 & 4 & 5 & 6 & 7 & 8 & 9 & 10 \\ \hline
    \ras{CX before \code{add}} & 10 & 9 & 8 & 7 & 6 & 5 & 4 & 3 & 2 & 1 \\ \hline
    \ras{AX after \code{add}} & 10 & 19 & 27 & 34 & 40 & 45 & 49 & 52 & 54 & \cell{orangeDeep}{\textbf{55}} \\ \hline
\end{tabular}
\end{center}
After the tenth turn CX = 0, so the loop stops, and AX = \hex{0037} = 55. Notice that we add backwards:
$10 + 9 + \dots + 1$, and the result is the same.

\begin{intabih}
    If we enter the loop with CX = 0, \code{loop} first decreases CX, which becomes \hex{FFFF}, and then repeats
    65535 more times! To avoid this, we put \code{jcxz end} before the loop: ``if CX = 0, jump to the end''.
\end{intabih}

\subsubsection{Finding the biggest value in an array}
\asm{ch3/maxarray.asm}

\begin{center}
\small
\begin{tabular}{|c|c|c|c|c|}
    \hline
    \ras{SI points to} & \ras{Element} & \ras{Element $>$ AL?} & \ras{AL after} & \ras{CX after \code{loop}} \\ \hline
    tab[0] & 14 & no (equal) & 14 & 5 \\ \hline
    tab[1] & 3 & no & 14 & 4 \\ \hline
    tab[2] & 27 & yes & 27 & 3 \\ \hline
    tab[3] & 9 & no & 27 & 2 \\ \hline
    tab[4] & 31 & yes & \cell{orangeDeep}{\textbf{31}} & 1 \\ \hline
    tab[5] & 8 & no & 31 & 0 \\ \hline
\end{tabular}
\end{center}
The result: AL = \hex{1F} = 31, stored in \code{max}. We used \code{jbe} because the values are \textbf{unsigned}.

\subsubsection{Repetition with the \code{rep} prefix}
String instructions work on DS:SI (the source) and ES:DI (the destination), and the prefix \code{rep}
repeats them CX times.

\asm{ch3/repmovs.asm}
\screen{ch3/repmovs.out}

\begin{sharh}
    \s{7}{\code{cld}: set DF = 0, so SI and DI increase after each copy.}
    \s{8}{\code{rep movsb} copies one byte from [SI] to [DI], and repeats 6 times: the five letters and \code{\$}.}
\end{sharh}

\subsection{Generating machine code}
The assembler turns each instruction into bytes. These are the real bytes of the program \code{machine.asm}:

\asm{ch3/machine.asm}

\begin{center}
\small
\begin{tabular}{|c|l|l|}
    \hline
    \ras{Address} & \ras{Machine code (hex)} & \ras{Instruction} \\ \hline
    0100 & {\codefont B8 34 12} & \code{mov ax, 1234h} \\ \hline
    0103 & {\codefont 89 C3} & \code{mov bx, ax} \\ \hline
    0105 & {\codefont 04 05} & \code{add al, 5} \\ \hline
    0107 & {\codefont 88 07} & \code{mov [bx], al} \\ \hline
    0109 & {\codefont 41} & \code{inc cx} \\ \hline
    010A & {\codefont EB 01} & \code{jmp fin} \\ \hline
    010C & {\codefont 90} & \code{nop} \\ \hline
    010D & {\codefont C3} & \code{ret} \\ \hline
\end{tabular}
\end{center}

\begin{khatwa}[: decoding B8 34 12]
    \begin{enumerate}
        \item B8 = \bin{1011\,1000}: the first five bits 10111 mean ``\code{mov} an immediate value into a 16-bit register'',
              and the last three bits 000 are the number of the register AX.
        \item 34 12: the value \hex{1234}, written with the small byte first (little endian).
    \end{enumerate}
\end{khatwa}

\begin{khatwa}[: decoding 89 C3 (\code{mov bx, ax})]
    The first byte 89 = \bin{100010\,0\,1}:
    \begin{itemize}
        \item 100010: the opcode of \code{mov} between a register and a register or memory.
        \item d = 0: the register named in the \emph{reg} field is the \textbf{source}.
        \item w = 1: we work on 16 bits (if it were 0 we would work on 8 bits).
    \end{itemize}
    The second byte C3 = \bin{11\,000\,011} is called ModR/M:
    \begin{center}
    \begin{tikzpicture}[f/.style={draw=brandDark, thick, minimum height=0.7cm, font=\codefont\bfseries}]
        \node[f, minimum width=1.6cm, fill=brandOrange!20] (m) at (0,0) {11};
        \node[f, minimum width=2.2cm, fill=brandTeal!20, anchor=west] (r) at (m.east) {000};
        \node[f, minimum width=2.2cm, fill=brandGreen!20, anchor=west] (rm) at (r.east) {011};
        \node[below=0.1cm of m, font=\small] {mod: register};
        \node[below=0.1cm of r, font=\small] {reg = AX};
        \node[below=0.1cm of rm, font=\small] {r/m = BX};
    \end{tikzpicture}
    \end{center}
    So: move AX (the source) into BX.
\end{khatwa}

\begin{center}
\small
\begin{tabular}{|c|c|c|c|c|c|c|c|c|}
    \hline
    \ras{Number} & 000 & 001 & 010 & 011 & 100 & 101 & 110 & 111 \\ \hline
    \ras{w = 1} & AX & CX & DX & BX & SP & BP & SI & DI \\ \hline
    \ras{w = 0} & AL & CL & DL & BL & AH & CH & DH & BH \\ \hline
\end{tabular}
\end{center}

\begin{mithal}[: 88 07 and EB 01]
    88 = \bin{100010\,0\,0}: \code{mov} on 8 bits. 07 = \bin{00\,000\,111}: mod = 00 (memory without offset),
    reg = 000 (AL), r/m = 111 ([BX]). So \code{mov [bx], al}.\\[3pt]
    EB 01: a short jump, and the offset is \textbf{relative}: after reading \code{jmp}, IP = \hex{010C}, and we add 1
    to reach \hex{010D} (\code{fin}). This is how we jumped over \code{nop}.
\end{mithal}

%================================================
\section{Input and output}
%================================================

\subsection{Ports: \code{in} and \code{out}}
The processor talks with devices through numbered \term{ports}{ports d'entrée/sortie}, from 0 to \hex{FFFF}.
\begin{itemize}
    \item \code{in al, 60h}: read one byte from port \hex{60} (the keyboard in a PC).
    \item \code{out 61h, al}: send AL to port \hex{61} (the speaker).
\end{itemize}
Normal programs rarely use ports directly; they ask the operating system through interrupts.

\subsection{DOS services: the \code{int 21h} interrupt}
\begin{center}
\small
\begin{tabularx}{\linewidth}{|c|R|R|}
    \hline
    \ras{AH} & \ras{Function} & \ras{Inputs and outputs} \\ \hline
    \hex{01} & read a key from the keyboard and show it & output: the character in AL \\ \hline
    \hex{02} & print a character & input: the character in DL \\ \hline
    \hex{09} & print a string ending with \code{\$} & input: its address in DX \\ \hline
    \hex{4C} & end the program & AL = exit code \\ \hline
\end{tabularx}
\end{center}

\asm{ch3/io.asm}
\screen{ch3/io.out}

\begin{sharh}
    \s{7--8}{Wait for the user to press a key (here \code{m}); its ASCII code arrives in AL: 109.}
    \s{9}{In the ASCII table, a capital letter is 32 less than the small one: $109 - 32 = 77$, which is the code of \code{'M'}.}
    \s{10}{Keep the letter in BL, because the printing on line 13 changes AL.}
    \s{14--16}{Function \hex{02} prints the character that is in DL.}
    \s{20}{\code{13, 10}: back to the start of the line, then a new line.}
\end{sharh}

%================================================
\section{The stack}
%================================================

\begin{tarif}
    The \term{stack}{pile}: a zone of memory where we put values and take them back with the rule
    \textbf{last in, first out} (LIFO). Its top is given by SS:SP.
\end{tarif}

\begin{tashbih}
    The stack is like a \textbf{pile of plates}: we put the new plate on top of the pile, and we always take the plate on top.
\end{tashbih}

\begin{itemize}
    \item \code{push ax}: \nql{SP}{SP - 2}, then write AX at [SP].
    \item \code{pop ax}: read [SP] into AX, then \nql{SP}{SP + 2}.
\end{itemize}

\begin{intabih}
    In the 8086 the stack \textbf{grows towards small addresses}: each \code{push} decreases SP. And we always
    put and take \textbf{words} (16 bits).
\end{intabih}

\asm{ch3/stack.asm}

\begin{center}
\begin{tikzpicture}[c/.style={draw=brandDark, fill=white, minimum width=1.6cm, minimum height=0.55cm, font=\codefont\small},
    a/.style={font=\codefont\scriptsize, text=tealDeep}]
    \foreach \t/\x/\v/\w/\sp in {{start}/0/{}/{}/FFFE, {after \code{push ax}}/3.3/1111/{}/FFFC,
                                {after \code{push bx}}/6.6/1111/2222/FFFA, {after \code{pop ax}}/9.9/1111/{}/FFFC} {
        \node[font=\small\bfseries] at (\x,1.1) {\t};
        \node[c] (u) at (\x,0.3) {\v};
        \node[c] (d) at (\x,-0.25) {\w};
        \node[a, left=0.05cm of u] {FFFC};
        \node[a, left=0.05cm of d] {FFFA};
        \node[font=\codefont\small] at (\x,-1) {SP = \sp};
    }
\end{tikzpicture}
\end{center}

\begin{center}
\small
\begin{tabular}{|c|l|c|c|c|}
    \hline
    \ras{Line} & \ras{Instruction} & \ras{AX} & \ras{BX} & \ras{SP} \\ \hline
    4--5 & \code{mov ax}, \code{mov bx} & 1111 & 2222 & FFFE \\ \hline
    6 & \code{push ax} & 1111 & 2222 & \cell{orangeDeep}{FFFC} \\ \hline
    7 & \code{push bx} & 1111 & 2222 & \cell{orangeDeep}{FFFA} \\ \hline
    8 & \code{pop ax} & \cell{orangeDeep}{2222} & 2222 & FFFC \\ \hline
    9 & \code{pop bx} & 2222 & \cell{orangeDeep}{1111} & FFFE \\ \hline
\end{tabular}
\end{center}
Because the last value in (\hex{2222}) is the first value out, AX and BX have swapped their values.

\subsection{Why do we need a stack?}
\begin{enumerate}
    \item \textbf{To keep return addresses}: when we call a procedure, the processor keeps in the stack where it must come back.
    \item \textbf{To keep registers for a while}: a procedure needs CX, but CX is used by a loop? We write
          \code{push cx} at its start and \code{pop cx} at its end, so the loop counter is not lost.
    \item \textbf{To pass parameters} to procedures (we see it soon).
    \item \textbf{Interrupts}: the processor keeps its state there before going to the interrupt handler.
\end{enumerate}

\begin{intabih}
    The number of \code{pop} must equal the number of \code{push}, in the opposite order. \code{push ax} then
    \code{push bx} must be matched by \code{pop bx} then \code{pop ax}. If you forget one \code{pop} inside a
    procedure, \code{ret} takes a wrong value as the return address, and the program jumps to a random place.
\end{intabih}

%================================================
\section{Functions and procedures}
%================================================

\begin{tarif}
    A \term{procedure}{procédure, sous-programme}: a group of instructions with a name. We call it with
    \code{call} and come back from it with \code{ret}. We write it once and use it many times.
\end{tarif}

\subsection{How does \code{ret} know where to go back?}
\begin{itemize}
    \item \code{call P}: puts on the stack \textbf{the address of the instruction after it} (IP), then jumps to P.
    \item \code{ret}: takes this address from the stack and puts it into IP, so execution goes back after \code{call}.
\end{itemize}

\begin{khatwa}[: the two calls in the next program (real addresses)]
    \begin{enumerate}
        \item The first \code{call print\_num} is at offset \hex{0103} and is 3 bytes long, so the next instruction is
              at \hex{0106}. The processor puts \hex{0106} on the stack, then jumps to \code{print\_num}.
        \item At the end of the procedure, \code{ret} takes \hex{0106} from the stack: execution goes back to \code{mov ax, 7}.
        \item The second call is at \hex{0109}, so it puts \hex{010C}, and this time \code{ret} goes back to \hex{010C}.
    \end{enumerate}
    The same procedure goes back to two different places, because the return address is kept in the stack each time.
\end{khatwa}

\subsection{A procedure that prints a number in decimal}
\code{int 21h} only prints characters. To print the number 2026, we must take out its digits one by one:

\begin{khatwa}[: taking out the digits of 2026]
    \begin{center}
    \begin{tabular}{|c|c|c|c|}
        \hline
        \ras{Operation} & \ras{Quotient (AX)} & \ras{Remainder (DX)} & \ras{Stack} \\ \hline
        $2026 \div 10$ & 202 & 6 & 6 \\ \hline
        $202 \div 10$ & 20 & 2 & 6, 2 \\ \hline
        $20 \div 10$ & 2 & 0 & 6, 2, 0 \\ \hline
        $2 \div 10$ & 0 (we stop) & 2 & 6, 2, 0, 2 \\ \hline
    \end{tabular}
    \end{center}
    The digits come out backwards (6, 2, 0, 2). We put them on the stack, then take them back: they come out in
    the right order 2, 0, 2, 6.
\end{khatwa}

\asm{ch3/procnum.asm}
\screen{ch3/procnum.out}

\begin{sharh}
    \s{4--5}{Put the number in AX (the parameter), then call the procedure.}
    \s{10}{\code{proc} and \code{endp} mark the start and the end of the procedure.}
    \s{14}{DX = 0 is needed before \code{div bx} (a 16-bit division divides DX:AX).}
    \s{16--17}{Put the digit on the stack, and count the digits in CX.}
    \s{21--25}{CX times: take a digit, turn it into a character by adding \code{'0'} (48), and print it.}
    \s{28}{\code{ret} goes back to line 6 the first time, and to line 8 the second time.}
\end{sharh}

\subsection{Passing parameters on the stack}
\asm{ch3/procstack.asm}

After \code{mov bp, sp} on line 14, the stack looks like this (values from the real run):

\begin{center}
\begin{tikzpicture}[c/.style={draw=brandDark, fill=white, minimum width=3.6cm, minimum height=0.6cm, font=\small},
    a/.style={font=\codefont\small, text=tealDeep}]
    \foreach \ad/\v/\l [count=\i from 0] in {FFFC/{30 (2nd parameter)}/{BP+6}, FFFA/{12 (1st parameter)}/{BP+4},
                                           FFF8/{return address 010B}/{BP+2}, FFF6/{old BP}/{BP}} {
        \node[c] (n\i) at (0,-0.6*\i) {\v};
        \node[a, left=0.15cm of n\i] {\ad};
        \node[font=\codefont\small\bfseries, right=0.15cm of n\i] {\l};
    }
    \draw[sahmO] (4.6,-1.8) node[right, font=\small] {SP = BP = FFF6} -- (3.0,-1.8);
\end{tikzpicture}
\end{center}

\begin{sharh}
    \s{4--7}{Put the two parameters on the stack: SP goes from \hex{FFFE} to \hex{FFFA}.}
    \s{8}{\code{call} puts the return address \hex{010B}: SP = \hex{FFF8}.}
    \s{13--14}{Save BP, then make BP stay fixed on the top of the stack (\hex{FFF6}).}
    \s{15--16}{The first parameter is at [BP+4], the second at [BP+6]: $12 + 30 = 42$.}
    \s{18}{\code{ret 4}: go back, then remove 4 bytes (the two parameters) from the stack. SP goes back to \hex{FFFE}.}
\end{sharh}

%================================================
\section{Interrupts}
%================================================

\begin{tarif}
    An \term{interrupt}{interruption}: an event that stops the current program for a moment, so the processor
    runs a special program (the \textbf{interrupt handler}), then comes back to continue where it stopped.
\end{tarif}

\begin{tashbih}
    You are reading a book and the \textbf{doorbell} rings. You put a bookmark in the page, open the door, then
    come back and continue reading from the same page. The bell is the interrupt, and the bookmark is what the
    processor keeps in the stack.
\end{tashbih}

\subsection{Kinds of interrupts}
\begin{center}
\small
\begin{tabularx}{\linewidth}{|>{\centering\arraybackslash}p{3cm}|R|R|}
    \hline
    \ras{Kind} & \ras{Who causes it?} & \ras{Examples} \\ \hline
    hardware & an outside device sends a signal on the INTR or NMI pin & a key press, the clock, the printer is ready \\ \hline
    software & the instruction \code{int n} in the program & \code{int 21h} for DOS services \\ \hline
    exceptions & the processor itself, after an error & division by zero (interrupt 0) \\ \hline
\end{tabularx}
\end{center}

\subsection{The interrupt vector table}
At the start of memory (from \hex{00000} to \hex{003FF}) there is a table of 256 entries. Each interrupt has
4 bytes: the offset, then the segment, of its handler.

\begin{tarif}
    \centering
    The vector of interrupt number $n$ is at address: \quad $4 \times n$
\end{tarif}

\begin{mithal}
    The vector of \code{int 21h}: $\hex{21} \times 4 = 33 \times 4 = 132 = \hex{84}$. So the address of the DOS
    handler is kept from \hex{00084} to \hex{00087}.
\end{mithal}

\subsection{What does the processor do on \code{int n}?}
\begin{enumerate}
    \item It puts FLAGS on the stack, then switches IF and TF off.
    \item It puts CS, then IP, on the stack (the bookmark in the book).
    \item It reads the new IP and CS from the entry $n \times 4$, so it jumps to the interrupt handler.
    \item At the end of the handler, \code{iret} takes IP, CS and FLAGS back from the stack, and the program continues.
\end{enumerate}

\subsection{A hardware interrupt, step by step}
\begin{khatwa}[: you press a key while a program runs]
    \begin{enumerate}
        \item The keyboard sends a signal to the \term{interrupt controller}{contrôleur d'interruptions} (PIC 8259).
        \item The controller sends a signal on the INTR input of the processor.
        \item If IF = 1, the processor first finishes the \textbf{current instruction}, then accepts the interrupt.
        \item The controller gives the processor the interrupt number (\hex{09} for the keyboard).
        \item The processor saves FLAGS, CS and IP, then jumps to the address written at $4 \times 9 = \hex{24}$.
        \item The handler reads the key from port \hex{60}, keeps it, then runs \code{iret}.
        \item The program continues from where it stopped, without ``noticing'' anything.
    \end{enumerate}
\end{khatwa}

\begin{tarif}
    \begin{itemize}
        \item \code{cli} (clear IF) blocks hardware interrupts for a while, and \code{sti} (set IF) allows them again.
              We use them when a sensitive part of the program must not be interrupted.
        \item The NMI input (Non Maskable Interrupt): an interrupt that \textbf{cannot be blocked}, used for serious
              errors such as a memory fault or a power failure.
    \end{itemize}
\end{tarif}

\subsection{A program: our own interrupt}
\asm{ch3/myint.asm}
\screen{ch3/myint.out}

\begin{sharh}
    \s{4--5}{ES = 0: we want to write in the vector table, which is in segment 0.}
    \s{6--8}{Write the address of \code{handler} in entry $\hex{60} \times 4 = \hex{180}$: the offset, then the segment.}
    \s{10--11}{Each \code{int 60h} runs \code{handler}: the first time it prints \code{A}, the second time \code{B}.}
    \s{19}{\code{iret} and not \code{ret}, because the interrupt also put FLAGS on the stack.}
\end{sharh}

%================================================
\section{Exercises}
%================================================

\begin{tamrin}
    Compute the physical address in each case: (a) DS = 3000H, BX = 0050H with \code{mov al, [bx]}.
    (b) SS = 4000H, SP = FFFEH.
\end{tamrin}
\begin{hall}
    (a) $\hex{30000} + \hex{0050} = \hex{30050}$. \quad (b) $\hex{40000} + \hex{FFFE} = \hex{4FFFE}$.
\end{hall}

\begin{tamrin}
    What is AX after these instructions? \code{mov ax, 0F0Fh}, then \code{and ax, 00FFh}, then \code{shl ax, 1}.
\end{tamrin}
\begin{hall}
    After \code{and}: \hex{000F}. After the shift (multiply by 2): \hex{001E} = 30.
\end{hall}

\begin{tamrin}
    Write a program that computes $5! = 1 \times 2 \times 3 \times 4 \times 5$ with the \code{loop} instruction, and puts the result in AX.
\end{tamrin}
\begin{hall}
    \asm{ch3/ex_fact.asm}
    The result: AX = 0078 = 120 (tested in the emulator).
\end{hall}

\begin{tamrin}
    After \code{mov sp, 0100h}, then \code{push ax}, then \code{push bx}, then \code{pop cx}: what is SP, and which value is in CX?
\end{tamrin}
\begin{hall}
    After the two pushes SP = \hex{00FC}, and after \code{pop} SP = \hex{00FE}. CX gets the last value in: the value of BX.
\end{hall}

\begin{tamrin}
    What is the vector address of interrupt number \hex{10} (BIOS screen services)?
\end{tamrin}
\begin{hall}
    $\hex{10} \times 4 = 16 \times 4 = 64 = \hex{40}$: the bytes from \hex{00040} to \hex{00043}.
\end{hall}

\begin{tamrin}
    Decode the instruction 89 D8.
\end{tamrin}
\begin{hall}
    89: \code{mov} on 16 bits, d = 0. D8 = \bin{11\,011\,000}: mod = 11 (register), reg = 011 (BX, the source),
    r/m = 000 (AX). So: \code{mov ax, bx}.
\end{hall}

\chapter{Case Study: 32-bit and 64-bit Processors}

\begin{khulasa}[: what will we learn in this chapter?]
\begin{itemize}
    \item Why we moved from 16 bits to 32 and then 64 bits, and what changed.
    \item How the architecture grew from IA-32 to AMD64.
    \item The registers of modern processors, and flat addressing.
    \item New instructions: SSE and AVX, which compute several numbers at the same time.
    \item The pipeline: how the processor runs several instructions at the same time.
\end{itemize}
\end{khulasa}

\begin{intabih}
    The programs of this chapter are written in \textbf{NASM} syntax for \textbf{Linux}, and we really ran them on an
    Intel Core i7-8665U processor (4 cores). To assemble and run:
    \code{nasm -f elf64 prog.asm}, then \code{ld prog.o -o prog}, then \code{./prog}.
\end{intabih}

%================================================
\section{From 16 to 32, then 64 bits}
%================================================

\begin{tarif}
    We say a processor is an \textbf{$n$-bit} processor when its \textbf{general registers} are $n$ bits wide: this is
    the size of the number it computes in one instruction, and the size of the address it can keep in a register.
\end{tarif}

\subsection{The main differences}
\begin{center}
\small
\begin{tabularx}{\linewidth}{|>{\centering\arraybackslash}p{3.2cm}|C|C|C|}
    \hline
    \ras{} & \ras{16 bits (8086)} & \ras{32 bits (IA-32)} & \ras{64 bits (x86-64)} \\ \hline
    biggest unsigned number & $2^{16}-1 = 65\,535$ & $2^{32}-1 \approx 4.3$ billion & $2^{64}-1 \approx 1.8 \times 10^{19}$ \\ \hline
    max memory & 1 MB (with segments) & 4 GB & 256 TB (48 bits in practice) \\ \hline
    general registers & 8 & 8 & \textbf{16} \\ \hline
    register names & AX, BX, \dots & EAX, EBX, \dots & RAX, RBX, \dots, R15 \\ \hline
    addressing & segmented & flat (in practice) & flat \\ \hline
\end{tabularx}
\end{center}

\subsection{Why do we need more than 32 bits?}
\begin{enumerate}
    \item \textbf{Memory}: with a 32-bit address we cannot use more than 4 GB. Today's computers have 16 or 32 GB,
          and servers have hundreds of gigabytes.
    \item \textbf{Performance}: a 64-bit register computes a big number in one instruction, while a 32-bit processor needs
          at least two instructions (like \code{add} and \code{adc} in Chapter 3). And twice as many registers (16 instead of 8)
          means fewer trips to memory.
    \item \textbf{Compatibility}: a 64-bit processor also runs old 32-bit programs without any change.
\end{enumerate}

\begin{khatwa}[: what does ``4 GB'' mean exactly?]
    \begin{enumerate}
        \item A 32-bit address can take $2^{32}$ different values.
        \item $2^{32} = 2^{2} \times 2^{30} = 4 \times 2^{30}$.
        \item And $2^{30}$ bytes $= 1$ GB, so the maximum size is 4 GB.
    \end{enumerate}
    With 48 bits: $2^{48} = 2^{8} \times 2^{40} = 256$ TB, that is 65536 times bigger!
\end{khatwa}

%================================================
\section{How the architecture grew: from IA-32 to AMD64}
%================================================

\begin{center}
\begin{tikzpicture}[e/.style={unit=#1, minimum width=2.9cm, minimum height=1.3cm, font=\small}]
    \node[e=brandMuted] (a) at (0,0) {8086 (1978)\\16 bits, segments};
    \node[e=brandPurple] (b) at (3.6,0) {80286 (1982)\\protected mode};
    \node[e=brandTeal] (c) at (7.2,0) {80386 (1985)\\\textbf{IA-32}};
    \node[e=brandOrange] (d) at (10.8,0) {AMD64 (2003)\\\textbf{64 bits}};
    \draw[sahmT] (a) -- (b); \draw[sahmT] (b) -- (c); \draw[sahmT] (c) -- (d);
\end{tikzpicture}
\end{center}

\begin{itemize}
    \item \textbf{IA-32} (Intel Architecture 32 bits): it started with the 80386. Registers became 32 bits wide
          (EAX: the E means \emph{Extended}), and the \textbf{protected mode} appeared: it stops programs from writing in each other's memory.
    \item \textbf{AMD64}: the surprise is that AMD (not Intel) designed the 64-bit extension in 2003, while keeping
          compatibility. Intel adopted it later as Intel 64. Today we call it x86-64.
\end{itemize}

\subsection{The working modes of the processor}
\begin{center}
\small
\begin{tabularx}{\linewidth}{|>{\centering\arraybackslash}p{3.3cm}|R|R|}
    \hline
    \ras{Mode} & \ras{Description} & \ras{When?} \\ \hline
    real mode & behaves like an 8086: 16 bits and 1 MB & when the computer starts, before the operating system \\ \hline
    protected mode & 32 bits, memory protection & 32-bit systems \\ \hline
    long mode & 64 bits, with a sub-mode to run 32-bit programs & all modern systems \\ \hline
\end{tabularx}
\end{center}

\begin{hal}
    Every modern computer wakes up ``thinking'' it is an 8086! It starts in real mode, then the operating system moves
    it to protected mode, then to long mode. This is how the architecture stays compatible since 1978.
\end{hal}

%================================================
\section{Registers in 32-bit and 64-bit processors}
%================================================

\subsection{Nested registers}
Each new register \textbf{contains} the old register in its right half (the low bits):

\begin{center}
\begin{tikzpicture}[b/.style={draw=brandDark, thick, minimum height=0.75cm, font=\codefont\bfseries}]
    \node[b, minimum width=12cm, fill=brandOrange!15, anchor=east] (r) at (6,1.6) {RAX: 64 bits};
    \node[b, minimum width=6cm, fill=brandTeal!15, anchor=east] (e) at (6,0.8) {EAX: 32 bits};
    \node[b, minimum width=3cm, fill=brandGreen!20, anchor=east] (x) at (6,0) {AX: 16};
    \node[b, minimum width=1.5cm, fill=brandPurple!20, anchor=east] (h) at (4.5,-0.8) {AH};
    \node[b, minimum width=1.5cm, fill=brandPurple!35, anchor=east] (l) at (6,-0.8) {AL};
    \foreach \x/\t in {6/0, 4.5/8, 3/16, 0/32, -6/64} \node[font=\codefont\scriptsize] at (\x,-1.45) {\t};
    \draw[brandMuted, dashed] (0,-1.25) -- (0,1.2); \draw[brandMuted, dashed] (3,-1.25) -- (3,0.4);
\end{tikzpicture}
\end{center}

\begin{center}
\small
\begin{tabular}{|c|c|c|c|l|}
    \hline
    \ras{64 bits} & \ras{32 bits} & \ras{16 bits} & \ras{8 bits} & \ras{Usual use} \\ \hline
    RAX & EAX & AX & AL & value returned by a function \\ \hline
    RBX & EBX & BX & BL & general \\ \hline
    RCX & ECX & CX & CL & counter, 4th parameter \\ \hline
    RDX & EDX & DX & DL & 3rd parameter \\ \hline
    RSI & ESI & SI & SIL & 2nd parameter \\ \hline
    RDI & EDI & DI & DIL & 1st parameter \\ \hline
    RSP & ESP & SP & SPL & top of the stack \\ \hline
    RBP & EBP & BP & BPL & base of the stack frame \\ \hline
    R8 \dots\ R15 & R8D \dots\ R15D & R8W \dots & R8B \dots & \textbf{new in 64 bits} \\ \hline
\end{tabular}
\end{center}

\begin{itemize}
    \item \textbf{RIP}: the instruction pointer (it was IP, then EIP).
    \item \textbf{RFLAGS}: the flags register; ZF, CF, SF, OF are in the same place as in the 8086.
    \item \textbf{Segment registers}: they still exist (CS, DS\dots) but their base is 0 in 64-bit mode, except FS and GS,
          which the system uses for data that belongs to each thread.
\end{itemize}

\subsection{A test: what happens to the other bits?}
\asm{ch4/regs64.asm}
\screen{ch4/regs64.out}

\begin{sharh}
    \s{5--6}{Put a 64-bit number in RAX and show it: 1122334455667788.}
    \s{9}{Writing EAX (32 bits) \textbf{clears} the high half: 00000000AABBCCDD.}
    \s{13}{Writing AX (16 bits) keeps the other bits: \dots 5566\textbf{1234}.}
    \s{17}{Writing AL changes only the last byte: \dots 77\textbf{FF}.}
    \s{24--42}{A procedure that shows RAX in hexadecimal: at each turn it rotates 4 bits, keeps them with the mask \code{0xF}, and turns them into a character.}
\end{sharh}

\begin{intabih}
    A strange but important rule: in 64-bit mode, \textbf{every write to a 32-bit register clears the high bits}, but a
    write to 8 or 16 bits does not. AMD designed this rule to make the processor faster.
\end{intabih}

%================================================
\section{Advanced addressing modes}
%================================================

\subsection{From segmented addressing to flat addressing}
\begin{center}
\begin{tikzpicture}
    \draw[thick, fill=brandLight] (-6,0) rectangle (-4,-4);
    \foreach \y/\c in {0/brandTeal, -1.3/brandOrange, -2.6/brandGreen} \fill[\c!30] (-6,\y) rectangle (-4,\y-1.2);
    \foreach \y/\t in {-0.6/CS, -1.9/DS, -3.2/SS} \node[font=\codefont] at (-5,\y) {\t};
    \node[font=\small\bfseries, align=center] at (-5,0.6) {segmented (8086)};
    \node[font=\scriptsize, text width=4cm, align=center] at (-5,-4.6) {zones of 64 KB, and address = segment:offset};
    \draw[thick, fill=brandTeal!15] (0,0) rectangle (2,-4);
    \node[font=\small, align=center] at (1,-2) {one single\\space};
    \node[font=\small\bfseries, align=center] at (1,0.6) {flat (x86-64)};
    \node[font=\scriptsize, text width=4cm, align=center] at (1,-4.6) {address = one number from 0 to the end};
    \draw[sahmO] (-3.6,-2) -- (-0.4,-2);
\end{tikzpicture}
\end{center}

In the \term{flat memory model}{modèle mémoire plat}, the program sees memory \textbf{as one long strip} of bytes,
numbered from 0. We do not need to compute ``segment $\times$ 16 + offset'': the address is just a number in a 64-bit register.

\begin{tashbih}
    Segmented addressing is like saying ``building 12, floor 3''. Flat addressing is like giving every flat in the
    town one unique number: ``flat number 73\,502''. Much simpler, as long as we have big enough numbers.
\end{tashbih}

\subsection{The general form of an address}
\begin{tarif}
    \centering
    $\text{address} = \text{base} + \text{index} \times \text{scale} + \text{displacement}$
    \par\small
    base and index: any two registers. scale: 1, 2, 4 or 8. displacement: a fixed number.
\end{tarif}

\begin{center}
\small
\begin{tabular}{|l|l|}
    \hline
    \ras{Writing (NASM)} & \ras{Use} \\ \hline
    \code{[rbx]} & base only: a pointer \\ \hline
    \code{[rbx + 16]} & base + displacement: a field in a structure (struct) \\ \hline
    \code{[tab + rcx*8]} & index $\times$ scale + displacement: element i of an array of 8-byte numbers \\ \hline
    \code{[rbx + rcx*4 + 8]} & everything together \\ \hline
    \code{[rel msg]} & relative to RIP: special to 64 bits \\ \hline
\end{tabular}
\end{center}

\begin{khatwa}[: where does \code{[tab + rcx*8]} read when RCX = 3?]
    \begin{enumerate}
        \item Each element of \code{tab} is a 64-bit number = 8 bytes (\code{dq}).
        \item Element number 3 (the fourth one, because we start at 0) is $3 \times 8 = 24$ bytes from the start.
        \item If \code{tab} is at address \hex{402000}, the address is $\hex{402000} + 24 = \hex{402018}$.
    \end{enumerate}
    The scale saves us from multiplying the index ourselves: the processor does it inside the instruction.
\end{khatwa}

\asm{ch4/sum64.asm}
\screen{ch4/sum64.out}

\begin{sharh}
    \s{5--6}{\code{xor} of a register with itself = 0: the sum and the index start at 0.}
    \s{8}{Addressing \textbf{with a scale}: read element i of the array, and add it to RAX.}
    \s{9--11}{\code{i++}, then ``while i < 5'' go back. This is a \code{while} loop, as in Chapter 3.}
    \s{18--36}{The same method as in Chapter 3 to print a number in decimal (divide by 10), but with 64-bit registers.}
    \s{39}{\code{dq} (define quad-word): 64-bit numbers. The first number ($10^{12}$) does not fit in 32 bits.}
\end{sharh}

\subsection{RIP-relative addressing}
In 64-bit mode we often write \code{lea rsi, [rel buf]}: the address is computed \textbf{relative to the current
instruction} ($\text{RIP} + \text{distance}$). The benefit: the program works wherever the system puts it in memory.
This is important for security (the system moves programs at each run so attackers cannot guess the addresses).

%================================================
\section{The instruction set}
%================================================

\subsection{The same instruction on 16, 32 and 64 bits}
We assembled the program below; these are the real machine-code bytes:

\asm{ch4/same.asm}

\begin{center}
\small
\begin{tabular}{|l|l|l|}
    \hline
    \ras{Instruction} & \ras{Machine code} & \ras{Note} \\ \hline
    \code{add ax, bx} & {\codefont \textcolor{orangeDeep}{66} 01 D8} & prefix 66: ``careful, 16 bits'' \\ \hline
    \code{add eax, ebx} & {\codefont 01 D8} & 32 bits is the default size \\ \hline
    \code{add rax, rbx} & {\codefont \textcolor{orangeDeep}{48} 01 D8} & prefix REX.W = 48: ``64 bits'' \\ \hline
    \code{add r8, r9} & {\codefont \textcolor{orangeDeep}{4D} 01 C8} & REX = 4D: 64 bits + new registers \\ \hline
\end{tabular}
\end{center}

\begin{khatwa}[: why is 01 D8 in the first three lines?]
    \begin{enumerate}
        \item 01 is the opcode ``add a register to a register or memory''. It is the same since the 8086.
        \item D8 = \bin{11\,011\,000} is the ModR/M byte we saw in Chapter 3: mod = 11 (register), reg = 011 (BX), r/m = 000 (AX).
        \item The size (16, 32 or 64) is not written in the instruction itself, but in a \textbf{prefix} before it. This is how
              the old instructions stayed valid, and the new sizes were added ``from the outside''.
        \item In 4D = \bin{0100\,1101}: 0100 means ``REX prefix'', then the bit W = 1 (64 bits), and two bits that add 8 to the
              two register numbers, so AX becomes R8 and BX becomes R9.
    \end{enumerate}
\end{khatwa}

\subsection{Comparison: a 32-bit program and a 64-bit program}
\textbf{In 32 bits:}
\asm{ch4/hello32.asm}
\screen{ch4/hello32.out}
\textbf{In 64 bits:}
\asm{ch4/hello64.asm}
\screen{ch4/hello64.out}

\begin{center}
\small
\begin{tabularx}{\linewidth}{|>{\centering\arraybackslash}p{3.8cm}|C|C|}
    \hline
    \ras{} & \ras{32 bits} & \ras{64 bits} \\ \hline
    asking the system for a service & \code{int 0x80} (an interrupt) & \code{syscall} (a special, faster instruction) \\ \hline
    number of the \emph{write} service & 4 in EAX & 1 in RAX \\ \hline
    parameters & EBX, ECX, EDX & RDI, RSI, RDX \\ \hline
    size of an address & 4 bytes & 8 bytes \\ \hline
\end{tabularx}
\end{center}

\begin{hal}
    We ran the two programs on the same 64-bit computer. The 32-bit program worked without any change: this is the
    compatibility sub-mode inside long mode.
\end{hal}

\subsection{What does the compiler write? From C to assembly}
We wrote a function in C, and asked the compiler gcc to give us the assembly it produces, with the command
\code{gcc -O1 -S -masm=intel sum.c}:

\begin{center}
\begin{minipage}[t]{0.47\linewidth}
    \cfile{ch4/sum.c}
\end{minipage}\hfill
\begin{minipage}[t]{0.49\linewidth}
    \asm{ch4/sum_gcc.s}
\end{minipage}
\end{center}

\begin{sharh}
    \s{1}{By the Linux 64-bit rules: the first parameter (\code{t}) arrives in RDI, and the second one (\code{n}) in RSI.}
    \s{2--3}{If $n \leq 0$ we jump to \code{.L4}: the result is 0.}
    \s{5}{\code{lea} computes the address of the end of the array $t + n \times 8$ \textbf{without reading memory}.}
    \s{8--11}{The loop: add the element, move 8 bytes forward, and stop at the end. The compiler used a pointer instead of the index \code{i}.}
    \s{13--14}{The result is always returned in RAX.}
\end{sharh}

\subsection{SSE and AVX: computing several numbers at once}
\begin{tarif}
    \term{SIMD}{une instruction, données multiples} (Single Instruction, Multiple Data): one instruction applies the
    same operation to \textbf{several numbers at the same time}.
\end{tarif}

\begin{center}
\begin{tikzpicture}[c/.style={draw=brandDark, thick, minimum width=1.3cm, minimum height=0.65cm, font=\codefont}]
    \foreach \v [count=\i from 0] in {1.5, 2.5, 3.5, 4.5} \node[c, fill=brandTeal!15] at (1.3*\i,1.6) {\v};
    \foreach \v [count=\i from 0] in {10, 20, 30, 40} \node[c, fill=brandOrange!15] at (1.3*\i,0.6) {\v};
    \foreach \v [count=\i from 0] in {11.5, 22.5, 33.5, 44.5} \node[c, fill=brandGreen!20] at (1.3*\i,-0.9) {\v};
    \draw[very thick] (-0.7,0.1) -- (4.6,0.1);
    \node[font=\codefont\bfseries] at (-1.1,0.6) {+};
    \node[font=\small, anchor=west] at (4.8,1.6) {XMM0 (128 bits = 4 numbers)};
    \node[font=\small, anchor=west] at (4.8,-0.9) {result: \textbf{one instruction} \code{addps}};
\end{tikzpicture}
\end{center}

\begin{center}
\small
\begin{tabular}{|c|c|c|c|c|}
    \hline
    \ras{Extension} & \ras{Year} & \ras{Registers} & \ras{Size} & \ras{floats at once} \\ \hline
    SSE & 1999 & XMM0--XMM15 & 128 bits & 4 \\ \hline
    AVX & 2011 & YMM0--YMM15 & 256 bits & 8 \\ \hline
    AVX-512 & 2016 & ZMM0--ZMM31 & 512 bits & 16 \\ \hline
\end{tabular}
\end{center}

\cfile{ch4/simd.c}
\screen{ch4/simd.out}

When we open the compiled program, we really find the instruction:
\begin{codebox}\codefont\small
    addps \ xmm0, XMMWORD PTR [rbp-0x50] \quad \textcolor{red!75!black}{; 4 additions at once}
\end{codebox}

And the AVX version adds 8 numbers with the instruction \code{vaddps ymm0, ymm0, [\dots]}:
\cfile{ch4/avx.c}
\screen{ch4/avx.out}

\begin{tashbih}
    A normal instruction is like a teacher who corrects one exam paper at a time. SIMD is like a teacher who puts
    8 papers side by side and corrects question 1 on all of them with one move. This is why it is used in image and
    video processing, games, and artificial intelligence, where we apply the same operation to millions of numbers.
\end{tashbih}

%================================================
\section{The instruction cycle and the pipeline}
%================================================

\subsection{The five stages}
In modern processors, each instruction is split into stages, and each stage has its own circuit:
\begin{center}
\small
\begin{tabular}{|c|c|l|}
    \hline
    \ras{Stage} & \ras{Name} & \ras{Work} \\ \hline
    IF & fetch & read the instruction from memory (address in RIP) \\ \hline
    ID & decode & understand the instruction and read the registers \\ \hline
    EX & execute & compute the result in the UAL, or compute an address \\ \hline
    MEM & memory & read or write memory if needed \\ \hline
    WB & write back & store the result in the register \\ \hline
\end{tabular}
\end{center}

\subsection{Without a pipeline and with a pipeline}
\begin{tashbih}
    In a laundry: washing takes one hour, drying one hour, ironing one hour. To wash 3 bags one after the other we
    need 9 hours. But when the first bag goes from the washing machine to the dryer, we put the second bag in the
    washing machine at once. The three machines work together, so we finish in only 5 hours. This is a pipeline.
\end{tashbih}

\begin{center}
\begin{tikzpicture}[s/.style={draw=brandDark, minimum width=0.95cm, minimum height=0.55cm, font=\codefont\scriptsize\bfseries}]
    \node[font=\small\bfseries] at (4.3,0.9) {with a pipeline: 4 instructions in 8 cycles};
    \foreach \c in {1,...,8} \node[font=\codefont\scriptsize] at (0.95*\c, 0.35) {\c};
    \foreach \i/\col in {1/brandTeal, 2/brandOrange, 3/brandGreen, 4/brandPurple} {
        \node[font=\small] at (-0.3,-0.6*\i+0.1) {I\i};
        \foreach \st [count=\k from 0] in {IF, ID, EX, MEM, WB} {
            \pgfmathsetmacro\x{0.95*(\i+\k)}
            \node[s, fill=\col!25] at (\x, -0.6*\i+0.1) {\st};
        }
    }
\end{tikzpicture}
\end{center}

\begin{tarif}
    For $n$ instructions and $k$ stages (one cycle per stage):
    \begin{itemize}
        \item without a pipeline: $n \times k$ cycles.
        \item with a pipeline: $k + (n - 1)$ cycles: the first one needs $k$ cycles, then one instruction finishes in \textbf{every} cycle.
    \end{itemize}
\end{tarif}

\begin{khatwa}[: 100 instructions in a 5-stage pipeline]
    \begin{enumerate}
        \item Without a pipeline: $100 \times 5 = 500$ cycles.
        \item With a pipeline: $5 + (100 - 1) = 104$ cycles.
        \item The speed-up: $\dfrac{500}{104} \approx 4.8$ times, close to the number of stages (5).
    \end{enumerate}
\end{khatwa}

\subsection{Hazards: when does the pipeline stop?}
\begin{enumerate}
    \item \textbf{Data hazard}: \code{add rax, rbx} then \code{add rcx, rax}. The second instruction needs the new RAX,
          but the first one has not written it yet. The fix: the processor ``forwards'' the result directly from EX to the
          second instruction, or waits one cycle.
    \item \textbf{Control hazard}: after \code{jmp} or \code{jne}, the instructions fetched after it may be wrong. The fix:
          \term{branch prediction}{prédiction de branchement}: the processor guesses if it will jump, and if it guessed wrong,
          it throws away the wrong instructions and starts again.
    \item \textbf{Structural hazard}: two instructions need the same circuit in the same cycle (like memory).
\end{enumerate}

\subsection{Example: \code{mov}, \code{add} and \code{jmp} in 64 bits}
We follow three instructions from the program \code{same.asm} and the programs of this chapter, with their real bytes:

\begin{center}
\small
\begin{tabularx}{\linewidth}{|>{\centering\arraybackslash}p{1.1cm}|>{\raggedright\arraybackslash}p{3.3cm}|R|R|}
    \hline
    \ras{Stage} & \ras{\code{mov rax, 60}} & \ras{\code{add rax, rbx}} & \ras{\code{jmp fin}} \\ \hline
    IF & fetch 7 bytes 48 C7 C0 3C 00 00 00 & fetch 3 bytes 48 01 D8 & fetch 2 bytes EB 00 \\ \hline
    ID & REX.W: 64 bits, destination RAX, the value 60 is in the instruction & 64 bits, source RBX, destination RAX; read both registers & short jump, offset 0 \\ \hline
    EX & nothing to compute (the value is ready) & UAL: $\text{RAX} + \text{RBX}$, and compute the flags & compute the target: $\text{RIP} + 0$ \\ \hline
    MEM & nothing & nothing & nothing \\ \hline
    WB & \nql{RAX}{60} & RAX gets the result; update RFLAGS & RIP gets the target; throw away wrongly fetched instructions \\ \hline
\end{tabularx}
\end{center}

\begin{hal}
    The Core i7 we used does much more than a 5-stage pipeline: it has more than 14 stages, and it can start
    \textbf{4 instructions in every cycle}. It even runs them \textbf{out of order} when they do not depend on each
    other, then puts the results back in order. All this happens without the programmer seeing any difference.
\end{hal}

%================================================
\section{Exercises}
%================================================

\begin{tamrin}
    How many bytes can a processor address with a 36-bit address? (This is what the PAE extension allowed in 32-bit processors.)
\end{tamrin}
\begin{hall}
    $2^{36} = 2^{6} \times 2^{30} = 64$ GB.
\end{hall}

\begin{tamrin}
    RAX = \hex{FFFFFFFFFFFFFFFF}. What is RAX after: (a) \code{mov eax, 1} (b) \code{mov ax, 1} (c) \code{mov al, 1}?
\end{tamrin}
\begin{hall}
    (a) \hex{0000000000000001} (writing 32 bits clears the high half).\\
    (b) \hex{FFFFFFFFFFFF0001}. \quad (c) \hex{FFFFFFFFFFFFFF01}.
\end{hall}

\begin{tamrin}
    An array of \code{int} (4 bytes each) starts at \hex{601000}. What is the address of element number 10? Write the addressing in NASM.
\end{tamrin}
\begin{hall}
    $\hex{601000} + 10 \times 4 = \hex{601000} + 40 = \hex{601028}$. The writing: \code{[tab + rcx*4]} with RCX = 10.
\end{hall}

\begin{tamrin}
    A 6-stage pipeline, and 1000 instructions. Compute the number of cycles without and with a pipeline, and the speed-up.
\end{tamrin}
\begin{hall}
    Without: $1000 \times 6 = 6000$. With: $6 + 999 = 1005$. Speed-up $= \dfrac{6000}{1005} \approx 5.97$.
\end{hall}

\begin{tamrin}
    In this program, where is the data hazard? \code{mov rax, [x]}, then \code{add rax, 5}, then \code{mov [y], rax}.
\end{tamrin}
\begin{hall}
    The second instruction needs the RAX that the first one reads from memory, and the third one needs the RAX that the
    second one computes: each instruction depends on the one just before it.
\end{hall}

\chapter{Modern Processor Architectures}

\begin{khulasa}[: what will we learn in this chapter?]
\begin{itemize}
    \item How processors grew, and why the frequency stopped going up.
    \item Multi-core processors and parallel processing: the benefits and the problems, and Amdahl's law.
    \item Cache memory and the memory hierarchy, and why the order in which we read data matters so much.
\end{itemize}
\end{khulasa}

\begin{intabih}
    All the numbers in this chapter were \textbf{really measured} on a laptop with an Intel Core i7-8665U processor
    (4 cores, 8 threads, L1 cache = 32 KB and L2 = 256 KB per core, and a shared L3 = 8 MB). If you repeat the tests on
    your computer, your numbers will be a little different, but the effects will be the same.
\end{intabih}

%================================================
\section{Modern processors: trends and changes}
%================================================

\subsection{Fifty years of progress}
\begin{center}
\small
\begin{tabular}{|c|c|c|c|c|}
    \hline
    \ras{Processor} & \ras{Year} & \ras{Transistors} & \ras{Frequency} & \ras{Cores} \\ \hline
    Intel 4004 & 1971 & 2\,300 & 0.74 MHz & 1 \\ \hline
    Intel 8086 & 1978 & 29\,000 & 5--10 MHz & 1 \\ \hline
    Intel 80386 & 1985 & 275\,000 & 16--33 MHz & 1 \\ \hline
    Pentium 4 & 2000 & 42 million & 1.3--3.8 GHz & 1 \\ \hline
    Core i7-8665U & 2019 & about 2 billion & 1.9--4.8 GHz & 4 \\ \hline
    processors of 2024 & 2024 & tens of billions & about 5 GHz & from 8 to more than 100 \\ \hline
\end{tabular}
\end{center}

\subsection{The frequency wall}
Until about 2004, each year brought a faster processor because the \textbf{frequency} went up. Then the frequency
stopped at about 4 to 5 GHz. The main reason is \textbf{heat}:
\begin{itemize}
    \item Every clock pulse uses energy. The higher the frequency, the more energy is used, and energy turns into heat.
    \item At about 4 GHz the processor gets so hot that a normal computer cannot cool it.
\end{itemize}

\begin{tashbih}
    A runner runs faster and faster, until he reaches a speed he cannot keep without falling. Instead of asking him to
    run faster, we add other runners who share the distance with him. This is exactly what processor makers did.
\end{tashbih}

\subsection{Today's trends}
\begin{enumerate}
    \item \textbf{More cores} instead of a higher frequency (next section).
    \item \textbf{Bigger and smarter cache memory}, because memory became much slower than the processor (last section).
    \item \textbf{Special instructions}: SIMD (Chapter 4), and speed-up for artificial intelligence and encryption.
    \item \textbf{Different cores on the same chip}: strong cores for heavy work and economical cores for light work, to save battery.
\end{enumerate}

%================================================
\section{Multi-core processors and parallelism}
%================================================

\subsection{The ideas}
\begin{tarif}
    \begin{itemize}
        \item A \term{core}{cœur}: a complete processor (control unit, arithmetic unit, registers, L1 cache).
        \item A \term{multi-core processor}{processeur multi-cœur}: several cores on one chip.
        \item A \term{thread}{thread, fil d'exécution}: a series of instructions that runs on its own. A program may have one thread or several.
        \item \term{Parallel processing}{traitement parallèle}: running several parts of the work \textbf{at the same moment} on different cores.
    \end{itemize}
\end{tarif}

\begin{center}
\begin{tikzpicture}[core/.style={unit, minimum width=2.4cm, minimum height=1.5cm}]
    \draw[brandDark, very thick, rounded corners=8pt, fill=brandLight] (-6.1,-2.6) rectangle (6.1,2.1);
    \node[font=\bfseries] at (0,1.75) {a processor chip with 4 cores};
    \foreach \x/\i in {-4.5/1, -1.5/2, 1.5/3, 4.5/4} {
        \node[core] at (\x,0.6) {core \i\\{\scriptsize L1, L2}};
    }
    \node[unit=brandOrange, minimum width=11.2cm] at (0,-1.1) {shared cache memory L3};
    \node[unit=brandPurple, minimum width=4cm] (ram) at (0,-3.6) {main memory RAM};
    \draw[sahmT, <->] (0,-1.6) -- (ram);
\end{tikzpicture}
\end{center}

\begin{tashbih}
    A kitchen with one cook (one core) prepares one dish after another. A kitchen with 4 cooks (4 cores) prepares 4 dishes
    at the same time. But they share the same fridge (L3 and RAM), and one of them may wait for his turn in front of it.
    And one complex dish only gets faster if we can split it into independent tasks.
\end{tashbih}

\subsection{Hyper-Threading: one core for two threads}
\term{Simultaneous multithreading}{multithreading simultané} (SMT, or Hyper-Threading) makes one core look like two cores:
when a thread waits for data from memory, the second thread uses the free arithmetic units. This is why our processor shows
8 ``logical processors'' even though it has only 4 cores. But if the two threads need the same arithmetic unit all the time,
we gain almost nothing, as we will see in the test.

\subsection{A test: the same work on 1, 2, 4 and 8 threads}
The program below computes the sum of the square roots of 400 million numbers, and shares the work between $k$ threads:

\cfile{ch5/threads.c}
\screen{ch5/threads.out}

\begin{sharh}
    \s{18--26}{Each thread takes the numbers $id$, $id + k$, $id + 2k$\dots so it computes a part of the sum that does not overlap the others.}
    \s{34--35}{\code{pthread\_create}: start $k$ threads; each one runs on a core.}
    \s{36--39}{\code{pthread\_join}: wait for each thread to finish, then add the partial results.}
    \s{40--42}{Speed-up = time with one thread $\div$ time with $k$ threads.}
\end{sharh}

\begin{center}
\begin{tikzpicture}[x=1.1cm, y=1.1cm]
    \draw[->, thick] (0,0) -- (8.6,0) node[right, font=\small] {threads};
    \draw[->, thick] (0,0) -- (0,4.4) node[above, font=\small] {speed-up};
    \foreach \y in {1,2,3,4} \draw (0.08,\y) -- (-0.08,\y) node[left, font=\scriptsize] {\y};
    \foreach \x in {1,2,4,8} \draw (\x,0.08) -- (\x,-0.08) node[below, font=\scriptsize] {\x};
    \draw[brandMuted, dashed, thick] (1,1) -- (4,4) node[above, font=\scriptsize, text=brandMuted] {ideal speed-up};
    \draw[orangeDeep, very thick] (1,1) -- (2,1.98) -- (4,3.02) -- (8,3.05);
    \foreach \x/\y in {1/1, 2/1.98, 4/3.02, 8/3.05} {
        \fill[orangeDeep] (\x,\y) circle (2.5pt);
        \node[font=\codefont\scriptsize, below right] at (\x,\y) {\y};
    }
    \draw[brandTeal, dashed] (4,0) -- (4,4.2);
    \node[font=\scriptsize, text=tealDeep, align=center] at (6,3.9) {only 4 cores:\\extra threads do not help here};
\end{tikzpicture}
\end{center}

\begin{khatwa}[: how do we read these results?]
    \begin{enumerate}
        \item \textbf{2 threads}: speed-up 1.98, very close to 2. The work was split into two fully independent halves.
        \item \textbf{4 threads}: speed-up 3.02, not 4. When the four cores work together the processor gets hot, so it lowers
              its frequency by itself. And a part of the program (starting the threads and adding the results) is not parallel.
        \item \textbf{8 threads}: speed-up 3.05, almost the same as 4 threads. Our processor has only 4 real cores, and the
              extra threads share the same arithmetic units, so the work does not go faster.
    \end{enumerate}
\end{khatwa}

\subsection{Amdahl's law: the limits of parallelism}
Most programs have a part that \textbf{cannot be shared} (reading the file, showing the result, waiting for the user\dots).

\begin{tarif}
    If the parallel part of the program is $p$ (between 0 and 1), and we have $n$ cores, the maximum speed-up is:
    \[ S(n) = \frac{1}{(1 - p) + \dfrac{p}{n}} \]
\end{tarif}

\begin{khatwa}[: a program that is 90\% parallel]
    $p = 0.9$, and the serial part $1 - p = 0.1$.
    \begin{enumerate}
        \item With 4 cores: $S = \dfrac{1}{0.1 + \frac{0.9}{4}} = \dfrac{1}{0.1 + 0.225} = \dfrac{1}{0.325} \approx 3.08$.
        \item With 16 cores: $S = \dfrac{1}{0.1 + 0.05625} \approx 6.4$.
        \item With an infinite number of cores: $\dfrac{p}{n} \rightarrow 0$, so $S = \dfrac{1}{0.1} = 10$.
    \end{enumerate}
    However many cores we add, the program will never be more than 10 times faster, because the 10\% serial part stays the same.
\end{khatwa}

\begin{hal}
    Notice that Amdahl's law with $p = 0.9$ and 4 cores gives 3.08, very close to the speed-up of 3.02 we measured!
    A simple law explains reality well.
\end{hal}

\subsection{Benefits and challenges}
\begin{center}
\small
\begin{tabularx}{\linewidth}{|R|R|}
    \hline
    \ras{Benefits} & \ras{Challenges} \\ \hline
    more performance without raising the frequency & the program must be written with several threads, which is harder \\ \hline
    less energy for the same work & the serial part limits the speed-up (Amdahl's law) \\ \hline
    several programs run smoothly at the same time & \textbf{synchronization}: two threads writing the same variable can break the result \\ \hline
    if one task fails, the other cores keep working & \textbf{cache coherence}: all the cores must see the same value of a variable \\ \hline
\end{tabularx}
\end{center}

\begin{mithal}[: the synchronization problem]
    Two threads run \code{x = x + 1} on the same variable, and x = 5 at the start:
    \begin{enumerate}
        \item Thread 1 reads x = 5.
        \item Thread 2 also reads x = 5 (before the first one writes).
        \item Thread 1 computes $5 + 1$ and writes 6.
        \item Thread 2 computes $5 + 1$ and also writes 6.
    \end{enumerate}
    The result is 6 instead of 7: one increase is lost! This is why we use \textbf{locks} (mutex) or \textbf{atomic}
    instructions (\code{lock add} in x86) that stop two threads from changing the variable at the same time. In our program
    we avoided the problem: each thread writes into its own box \code{part[id]}.
\end{mithal}

%================================================
\section{Cache memory and the memory hierarchy}
%================================================

\subsection{The problem: memory is very slow}
The processor can run several instructions every nanosecond, but reading the main memory takes about 100 nanoseconds.
If the processor waited for memory every time, it would do nothing most of the time.

\begin{tarif}
    \term{Cache memory}{mémoire cache, antémémoire}: a small and very fast memory (SRAM) close to the processor. It keeps a
    \textbf{copy} of the data used recently, or of the data we will need soon.
\end{tarif}

\begin{tashbih}
    When you revise for an exam, you do not put all the books of the library on your desk. On the desk you put the books you
    need now (L1), in your bag the books of the week (L2), in your cupboard at home the books of the year (L3), and the rest
    stays in the far library (RAM). The closer the place, the smaller and the faster it is.
\end{tashbih}

\subsection{The memory hierarchy}
\begin{center}
\begin{tikzpicture}
    \foreach \y/\w/\t/\c/\d in {
        3.6/1.6/{registers}/brandOrange/{less than 1 KB, less than 1 ns},
        2.7/2.8/{L1}/brandOrange!70!brandTeal/{32 KB, about 1 ns},
        1.8/4.0/{L2}/brandTeal/{256 KB, about 4 ns},
        0.9/5.2/{L3}/brandTeal!70!brandPurple/{8 MB, about 10 ns},
        0.0/6.4/{main memory RAM}/brandPurple/{8 GB, about 100 ns},
        -0.9/7.6/{SSD disk}/brandMuted/{256 GB, about 100 µs}} {
        \fill[\c!35, draw=\c, thick] (-\w/2,\y) rectangle (\w/2,\y+0.8);
        \node[font=\small\bfseries] at (0,\y+0.4) {\t};
        \node[font=\scriptsize, anchor=west] at (4.1,\y+0.4) {\d};
    }
    \draw[sahmO] (-4.3,-0.9) -- (-4.3,4.4) node[midway, left, font=\small, align=right] {faster\\smaller\\more expensive};
\end{tikzpicture}
\end{center}

\subsection{A test: we can see the memory levels}
The program below reads an array in a random order (each read gives the address of the next read, so the processor cannot
guess), and measures the average time of one read, for different array sizes:

\cfile{ch5/latency.c}
\screen{ch5/latency.out}

\begin{center}
\begin{tikzpicture}[x=1.45cm, y=0.042cm]
    \draw[->, thick] (0,0) -- (7.4,0) node[right, font=\small] {data size};
    \draw[->, thick] (0,0) -- (0,125) node[above, font=\small] {time (ns)};
    \foreach \y in {0,25,50,75,100} \draw (0.06,\y) -- (-0.06,\y) node[left, font=\scriptsize] {\y};
    \foreach \x/\t in {1/16 KB, 2/64 KB, 3/256 KB, 4/1 MB, 5/4 MB, 6/16 MB, 7/64 MB}
        \node[font=\codefont\scriptsize, below] at (\x,-2) {\t};
    \fill[brandOrange!15] (0.5,0) rectangle (1.6,125); \node[font=\scriptsize] at (1.05,118) {L1};
    \fill[brandTeal!15] (1.6,0) rectangle (3.4,125); \node[font=\scriptsize] at (2.5,118) {L2};
    \fill[brandPurple!12] (3.4,0) rectangle (5.6,125); \node[font=\scriptsize] at (4.5,118) {L3};
    \fill[brandMuted!12] (5.6,0) rectangle (7.3,125); \node[font=\scriptsize] at (6.45,118) {RAM};
    \draw[orangeDeep, very thick] (1,1.7) -- (2,2.2) -- (3,4.1) -- (4,9.2) -- (5,12.8) -- (6,83.5) -- (7,112.7);
    \foreach \x/\y in {1/1.7, 2/2.2, 3/4.1, 4/9.2, 5/12.8, 6/83.5, 7/112.7} {
        \fill[orangeDeep] (\x,\y) circle (2.5pt);
        \node[font=\codefont\scriptsize, above] at (\x,\y+2) {\y};
    }
\end{tikzpicture}
\end{center}

\begin{khatwa}[: what do we see in the curve?]
    \begin{enumerate}
        \item \textbf{Up to 32 KB}: the whole array is in L1, and a read takes only 1.7 ns.
        \item \textbf{64 KB and 256 KB}: the array does not fit in L1 any more, so we read from L2: the time goes up to 2 then 4 ns.
        \item \textbf{1 MB and 4 MB}: we read from the shared L3: about 9 to 13 ns.
        \item \textbf{16 MB and more}: the array is bigger than L3 (8 MB), so we go to RAM: the time jumps to 84 then 113 ns,
              about \textbf{65 times} slower than L1!
    \end{enumerate}
    The steps in the curve fall exactly at the cache sizes of our processor.
\end{khatwa}

\subsection{The principle of locality: why does cache memory work?}
Cache memory is small, but it usually holds what we need, thanks to two natural habits of programs:
\begin{itemize}
    \item \term{Temporal locality}{localité temporelle}: the data we used now will often be used \textbf{again soon}
          (the loop counter, the sum\dots).
    \item \term{Spatial locality}{localité spatiale}: after reading an address, we often read the \textbf{nearby addresses}
          (array elements, instructions that follow each other).
\end{itemize}

\begin{tarif}
    A \term{cache line}{ligne de cache}: the cache does not move one byte, but a whole block of \textbf{64 bytes} next to
    each other. When we read \code{t[0]}, \code{t[1]} to \code{t[15]} (for 4-byte numbers) come with it for free.
\end{tarif}

\subsection{A test: the order of reading changes everything}
\cfile{ch5/locality.c}
\screen{ch5/locality.out}

\begin{center}
\begin{tikzpicture}[c/.style={draw=brandDark, minimum size=0.55cm}]
    \foreach \i in {0,...,3} \foreach \j in {0,...,3} \node[c] at (\j*0.55-6.0, -\i*0.55) {};
    \foreach \i in {0,...,3} \draw[sahmT] (-6.2, -\i*0.55) -- (-4.1, -\i*0.55);
    \node[font=\small, align=center] at (-5.2,0.75) {row by row: each read\\is next to the one before};
    \foreach \i in {0,...,3} \foreach \j in {0,...,3} \node[c] at (\j*0.55+3.0, -\i*0.55) {};
    \foreach \j in {0,...,3} \draw[sahmO] (\j*0.55+3.0, 0.2) -- (\j*0.55+3.0, -1.85);
    \node[font=\small, align=center] at (3.8,0.75) {column by column: each read\\jumps 16 KB};
\end{tikzpicture}
\end{center}

\begin{khatwa}[: why are columns 33 times slower?]
    \begin{enumerate}
        \item In C, the array \code{m[4096][4096]} is stored \textbf{row after row}: \code{m[0][0]}, then \code{m[0][1]}, then\dots
        \item \textbf{Row by row}: each cache line (64 bytes) holds 16 numbers next to each other. We read from RAM once, then use
              16 numbers from the cache. And the processor notices we move forward regularly, so it fetches the next lines in advance.
        \item \textbf{Column by column}: \code{m[0][j]} and \code{m[1][j]} are $4096 \times 4 = 16\,384$ bytes apart. Every read is
              in a new cache line, so we use only one number out of the 16 we moved.
        \item The real result: 0.009 s against 0.301 s, about \textbf{33 times} slower, for exactly the same operations!
    \end{enumerate}
\end{khatwa}

\subsection{Hits, misses and average access time}
\begin{tarif}
    \begin{itemize}
        \item A \term{hit}{succès}: the data is in the cache.
        \item A \term{miss}{défaut}: the data is not in the cache, so we go to the next (slower) level.
        \item The \textbf{hit rate} $h$: the share of reads that found their data in the cache.
    \end{itemize}
    \[ \text{average time} = h \times t_{cache} + (1 - h) \times t_{RAM} \]
\end{tarif}

\begin{khatwa}[: $t_{cache} = 1$ ns, $t_{RAM} = 100$ ns]
    \begin{enumerate}
        \item Hit rate 90\%: $0.9 \times 1 + 0.1 \times 100 = 0.9 + 10 = 10.9$ ns.
        \item Hit rate 99\%: $0.99 \times 1 + 0.01 \times 100 = 0.99 + 1 = 1.99$ ns.
    \end{enumerate}
    Going from 90\% to 99\% makes memory \textbf{more than 5 times faster}. This is why modern processors work hard to improve the hit rate.
\end{khatwa}

\subsection{Where do we put data in the cache? (mapping)}
The cache is split into lines. We must decide in which line we put a block coming from memory:

\begin{center}
\small
\begin{tabularx}{\linewidth}{|>{\centering\arraybackslash}p{3.4cm}|R|R|}
    \hline
    \ras{Kind} & \ras{Rule} & \ras{Advantage and drawback} \\ \hline
    direct mapped & each block has only one possible line: block number mod number of lines & simple and fast, but two blocks may fight for the same line \\ \hline
    fully associative & any block in any line & fewer conflicts, but searching all the lines is expensive \\ \hline
    set associative & the block goes to a given set, and inside it to any of $k$ lines & the middle way, used in all modern processors \\ \hline
\end{tabularx}
\end{center}

\begin{khatwa}[: splitting an address in a direct-mapped cache]
    A 1 KB cache, with 16-byte lines, and 16-bit addresses.
    \begin{enumerate}
        \item Number of lines: $1024 \div 16 = 64 = 2^6$ lines, so we need \textbf{6 bits} for the line number.
        \item Inside the line there are $16 = 2^4$ bytes, so \textbf{4 bits} for the offset.
        \item The rest: $16 - 6 - 4 = 6$ bits, called the \textbf{tag}: it tells which memory block is in the line now.
    \end{enumerate}
    \begin{center}
    \begin{tikzpicture}[f/.style={draw=brandDark, thick, minimum height=0.7cm, font=\small\bfseries}]
        \node[f, minimum width=3cm, fill=brandOrange!20] (t) at (0,0) {tag: 6 bits};
        \node[f, minimum width=3cm, fill=brandTeal!20, anchor=west] (l) at (t.east) {line: 6 bits};
        \node[f, minimum width=2cm, fill=brandGreen!20, anchor=west] (o) at (l.east) {offset: 4 bits};
        \node[font=\codefont\scriptsize, below=1pt of t.south west, anchor=north west] {bit 15};
        \node[font=\codefont\scriptsize, below=1pt of o.south east, anchor=north east] {bit 0};
    \end{tikzpicture}
    \end{center}
    Example: $\hex{1A2C} = \bin{000110\,100010\,1100}$. Tag $\bin{000110} = 6$, line $\bin{100010} = 34$, and offset $\bin{1100} = 12$.
\end{khatwa}

\subsection{Replacement and writing}
\begin{itemize}
    \item \textbf{Replacement policy}: when a set is full, which line do we throw out? Usually \textbf{LRU} (Least Recently Used):
          the line that has not been used for the longest time.
    \item \textbf{Write policy}:
    \begin{itemize}
        \item \textbf{Write-through}: we write in the cache and in RAM together. Simple but slow.
        \item \textbf{Write-back}: we write only in the cache, and write in RAM only when the line is thrown out. Faster, and used today.
    \end{itemize}
\end{itemize}

\begin{hal}
    Remember the Harvard model from Chapter 1: the L1 cache is split into two parts, L1i for instructions and L1d for data
    (32 KB each in our processor), so the processor can fetch an instruction and read data at the same time. But L2, L3 and RAM
    are shared by both, as in the Von Neumann model.
\end{hal}

%================================================
\section{Exercises}
%================================================

\begin{tamrin}
    A program is 75\% parallel. Compute the speed-up with 2, 4 and 8 cores, and the maximum speed-up.
\end{tamrin}
\begin{hall}
    $1 - p = 0.25$. \quad $S(2) = \dfrac{1}{0.25 + 0.375} = 1.6$, \quad $S(4) = \dfrac{1}{0.25 + 0.1875} \approx 2.29$,
    \quad $S(8) = \dfrac{1}{0.25 + 0.09375} \approx 2.91$. Maximum $= \dfrac{1}{0.25} = 4$.
\end{hall}

\begin{tamrin}
    A cache with a time of 2 ns, a main memory of 80 ns, and a hit rate of 95\%. What is the average time?
\end{tamrin}
\begin{hall}
    $0.95 \times 2 + 0.05 \times 80 = 1.9 + 4 = 5.9$ ns.
\end{hall}

\begin{tamrin}
    A direct-mapped cache: 4 KB, 32-byte lines, 32-bit addresses. How many bits for the offset, the line and the tag?
\end{tamrin}
\begin{hall}
    Offset: $32 = 2^5$, so 5 bits. Lines: $4096 \div 32 = 128 = 2^7$, so 7 bits. Tag: $32 - 7 - 5 = 20$ bits.
\end{hall}

\begin{tamrin}
    Explain why this code is slow, and how to fix it:
    \begin{codebox}\codefont\small
    for (j = 0; j < N; j++) \\
    \hspace*{2em}for (i = 0; i < N; i++) \\
    \hspace*{4em}s += m[i][j];
    \end{codebox}
\end{tamrin}
\begin{hall}
    It reads the array column by column, so every read is in a different cache line (no spatial locality). We fix it by
    swapping the two loops: \code{for (i ...) for (j ...)}, so we read row by row.
\end{hall}

\begin{tamrin}
    Why did the program \code{threads.c} not get better when going from 4 to 8 threads on our processor?
\end{tamrin}
\begin{hall}
    Because the processor has only 4 real cores. Threads 5 to 8 run with Hyper-Threading on the same cores, and the program
    uses the arithmetic unit all the time, so there are no free units for the second thread to use.
\end{hall}

\appendix

%================================================
\chapter{8086 Instruction Summary}
%================================================

This appendix puts in one place all the instructions we studied in Chapter 3, so you can find them quickly while
you write your programs. Symbols: d = destination, s = source, r = register, m = memory box.

\section*{The most used DOS functions (\code{int 21h})}
\begin{center}
\small
\begin{tabular}{|c|l|l|}
    \hline
    \ras{AH} & \ras{Function} & \ras{Inputs / outputs} \\ \hline
    \hex{01} & read a key and show it & output in AL \\ \hline
    \hex{02} & print a character & the character in DL \\ \hline
    \hex{09} & print a string ending with \code{\$} & the address in DX \\ \hline
    \hex{0A} & read a line into a buffer & the address in DX \\ \hline
    \hex{4C} & end the program & the exit code in AL \\ \hline
\end{tabular}
\end{center}

\section*{The instructions}
\begin{center}
\small\renewcommand{\arraystretch}{1.12}
\begin{tabularx}{\linewidth}{|l|X|l|}
    \hline
    \ras{Instruction} & \ras{Meaning} & \ras{Flags changed} \\ \hline
    \multicolumn{3}{|c|}{\cellcolor{brandOrange!20}\textbf{Transfer}} \\ \hline
    \code{mov d, s} & \nql{d}{s} & none \\ \hline
    \code{xchg a, b} & swap a and b & none \\ \hline
    \code{lea r, m} & r gets the address of m & none \\ \hline
    \code{push s} & \nql{SP}{SP - 2}, then \nql{[SP]}{s} & none \\ \hline
    \code{pop d} & \nql{d}{[SP]}, then \nql{SP}{SP + 2} & none \\ \hline
    \code{in al, port}, \code{out port, al} & read from a port, write to a port & none \\ \hline
    \multicolumn{3}{|c|}{\cellcolor{brandOrange!20}\textbf{Arithmetic}} \\ \hline
    \code{add d, s} & \nql{d}{d + s} & all \\ \hline
    \code{adc d, s} & \nql{d}{d + s + CF} & all \\ \hline
    \code{sub d, s} & \nql{d}{d - s} & all \\ \hline
    \code{sbb d, s} & \nql{d}{d - s - CF} & all \\ \hline
    \code{inc d} & \nql{d}{d + 1} & all except CF \\ \hline
    \code{dec d} & \nql{d}{d - 1} & all except CF \\ \hline
    \code{neg d} & \nql{d}{-d} & all \\ \hline
    \code{cmp a, b} & computes a - b without keeping it & all \\ \hline
    \code{mul s} (8 bits) & \nql{AX}{AL * s} & CF, OF \\ \hline
    \code{mul s} (16 bits) & \nql{DX:AX}{AX * s} & CF, OF \\ \hline
    \code{div s} (8 bits) & \nql{AL}{AX / s}, \nql{AH}{AX mod s} & undefined \\ \hline
    \code{div s} (16 bits) & \nql{AX}{DX:AX / s}, \nql{DX}{DX:AX mod s} & undefined \\ \hline
    \multicolumn{3}{|c|}{\cellcolor{brandOrange!20}\textbf{Logic and shifts}} \\ \hline
    \code{and}, \code{or}, \code{xor} & bit by bit operations & CF = OF = 0, ZF, SF \\ \hline
    \code{not d} & flip all the bits & none \\ \hline
    \code{test a, b} & \code{and} without keeping the result & like \code{and} \\ \hline
    \code{shl d, n}, \code{shr d, n} & shift: multiply or divide by $2^n$ & CF = last bit out \\ \hline
    \code{rol d, n}, \code{ror d, n} & rotate & CF \\ \hline
    \multicolumn{3}{|c|}{\cellcolor{brandOrange!20}\textbf{Flow control}} \\ \hline
    \code{jmp L} & always jump & none \\ \hline
    \code{je}, \code{jne}, \code{jz}, \code{jnz} & depends on ZF & -- \\ \hline
    \code{jg}, \code{jge}, \code{jl}, \code{jle} & \textbf{signed} comparison & -- \\ \hline
    \code{ja}, \code{jae}, \code{jb}, \code{jbe} & \textbf{unsigned} comparison & -- \\ \hline
    \code{loop L} & \nql{CX}{CX - 1}, then jump if CX $\neq 0$ & none \\ \hline
    \code{jcxz L} & jump if CX = 0 & -- \\ \hline
    \code{call P}, \code{ret} & call a procedure, come back from it & none \\ \hline
    \code{int n}, \code{iret} & interrupt, come back from it & IF = TF = 0 / restored \\ \hline
    \multicolumn{3}{|c|}{\cellcolor{brandOrange!20}\textbf{Strings}} \\ \hline
    \code{movsb}, \code{movsw} & \nql{ES:[DI]}{DS:[SI]}, then move SI and DI forward & none \\ \hline
    \code{rep} & repeat the instruction CX times & -- \\ \hline
    \code{cld}, \code{std} & DF = 0 (forward), DF = 1 & DF \\ \hline
\end{tabularx}
\end{center}

%================================================
\chapter{EMU8086 and NASM: differences in writing}
%================================================

The programs of Chapter 3 are written for EMU8086 (a syntax close to MASM), and the programs of Chapter 4 for NASM.
The processor is the same, but the two assemblers do not write some things the same way:

\begin{center}
\small
\begin{tabularx}{\linewidth}{|X|l|l|}
    \hline
    \ras{What we want} & \ras{EMU8086 / MASM} & \ras{NASM} \\ \hline
    address of a variable in a register & \code{lea dx, msg} or \code{mov dx, offset msg} & \code{mov dx, msg} \\ \hline
    value of a variable & \code{mov al, x} or \code{mov al, [x]} & \code{mov al, [x]} (brackets needed) \\ \hline
    giving the size & \code{mov byte ptr [bx], 5} & \code{mov byte [bx], 5} \\ \hline
    a hex number & \code{0FFh} & \code{0xFF} or \code{0FFh} \\ \hline
    a procedure & \code{p proc} \dots{} \code{p endp} & a simple label \code{p:} \\ \hline
    a constant & \code{N equ 10} & \code{N equ 10} \\ \hline
    an array of 10 zeros & \code{t db 10 dup(0)} & \code{t times 10 db 0} \\ \hline
    assembling & the \emph{Emulate} button & \code{nasm -f bin prog.asm -o prog.com} \\ \hline
\end{tabularx}
\end{center}

\begin{intabih}
    A good habit: always write the brackets \code{[ ]} when you mean the content of memory, even in EMU8086. Your program
    stays clear, and works with both assemblers with few changes. This is what we did in all the programs of this course.
\end{intabih}

%================================================
\chapter{Glossary}
%================================================

\begin{center}
\small\renewcommand{\arraystretch}{1.08}
\begin{tabular}{|l|l|}
    \hline
    \ras{English} & \ras{Français} \\ \hline
    computer & ordinateur \\ \hline
    processor, CPU & processeur, UCT \\ \hline
    central unit & unité centrale \\ \hline
    main memory & mémoire centrale \\ \hline
    volatile memory & mémoire volatile \\ \hline
    bus & bus \\ \hline
    address bus / data bus / control bus & bus d'adresses / de données / de contrôle \\ \hline
    register & registre \\ \hline
    program counter (PC) & compteur ordinal (CO) \\ \hline
    instruction register (IR) & registre d'instruction (RI) \\ \hline
    memory address register (MAR) & registre d'adresse mémoire (RAM) \\ \hline
    memory data register (MDR) & registre d'information mémoire (RIM) \\ \hline
    accumulator & accumulateur \\ \hline
    flag & indicateur, drapeau \\ \hline
    arithmetic logic unit (ALU) & unité arithmétique et logique (UAL) \\ \hline
    control unit & unité de commande (UC) \\ \hline
    decoder / sequencer & décodeur / séquenceur \\ \hline
    half adder / full adder & demi-additionneur / additionneur complet \\ \hline
    two's complement & complément à deux \\ \hline
    instruction set & jeu d'instructions \\ \hline
    opcode & code opération \\ \hline
    addressing mode & mode d'adressage \\ \hline
    effective address / physical address & adresse effective / adresse physique \\ \hline
    segment, offset & segment, déplacement \\ \hline
    instruction cycle & cycle d'instruction \\ \hline
    fetch, decode, execute & recherche, décodage, exécution \\ \hline
    assembly language / assembler & langage d'assemblage / assembleur \\ \hline
    stack & pile \\ \hline
    procedure, subroutine & procédure, sous-programme \\ \hline
    interrupt / interrupt vector table & interruption / table des vecteurs d'interruption \\ \hline
    I/O port & port d'entrée/sortie \\ \hline
    protected mode & mode protégé \\ \hline
    flat memory model & modèle mémoire plat \\ \hline
    pipeline / branch prediction & pipeline / prédiction de branchement \\ \hline
    core / thread & cœur / thread, fil d'exécution \\ \hline
    cache memory / cache line & mémoire cache / ligne de cache \\ \hline
    hit / miss & succès / défaut \\ \hline
    temporal and spatial locality & localité temporelle et spatiale \\ \hline
\end{tabular}
\end{center}

\section*{References}
\begin{enumerate}
    \item A. S. Tanenbaum, \emph{Structured Computer Organization}, Pearson, 2012.
    \item D. A. Patterson, J. L. Hennessy, \emph{Computer Organization and Design: The Hardware/Software Interface}, Morgan Kaufmann.
    \item P. Zanella, Y. Ligier, E. Lazard, \emph{Architecture et technologie des ordinateurs}, 6e édition, Dunod, 2018.
    \item B. B. Brey, \emph{The Intel Microprocessors: 8086/8088 \dots{} Core2 with 64-bit Extensions}, Pearson.
    \item EMU8086, Microprocessor Emulator.
    \item NASM, The Netwide Assembler: https://www.nasm.us/
    \item MASM, Microsoft Macro Assembler reference: \\ {\small https://learn.microsoft.com/en-us/cpp/assembler/masm/}
    \item GAS, The GNU Assembler: https://sourceware.org/binutils/docs/as/
\end{enumerate}


\end{document}
